% ---------------------------------------------------------------------------
% metabolites.tex --- key metabolites of the extended glycolysis network.
%
% Source fragment: loaded by SI/si.tex (and exported standalone by
% Figure/export_tikz.py), replacing the four auto-generated chemfig dumps in
% Figure/moltex/, which drew every hydrogen at four different scales.
%
% Every structure is the constitution and configuration of the SMILES listed
% in Table S2 (tokens C, S, H, J).  Sugars use Haworth projections: bold
% front bonds, substituents up or down.  Phosphoryl groups are drawn in the
% shared highlight colour, matching the phosphate convention of the
% donor-selectivity figure.  The anomeric centre of fructose 1,6-bisphosphate
% is left unspecified because the listed SMILES does not specify it.
% ---------------------------------------------------------------------------
\IfFileExists{synstyle.tex}{% ---------------------------------------------------------------------------
% synstyle.tex --- shared look for every TikZ figure of the SynKit CRN paper.
%
% Input it from a figure file:
%     \IfFileExists{synstyle.tex}{% ---------------------------------------------------------------------------
% synstyle.tex --- shared look for every TikZ figure of the SynKit CRN paper.
%
% Input it from a figure file:
%     \IfFileExists{synstyle.tex}{% ---------------------------------------------------------------------------
% synstyle.tex --- shared look for every TikZ figure of the SynKit CRN paper.
%
% Input it from a figure file:
%     \IfFileExists{synstyle.tex}{\input{synstyle}}{\input{Figure/synstyle}}
% so that the figure compiles both from Figure/ and from the paper root.
%
% The visual vocabulary follows ../../../../Style/style.tex: the
% colour-blind-safe Okabe--Ito accents, quiet neutral rules and wash, circular
% panel badges, and a compact technical type scale are shared by the companion
% papers.  Semantic aliases below keep the CRN meaning explicit.
%
% DESIGN CONTRACT -----------------------------------------------------------
% 1. Geometry.  Every figure is laid out on a canvas exactly \scW wide (the
%    sn-jnl text width, 372 pt) and is included unscaled, so one point of type
%    is the same physical size in all five figures.  Open each picture with
%    \scCanvas{<height in cm>} to pin the bounding box.
% 2. Colour is semantic, never decorative.  Four roles only:
%       scAccent (green)       structure, the normal/expected path
%       scMark   (orange)      catalyst, certificate, or highlight
%       scAlert  (vermillion)  failure, violation, or uncovered set
%       scRef    (blue)        secondary or reference series
%    A panel gets at most one accent.  Titles are always ink; panel identity
%    is also encoded by a letter inside a circular badge, not by colour alone.
% 3. Type scale.  Three sizes, no others, and nothing below 6 pt:
%       \scTitle (7 pt bold)  panel titles
%       \scBody  (6.5 pt)     labels, values, in-figure prose
%       \scSmall (6 pt)       units, footnotes, dense graph labels
% ---------------------------------------------------------------------------
% Load guard.  The paper inputs every figure source inline, so this file is
% reached once per figure; without the guard the second \newlength below
% would abort the run.
\ifdefined\SynKitStyleLoaded\endinput\fi
\newcommand{\SynKitStyleLoaded}{}

% Load in the document preamble. XeLaTeX/LuaLaTeX and an installed Arial
% family are required; fontspec reports an error instead of substituting.
\RequirePackage[no-math]{fontspec}
\newfontfamily\scArial{Arial}[NFSSFamily=scArial,Ligatures=TeX]
% Keep manuscript mathematics unchanged; select Arial letters and numerals
% only inside figures. TeX supplies specialist mathematical symbols.
\RequirePackage[subdued,italic]{mathastext}
\MTDeclareVersion{synarial}{TU}{scArial}{m}{n}
\SetMathAlphabet{\Mathsf}{synarial}{TU}{scArial}{m}{n}
\SetMathAlphabet{\Mathtt}{synarial}{TU}{scArial}{m}{n}

\usetikzlibrary{positioning,calc,fit,backgrounds,arrows.meta,shapes.geometric,
                shapes.misc,decorations.pathreplacing,patterns,math}

% --- canvas ---------------------------------------------------------------
% 372 pt text width minus the 1 pt standalone border on each side.
\newlength{\scW}\setlength{\scW}{13.0326cm}
\providecommand{\scWnum}{13.0326}
% \scCanvas{<height>} pins the picture to the full text width, y running down
% from 0, so every figure shares one coordinate convention.  It *sets* the
% bounding box rather than contributing to it, so the emitted page is exactly
% the text width and every figure is included at 1:1 -- otherwise a stray
% invisible extent (pgfplots reserves a little around an axis) rescales one
% figure relative to the others and the same \scSmall prints at two sizes.
% Anything drawn outside the canvas is therefore cropped: keep content inside,
% and check with
%     \providecommand{\scCanvas}[1]{\path[draw=red] (0,0) rectangle (\scWnum,-#1);}
\providecommand{\scCanvas}[1]{%
  \path[use as bounding box] (0,0) rectangle (\scWnum,-#1);}

% --- palette --------------------------------------------------------------
% Canonical companion-paper values (../../../../Style/style.tex).
\definecolor{npgOrange}{HTML}{E69F00}
\definecolor{npgSky}{HTML}{56B4E9}
\definecolor{npgGreen}{HTML}{009E73}
\definecolor{npgYellow}{HTML}{F0E442}
\definecolor{npgBlue}{HTML}{0072B2}
\definecolor{npgVermillion}{HTML}{E64B35}
\definecolor{npgPurple}{HTML}{CC79A7}
\definecolor{figInk}{HTML}{1A1A1A}
\definecolor{figMuted}{HTML}{666666}
\definecolor{figHair}{HTML}{D9D9D9}
\definecolor{figWash}{HTML}{F7F7F7}

\colorlet{scInk}{figInk}              % primary text
\colorlet{scInkM}{figMuted}           % secondary text
\colorlet{scPaper}{figWash}           % quiet cards and note strips
\colorlet{scSpF}{npgSky!12}           % species fill
\colorlet{scSpE}{npgBlue}             % species edge
\colorlet{scRxF}{npgOrange!12}        % reaction fill
\colorlet{scRxE}{npgOrange!88!black}  % reaction edge
\colorlet{scReac}{npgBlue!78!black}   % reactant edge
\colorlet{scProd}{npgGreen!82!black}  % product edge
\colorlet{scHiF}{npgOrange!13}        % highlight fill
\colorlet{scHiE}{npgOrange!82!black}  % highlight edge
% Stage / accent tints.  stDE remains available for the fifth plot series.
\colorlet{stAF}{npgGreen!9}\colorlet{stAE}{npgGreen}
\colorlet{stBF}{npgOrange!11}\colorlet{stBE}{npgOrange}
\colorlet{stCF}{npgBlue!9}\colorlet{stCE}{npgBlue}
\colorlet{stDF}{npgPurple!10}\colorlet{stDE}{npgPurple!85!black}
\colorlet{stEF}{npgVermillion!9}\colorlet{stEE}{npgVermillion}

% Semantic aliases.  Figures should speak in these, not in stAE/stBE/...
\colorlet{scAccent}{stAE}\colorlet{scAccentF}{stAF}   % structure / normal
\colorlet{scMark}{stBE}\colorlet{scMarkF}{stBF}       % highlight / certificate
\colorlet{scRef}{stCE}\colorlet{scRefF}{stCF}         % secondary series
\colorlet{scAlert}{stEE}\colorlet{scAlertF}{stEF}     % failure / violation
\colorlet{scRule}{figHair}                            % card and axis hairlines
\colorlet{scRuleL}{figHair!55}                        % grid, dividers

% --- text styles ----------------------------------------------------------
% Three sizes, fixed in points so they do not drift with the document class.
% Arial is shared by labels, headings, panel badges, and chemistry labels.
\providecommand{\scTitle}{\scArial\fontsize{7}{8.4}\selectfont\bfseries\color{scInk}}
\providecommand{\scBody}{\scArial\fontsize{6.5}{7.8}\selectfont\color{scInkM}}
\providecommand{\scSmall}{\scArial\fontsize{6}{7.2}\selectfont\color{scInkM}}
\providecommand{\scSub}[1]{{\scSmall\itshape #1}}
\providecommand{\scPanel}[1]{{\scTitle(#1)}}
% Numerals that carry a result: ink, not colour, so the eye reads the value.
\providecommand{\scNum}{\scArial\fontsize{10}{11}\selectfont\bfseries\color{scInk}}

% --- shared TikZ styles ---------------------------------------------------
\tikzset{
  % text never hyphenates inside a figure
  every node/.append style={execute at begin node={%
      \hyphenpenalty=10000 \exhyphenpenalty=10000 \hbadness=10000 }},
  syn/.style={font=\scArial,color=scInk,
              execute at begin picture={\MTversion*{synarial}},
              >={Stealth[round,length=2.4mm,width=1.9mm]}},
  % containers ------------------------------------------------------------
  % One card chrome for the whole paper: white, 3 pt radius, one hairline.
  panel/.style={rounded corners=2.5pt,fill=none,draw=scRule,line width=0.45pt},
  card/.style={rounded corners=2.2pt,fill=scPaper,draw=scRule,line width=0.45pt,
               inner sep=5pt,align=left,anchor=north west},
  % a quiet strip for a panel's closing fact
  note/.style={rounded corners=2.5pt,fill=scPaper,draw=scRuleL,line width=0.4pt,
               inner xsep=6pt,inner ysep=3.5pt,align=left,font=\scSmall},
  % small labelled value pill
  chip/.style={rounded corners=2pt,fill=scPaper,draw=scRule,line width=0.4pt,
               inner xsep=4pt,inner ysep=2pt,font=\scSmall},
  % Circular letter badge: the same accessible panel marker used in Style/*.
  hbullet/.style={circle,fill=scInk,text=white,draw=none,
                  font=\scArial\fontsize{6}{6}\selectfont\bfseries,
                  minimum size=4.2mm,inner sep=0pt},
  ctag/.style={font=\scTitle,anchor=north west,text=scInk},
  % CRN primitives --------------------------------------------------------
  species/.style={rounded rectangle,fill=scSpF,draw=scSpE,line width=0.45pt,
                  inner xsep=2.4pt,inner ysep=1.6pt,font=\scSmall,
                  text=scInk,minimum height=3.6mm},
  cofactor/.style={species,fill=scRefF,draw=scRef},
  catalyst/.style={species,fill=scHiF,draw=scHiE},
  rxn/.style={rectangle,fill=scRxF,draw=scRxE,line width=0.45pt,
              inner sep=1.2pt,font=\scSmall,text=scInk,
              minimum width=3.0mm,minimum height=3.0mm},
  place/.style={circle,fill=scSpF,draw=scSpE,line width=0.45pt,
                minimum size=4.0mm,inner sep=0pt},
  trans/.style={rectangle,fill=scRxE,draw=scRxE,minimum width=1.1mm,
                minimum height=3.6mm,inner sep=0pt},
  token/.style={circle,fill=scInk,minimum size=1.1mm,inner sep=0pt},
  % edges -----------------------------------------------------------------
  reac/.style={->,line width=0.5pt,draw=scReac,shorten >=0.8pt,shorten <=0.8pt},
  prod/.style={->,line width=0.5pt,draw=scProd,shorten >=0.8pt,shorten <=0.8pt},
  cat/.style={->,line width=0.5pt,draw=scHiE,shorten >=0.8pt,shorten <=0.8pt},
  flow/.style={-{Stealth[round,length=2.2mm,width=1.8mm]},line width=1.1pt,
               draw=scInkM!45},
  flab/.style={font=\scSmall,text=scInkM,inner sep=2pt},
  stoi/.style={font=\scSmall,text=scInkM,inner sep=0.8pt},
  % verdict marks: shape carries the meaning, colour only reinforces it
  vpass/.style={circle,fill=scAccentF,draw=scAccent,text=scAccent,
                font=\fontsize{6}{6}\selectfont\bfseries,
                minimum size=3.2mm,inner sep=0pt},
  vfail/.style={circle,fill=scAlertF,draw=scAlert,text=scAlert,
                font=\fontsize{6}{6}\selectfont\bfseries,
                minimum size=3.2mm,inner sep=0pt},
  callout/.style={font=\scSmall,text=scInkM,align=left},
}

% \scHead[<badge colour>]{<anchor coord>}{<letter>}{<title>}{<subtitle or empty>}
% One header treatment for every panel in the paper.
\providecommand{\scHead}[5][scInk]{%
  \node[hbullet,fill=#1,anchor=north west] (scbadge#3) at ($#2+(0,-0.015)$)
    {\MakeUppercase{#3}};
  \node[ctag,anchor=west] at ($(scbadge#3.east)+(0.14,0)$) {#4};
  \ifx\relax#5\relax\else
    \node[anchor=north west,font=\scSmall,text=scInkM]
      at ($#2+(0.56,-0.245)$) {#5};
  \fi}

% --- shared pgfplots look (only if the figure loaded pgfplots) ------------
\ifdefined\pgfplotsset
  \pgfplotsset{compat=1.18}
  \pgfplotsset{
    synaxis/.style={
      scale only axis,
      tick align=outside, tick pos=left,
      axis line style={draw=scInkM!55,line width=0.4pt},
      % One grid, on the value axis only.  No minor grid: at these sizes it
      % turns the plot frame into a grey wash.
      ymajorgrids=true, xmajorgrids=false,
      major grid style={draw=scRuleL,line width=0.3pt},
      major tick length=1.6pt,
      label style={font=\scBody,color=scInk},
      tick label style={font=\scSmall,color=scInkM},
      legend style={font=\scSmall,draw=none,fill=none,inner sep=1pt,
                    row sep=0.4pt,legend columns=1},
      legend cell align=left,
      every axis plot/.append style={line width=0.7pt,mark size=1.2pt},
    },
  }
\fi
}{% ---------------------------------------------------------------------------
% synstyle.tex --- shared look for every TikZ figure of the SynKit CRN paper.
%
% Input it from a figure file:
%     \IfFileExists{synstyle.tex}{\input{synstyle}}{\input{Figure/synstyle}}
% so that the figure compiles both from Figure/ and from the paper root.
%
% The visual vocabulary follows ../../../../Style/style.tex: the
% colour-blind-safe Okabe--Ito accents, quiet neutral rules and wash, circular
% panel badges, and a compact technical type scale are shared by the companion
% papers.  Semantic aliases below keep the CRN meaning explicit.
%
% DESIGN CONTRACT -----------------------------------------------------------
% 1. Geometry.  Every figure is laid out on a canvas exactly \scW wide (the
%    sn-jnl text width, 372 pt) and is included unscaled, so one point of type
%    is the same physical size in all five figures.  Open each picture with
%    \scCanvas{<height in cm>} to pin the bounding box.
% 2. Colour is semantic, never decorative.  Four roles only:
%       scAccent (green)       structure, the normal/expected path
%       scMark   (orange)      catalyst, certificate, or highlight
%       scAlert  (vermillion)  failure, violation, or uncovered set
%       scRef    (blue)        secondary or reference series
%    A panel gets at most one accent.  Titles are always ink; panel identity
%    is also encoded by a letter inside a circular badge, not by colour alone.
% 3. Type scale.  Three sizes, no others, and nothing below 6 pt:
%       \scTitle (7 pt bold)  panel titles
%       \scBody  (6.5 pt)     labels, values, in-figure prose
%       \scSmall (6 pt)       units, footnotes, dense graph labels
% ---------------------------------------------------------------------------
% Load guard.  The paper inputs every figure source inline, so this file is
% reached once per figure; without the guard the second \newlength below
% would abort the run.
\ifdefined\SynKitStyleLoaded\endinput\fi
\newcommand{\SynKitStyleLoaded}{}

% Load in the document preamble. XeLaTeX/LuaLaTeX and an installed Arial
% family are required; fontspec reports an error instead of substituting.
\RequirePackage[no-math]{fontspec}
\newfontfamily\scArial{Arial}[NFSSFamily=scArial,Ligatures=TeX]
% Keep manuscript mathematics unchanged; select Arial letters and numerals
% only inside figures. TeX supplies specialist mathematical symbols.
\RequirePackage[subdued,italic]{mathastext}
\MTDeclareVersion{synarial}{TU}{scArial}{m}{n}
\SetMathAlphabet{\Mathsf}{synarial}{TU}{scArial}{m}{n}
\SetMathAlphabet{\Mathtt}{synarial}{TU}{scArial}{m}{n}

\usetikzlibrary{positioning,calc,fit,backgrounds,arrows.meta,shapes.geometric,
                shapes.misc,decorations.pathreplacing,patterns,math}

% --- canvas ---------------------------------------------------------------
% 372 pt text width minus the 1 pt standalone border on each side.
\newlength{\scW}\setlength{\scW}{13.0326cm}
\providecommand{\scWnum}{13.0326}
% \scCanvas{<height>} pins the picture to the full text width, y running down
% from 0, so every figure shares one coordinate convention.  It *sets* the
% bounding box rather than contributing to it, so the emitted page is exactly
% the text width and every figure is included at 1:1 -- otherwise a stray
% invisible extent (pgfplots reserves a little around an axis) rescales one
% figure relative to the others and the same \scSmall prints at two sizes.
% Anything drawn outside the canvas is therefore cropped: keep content inside,
% and check with
%     \providecommand{\scCanvas}[1]{\path[draw=red] (0,0) rectangle (\scWnum,-#1);}
\providecommand{\scCanvas}[1]{%
  \path[use as bounding box] (0,0) rectangle (\scWnum,-#1);}

% --- palette --------------------------------------------------------------
% Canonical companion-paper values (../../../../Style/style.tex).
\definecolor{npgOrange}{HTML}{E69F00}
\definecolor{npgSky}{HTML}{56B4E9}
\definecolor{npgGreen}{HTML}{009E73}
\definecolor{npgYellow}{HTML}{F0E442}
\definecolor{npgBlue}{HTML}{0072B2}
\definecolor{npgVermillion}{HTML}{E64B35}
\definecolor{npgPurple}{HTML}{CC79A7}
\definecolor{figInk}{HTML}{1A1A1A}
\definecolor{figMuted}{HTML}{666666}
\definecolor{figHair}{HTML}{D9D9D9}
\definecolor{figWash}{HTML}{F7F7F7}

\colorlet{scInk}{figInk}              % primary text
\colorlet{scInkM}{figMuted}           % secondary text
\colorlet{scPaper}{figWash}           % quiet cards and note strips
\colorlet{scSpF}{npgSky!12}           % species fill
\colorlet{scSpE}{npgBlue}             % species edge
\colorlet{scRxF}{npgOrange!12}        % reaction fill
\colorlet{scRxE}{npgOrange!88!black}  % reaction edge
\colorlet{scReac}{npgBlue!78!black}   % reactant edge
\colorlet{scProd}{npgGreen!82!black}  % product edge
\colorlet{scHiF}{npgOrange!13}        % highlight fill
\colorlet{scHiE}{npgOrange!82!black}  % highlight edge
% Stage / accent tints.  stDE remains available for the fifth plot series.
\colorlet{stAF}{npgGreen!9}\colorlet{stAE}{npgGreen}
\colorlet{stBF}{npgOrange!11}\colorlet{stBE}{npgOrange}
\colorlet{stCF}{npgBlue!9}\colorlet{stCE}{npgBlue}
\colorlet{stDF}{npgPurple!10}\colorlet{stDE}{npgPurple!85!black}
\colorlet{stEF}{npgVermillion!9}\colorlet{stEE}{npgVermillion}

% Semantic aliases.  Figures should speak in these, not in stAE/stBE/...
\colorlet{scAccent}{stAE}\colorlet{scAccentF}{stAF}   % structure / normal
\colorlet{scMark}{stBE}\colorlet{scMarkF}{stBF}       % highlight / certificate
\colorlet{scRef}{stCE}\colorlet{scRefF}{stCF}         % secondary series
\colorlet{scAlert}{stEE}\colorlet{scAlertF}{stEF}     % failure / violation
\colorlet{scRule}{figHair}                            % card and axis hairlines
\colorlet{scRuleL}{figHair!55}                        % grid, dividers

% --- text styles ----------------------------------------------------------
% Three sizes, fixed in points so they do not drift with the document class.
% Arial is shared by labels, headings, panel badges, and chemistry labels.
\providecommand{\scTitle}{\scArial\fontsize{7}{8.4}\selectfont\bfseries\color{scInk}}
\providecommand{\scBody}{\scArial\fontsize{6.5}{7.8}\selectfont\color{scInkM}}
\providecommand{\scSmall}{\scArial\fontsize{6}{7.2}\selectfont\color{scInkM}}
\providecommand{\scSub}[1]{{\scSmall\itshape #1}}
\providecommand{\scPanel}[1]{{\scTitle(#1)}}
% Numerals that carry a result: ink, not colour, so the eye reads the value.
\providecommand{\scNum}{\scArial\fontsize{10}{11}\selectfont\bfseries\color{scInk}}

% --- shared TikZ styles ---------------------------------------------------
\tikzset{
  % text never hyphenates inside a figure
  every node/.append style={execute at begin node={%
      \hyphenpenalty=10000 \exhyphenpenalty=10000 \hbadness=10000 }},
  syn/.style={font=\scArial,color=scInk,
              execute at begin picture={\MTversion*{synarial}},
              >={Stealth[round,length=2.4mm,width=1.9mm]}},
  % containers ------------------------------------------------------------
  % One card chrome for the whole paper: white, 3 pt radius, one hairline.
  panel/.style={rounded corners=2.5pt,fill=none,draw=scRule,line width=0.45pt},
  card/.style={rounded corners=2.2pt,fill=scPaper,draw=scRule,line width=0.45pt,
               inner sep=5pt,align=left,anchor=north west},
  % a quiet strip for a panel's closing fact
  note/.style={rounded corners=2.5pt,fill=scPaper,draw=scRuleL,line width=0.4pt,
               inner xsep=6pt,inner ysep=3.5pt,align=left,font=\scSmall},
  % small labelled value pill
  chip/.style={rounded corners=2pt,fill=scPaper,draw=scRule,line width=0.4pt,
               inner xsep=4pt,inner ysep=2pt,font=\scSmall},
  % Circular letter badge: the same accessible panel marker used in Style/*.
  hbullet/.style={circle,fill=scInk,text=white,draw=none,
                  font=\scArial\fontsize{6}{6}\selectfont\bfseries,
                  minimum size=4.2mm,inner sep=0pt},
  ctag/.style={font=\scTitle,anchor=north west,text=scInk},
  % CRN primitives --------------------------------------------------------
  species/.style={rounded rectangle,fill=scSpF,draw=scSpE,line width=0.45pt,
                  inner xsep=2.4pt,inner ysep=1.6pt,font=\scSmall,
                  text=scInk,minimum height=3.6mm},
  cofactor/.style={species,fill=scRefF,draw=scRef},
  catalyst/.style={species,fill=scHiF,draw=scHiE},
  rxn/.style={rectangle,fill=scRxF,draw=scRxE,line width=0.45pt,
              inner sep=1.2pt,font=\scSmall,text=scInk,
              minimum width=3.0mm,minimum height=3.0mm},
  place/.style={circle,fill=scSpF,draw=scSpE,line width=0.45pt,
                minimum size=4.0mm,inner sep=0pt},
  trans/.style={rectangle,fill=scRxE,draw=scRxE,minimum width=1.1mm,
                minimum height=3.6mm,inner sep=0pt},
  token/.style={circle,fill=scInk,minimum size=1.1mm,inner sep=0pt},
  % edges -----------------------------------------------------------------
  reac/.style={->,line width=0.5pt,draw=scReac,shorten >=0.8pt,shorten <=0.8pt},
  prod/.style={->,line width=0.5pt,draw=scProd,shorten >=0.8pt,shorten <=0.8pt},
  cat/.style={->,line width=0.5pt,draw=scHiE,shorten >=0.8pt,shorten <=0.8pt},
  flow/.style={-{Stealth[round,length=2.2mm,width=1.8mm]},line width=1.1pt,
               draw=scInkM!45},
  flab/.style={font=\scSmall,text=scInkM,inner sep=2pt},
  stoi/.style={font=\scSmall,text=scInkM,inner sep=0.8pt},
  % verdict marks: shape carries the meaning, colour only reinforces it
  vpass/.style={circle,fill=scAccentF,draw=scAccent,text=scAccent,
                font=\fontsize{6}{6}\selectfont\bfseries,
                minimum size=3.2mm,inner sep=0pt},
  vfail/.style={circle,fill=scAlertF,draw=scAlert,text=scAlert,
                font=\fontsize{6}{6}\selectfont\bfseries,
                minimum size=3.2mm,inner sep=0pt},
  callout/.style={font=\scSmall,text=scInkM,align=left},
}

% \scHead[<badge colour>]{<anchor coord>}{<letter>}{<title>}{<subtitle or empty>}
% One header treatment for every panel in the paper.
\providecommand{\scHead}[5][scInk]{%
  \node[hbullet,fill=#1,anchor=north west] (scbadge#3) at ($#2+(0,-0.015)$)
    {\MakeUppercase{#3}};
  \node[ctag,anchor=west] at ($(scbadge#3.east)+(0.14,0)$) {#4};
  \ifx\relax#5\relax\else
    \node[anchor=north west,font=\scSmall,text=scInkM]
      at ($#2+(0.56,-0.245)$) {#5};
  \fi}

% --- shared pgfplots look (only if the figure loaded pgfplots) ------------
\ifdefined\pgfplotsset
  \pgfplotsset{compat=1.18}
  \pgfplotsset{
    synaxis/.style={
      scale only axis,
      tick align=outside, tick pos=left,
      axis line style={draw=scInkM!55,line width=0.4pt},
      % One grid, on the value axis only.  No minor grid: at these sizes it
      % turns the plot frame into a grey wash.
      ymajorgrids=true, xmajorgrids=false,
      major grid style={draw=scRuleL,line width=0.3pt},
      major tick length=1.6pt,
      label style={font=\scBody,color=scInk},
      tick label style={font=\scSmall,color=scInkM},
      legend style={font=\scSmall,draw=none,fill=none,inner sep=1pt,
                    row sep=0.4pt,legend columns=1},
      legend cell align=left,
      every axis plot/.append style={line width=0.7pt,mark size=1.2pt},
    },
  }
\fi
}
% so that the figure compiles both from Figure/ and from the paper root.
%
% The visual vocabulary follows ../../../../Style/style.tex: the
% colour-blind-safe Okabe--Ito accents, quiet neutral rules and wash, circular
% panel badges, and a compact technical type scale are shared by the companion
% papers.  Semantic aliases below keep the CRN meaning explicit.
%
% DESIGN CONTRACT -----------------------------------------------------------
% 1. Geometry.  Every figure is laid out on a canvas exactly \scW wide (the
%    sn-jnl text width, 372 pt) and is included unscaled, so one point of type
%    is the same physical size in all five figures.  Open each picture with
%    \scCanvas{<height in cm>} to pin the bounding box.
% 2. Colour is semantic, never decorative.  Four roles only:
%       scAccent (green)       structure, the normal/expected path
%       scMark   (orange)      catalyst, certificate, or highlight
%       scAlert  (vermillion)  failure, violation, or uncovered set
%       scRef    (blue)        secondary or reference series
%    A panel gets at most one accent.  Titles are always ink; panel identity
%    is also encoded by a letter inside a circular badge, not by colour alone.
% 3. Type scale.  Three sizes, no others, and nothing below 6 pt:
%       \scTitle (7 pt bold)  panel titles
%       \scBody  (6.5 pt)     labels, values, in-figure prose
%       \scSmall (6 pt)       units, footnotes, dense graph labels
% ---------------------------------------------------------------------------
% Load guard.  The paper inputs every figure source inline, so this file is
% reached once per figure; without the guard the second \newlength below
% would abort the run.
\ifdefined\SynKitStyleLoaded\endinput\fi
\newcommand{\SynKitStyleLoaded}{}

% Load in the document preamble. XeLaTeX/LuaLaTeX and an installed Arial
% family are required; fontspec reports an error instead of substituting.
\RequirePackage[no-math]{fontspec}
\newfontfamily\scArial{Arial}[NFSSFamily=scArial,Ligatures=TeX]
% Keep manuscript mathematics unchanged; select Arial letters and numerals
% only inside figures. TeX supplies specialist mathematical symbols.
\RequirePackage[subdued,italic]{mathastext}
\MTDeclareVersion{synarial}{TU}{scArial}{m}{n}
\SetMathAlphabet{\Mathsf}{synarial}{TU}{scArial}{m}{n}
\SetMathAlphabet{\Mathtt}{synarial}{TU}{scArial}{m}{n}

\usetikzlibrary{positioning,calc,fit,backgrounds,arrows.meta,shapes.geometric,
                shapes.misc,decorations.pathreplacing,patterns,math}

% --- canvas ---------------------------------------------------------------
% 372 pt text width minus the 1 pt standalone border on each side.
\newlength{\scW}\setlength{\scW}{13.0326cm}
\providecommand{\scWnum}{13.0326}
% \scCanvas{<height>} pins the picture to the full text width, y running down
% from 0, so every figure shares one coordinate convention.  It *sets* the
% bounding box rather than contributing to it, so the emitted page is exactly
% the text width and every figure is included at 1:1 -- otherwise a stray
% invisible extent (pgfplots reserves a little around an axis) rescales one
% figure relative to the others and the same \scSmall prints at two sizes.
% Anything drawn outside the canvas is therefore cropped: keep content inside,
% and check with
%     \providecommand{\scCanvas}[1]{\path[draw=red] (0,0) rectangle (\scWnum,-#1);}
\providecommand{\scCanvas}[1]{%
  \path[use as bounding box] (0,0) rectangle (\scWnum,-#1);}

% --- palette --------------------------------------------------------------
% Canonical companion-paper values (../../../../Style/style.tex).
\definecolor{npgOrange}{HTML}{E69F00}
\definecolor{npgSky}{HTML}{56B4E9}
\definecolor{npgGreen}{HTML}{009E73}
\definecolor{npgYellow}{HTML}{F0E442}
\definecolor{npgBlue}{HTML}{0072B2}
\definecolor{npgVermillion}{HTML}{E64B35}
\definecolor{npgPurple}{HTML}{CC79A7}
\definecolor{figInk}{HTML}{1A1A1A}
\definecolor{figMuted}{HTML}{666666}
\definecolor{figHair}{HTML}{D9D9D9}
\definecolor{figWash}{HTML}{F7F7F7}

\colorlet{scInk}{figInk}              % primary text
\colorlet{scInkM}{figMuted}           % secondary text
\colorlet{scPaper}{figWash}           % quiet cards and note strips
\colorlet{scSpF}{npgSky!12}           % species fill
\colorlet{scSpE}{npgBlue}             % species edge
\colorlet{scRxF}{npgOrange!12}        % reaction fill
\colorlet{scRxE}{npgOrange!88!black}  % reaction edge
\colorlet{scReac}{npgBlue!78!black}   % reactant edge
\colorlet{scProd}{npgGreen!82!black}  % product edge
\colorlet{scHiF}{npgOrange!13}        % highlight fill
\colorlet{scHiE}{npgOrange!82!black}  % highlight edge
% Stage / accent tints.  stDE remains available for the fifth plot series.
\colorlet{stAF}{npgGreen!9}\colorlet{stAE}{npgGreen}
\colorlet{stBF}{npgOrange!11}\colorlet{stBE}{npgOrange}
\colorlet{stCF}{npgBlue!9}\colorlet{stCE}{npgBlue}
\colorlet{stDF}{npgPurple!10}\colorlet{stDE}{npgPurple!85!black}
\colorlet{stEF}{npgVermillion!9}\colorlet{stEE}{npgVermillion}

% Semantic aliases.  Figures should speak in these, not in stAE/stBE/...
\colorlet{scAccent}{stAE}\colorlet{scAccentF}{stAF}   % structure / normal
\colorlet{scMark}{stBE}\colorlet{scMarkF}{stBF}       % highlight / certificate
\colorlet{scRef}{stCE}\colorlet{scRefF}{stCF}         % secondary series
\colorlet{scAlert}{stEE}\colorlet{scAlertF}{stEF}     % failure / violation
\colorlet{scRule}{figHair}                            % card and axis hairlines
\colorlet{scRuleL}{figHair!55}                        % grid, dividers

% --- text styles ----------------------------------------------------------
% Three sizes, fixed in points so they do not drift with the document class.
% Arial is shared by labels, headings, panel badges, and chemistry labels.
\providecommand{\scTitle}{\scArial\fontsize{7}{8.4}\selectfont\bfseries\color{scInk}}
\providecommand{\scBody}{\scArial\fontsize{6.5}{7.8}\selectfont\color{scInkM}}
\providecommand{\scSmall}{\scArial\fontsize{6}{7.2}\selectfont\color{scInkM}}
\providecommand{\scSub}[1]{{\scSmall\itshape #1}}
\providecommand{\scPanel}[1]{{\scTitle(#1)}}
% Numerals that carry a result: ink, not colour, so the eye reads the value.
\providecommand{\scNum}{\scArial\fontsize{10}{11}\selectfont\bfseries\color{scInk}}

% --- shared TikZ styles ---------------------------------------------------
\tikzset{
  % text never hyphenates inside a figure
  every node/.append style={execute at begin node={%
      \hyphenpenalty=10000 \exhyphenpenalty=10000 \hbadness=10000 }},
  syn/.style={font=\scArial,color=scInk,
              execute at begin picture={\MTversion*{synarial}},
              >={Stealth[round,length=2.4mm,width=1.9mm]}},
  % containers ------------------------------------------------------------
  % One card chrome for the whole paper: white, 3 pt radius, one hairline.
  panel/.style={rounded corners=2.5pt,fill=none,draw=scRule,line width=0.45pt},
  card/.style={rounded corners=2.2pt,fill=scPaper,draw=scRule,line width=0.45pt,
               inner sep=5pt,align=left,anchor=north west},
  % a quiet strip for a panel's closing fact
  note/.style={rounded corners=2.5pt,fill=scPaper,draw=scRuleL,line width=0.4pt,
               inner xsep=6pt,inner ysep=3.5pt,align=left,font=\scSmall},
  % small labelled value pill
  chip/.style={rounded corners=2pt,fill=scPaper,draw=scRule,line width=0.4pt,
               inner xsep=4pt,inner ysep=2pt,font=\scSmall},
  % Circular letter badge: the same accessible panel marker used in Style/*.
  hbullet/.style={circle,fill=scInk,text=white,draw=none,
                  font=\scArial\fontsize{6}{6}\selectfont\bfseries,
                  minimum size=4.2mm,inner sep=0pt},
  ctag/.style={font=\scTitle,anchor=north west,text=scInk},
  % CRN primitives --------------------------------------------------------
  species/.style={rounded rectangle,fill=scSpF,draw=scSpE,line width=0.45pt,
                  inner xsep=2.4pt,inner ysep=1.6pt,font=\scSmall,
                  text=scInk,minimum height=3.6mm},
  cofactor/.style={species,fill=scRefF,draw=scRef},
  catalyst/.style={species,fill=scHiF,draw=scHiE},
  rxn/.style={rectangle,fill=scRxF,draw=scRxE,line width=0.45pt,
              inner sep=1.2pt,font=\scSmall,text=scInk,
              minimum width=3.0mm,minimum height=3.0mm},
  place/.style={circle,fill=scSpF,draw=scSpE,line width=0.45pt,
                minimum size=4.0mm,inner sep=0pt},
  trans/.style={rectangle,fill=scRxE,draw=scRxE,minimum width=1.1mm,
                minimum height=3.6mm,inner sep=0pt},
  token/.style={circle,fill=scInk,minimum size=1.1mm,inner sep=0pt},
  % edges -----------------------------------------------------------------
  reac/.style={->,line width=0.5pt,draw=scReac,shorten >=0.8pt,shorten <=0.8pt},
  prod/.style={->,line width=0.5pt,draw=scProd,shorten >=0.8pt,shorten <=0.8pt},
  cat/.style={->,line width=0.5pt,draw=scHiE,shorten >=0.8pt,shorten <=0.8pt},
  flow/.style={-{Stealth[round,length=2.2mm,width=1.8mm]},line width=1.1pt,
               draw=scInkM!45},
  flab/.style={font=\scSmall,text=scInkM,inner sep=2pt},
  stoi/.style={font=\scSmall,text=scInkM,inner sep=0.8pt},
  % verdict marks: shape carries the meaning, colour only reinforces it
  vpass/.style={circle,fill=scAccentF,draw=scAccent,text=scAccent,
                font=\fontsize{6}{6}\selectfont\bfseries,
                minimum size=3.2mm,inner sep=0pt},
  vfail/.style={circle,fill=scAlertF,draw=scAlert,text=scAlert,
                font=\fontsize{6}{6}\selectfont\bfseries,
                minimum size=3.2mm,inner sep=0pt},
  callout/.style={font=\scSmall,text=scInkM,align=left},
}

% \scHead[<badge colour>]{<anchor coord>}{<letter>}{<title>}{<subtitle or empty>}
% One header treatment for every panel in the paper.
\providecommand{\scHead}[5][scInk]{%
  \node[hbullet,fill=#1,anchor=north west] (scbadge#3) at ($#2+(0,-0.015)$)
    {\MakeUppercase{#3}};
  \node[ctag,anchor=west] at ($(scbadge#3.east)+(0.14,0)$) {#4};
  \ifx\relax#5\relax\else
    \node[anchor=north west,font=\scSmall,text=scInkM]
      at ($#2+(0.56,-0.245)$) {#5};
  \fi}

% --- shared pgfplots look (only if the figure loaded pgfplots) ------------
\ifdefined\pgfplotsset
  \pgfplotsset{compat=1.18}
  \pgfplotsset{
    synaxis/.style={
      scale only axis,
      tick align=outside, tick pos=left,
      axis line style={draw=scInkM!55,line width=0.4pt},
      % One grid, on the value axis only.  No minor grid: at these sizes it
      % turns the plot frame into a grey wash.
      ymajorgrids=true, xmajorgrids=false,
      major grid style={draw=scRuleL,line width=0.3pt},
      major tick length=1.6pt,
      label style={font=\scBody,color=scInk},
      tick label style={font=\scSmall,color=scInkM},
      legend style={font=\scSmall,draw=none,fill=none,inner sep=1pt,
                    row sep=0.4pt,legend columns=1},
      legend cell align=left,
      every axis plot/.append style={line width=0.7pt,mark size=1.2pt},
    },
  }
\fi
}{% ---------------------------------------------------------------------------
% synstyle.tex --- shared look for every TikZ figure of the SynKit CRN paper.
%
% Input it from a figure file:
%     \IfFileExists{synstyle.tex}{% ---------------------------------------------------------------------------
% synstyle.tex --- shared look for every TikZ figure of the SynKit CRN paper.
%
% Input it from a figure file:
%     \IfFileExists{synstyle.tex}{\input{synstyle}}{\input{Figure/synstyle}}
% so that the figure compiles both from Figure/ and from the paper root.
%
% The visual vocabulary follows ../../../../Style/style.tex: the
% colour-blind-safe Okabe--Ito accents, quiet neutral rules and wash, circular
% panel badges, and a compact technical type scale are shared by the companion
% papers.  Semantic aliases below keep the CRN meaning explicit.
%
% DESIGN CONTRACT -----------------------------------------------------------
% 1. Geometry.  Every figure is laid out on a canvas exactly \scW wide (the
%    sn-jnl text width, 372 pt) and is included unscaled, so one point of type
%    is the same physical size in all five figures.  Open each picture with
%    \scCanvas{<height in cm>} to pin the bounding box.
% 2. Colour is semantic, never decorative.  Four roles only:
%       scAccent (green)       structure, the normal/expected path
%       scMark   (orange)      catalyst, certificate, or highlight
%       scAlert  (vermillion)  failure, violation, or uncovered set
%       scRef    (blue)        secondary or reference series
%    A panel gets at most one accent.  Titles are always ink; panel identity
%    is also encoded by a letter inside a circular badge, not by colour alone.
% 3. Type scale.  Three sizes, no others, and nothing below 6 pt:
%       \scTitle (7 pt bold)  panel titles
%       \scBody  (6.5 pt)     labels, values, in-figure prose
%       \scSmall (6 pt)       units, footnotes, dense graph labels
% ---------------------------------------------------------------------------
% Load guard.  The paper inputs every figure source inline, so this file is
% reached once per figure; without the guard the second \newlength below
% would abort the run.
\ifdefined\SynKitStyleLoaded\endinput\fi
\newcommand{\SynKitStyleLoaded}{}

% Load in the document preamble. XeLaTeX/LuaLaTeX and an installed Arial
% family are required; fontspec reports an error instead of substituting.
\RequirePackage[no-math]{fontspec}
\newfontfamily\scArial{Arial}[NFSSFamily=scArial,Ligatures=TeX]
% Keep manuscript mathematics unchanged; select Arial letters and numerals
% only inside figures. TeX supplies specialist mathematical symbols.
\RequirePackage[subdued,italic]{mathastext}
\MTDeclareVersion{synarial}{TU}{scArial}{m}{n}
\SetMathAlphabet{\Mathsf}{synarial}{TU}{scArial}{m}{n}
\SetMathAlphabet{\Mathtt}{synarial}{TU}{scArial}{m}{n}

\usetikzlibrary{positioning,calc,fit,backgrounds,arrows.meta,shapes.geometric,
                shapes.misc,decorations.pathreplacing,patterns,math}

% --- canvas ---------------------------------------------------------------
% 372 pt text width minus the 1 pt standalone border on each side.
\newlength{\scW}\setlength{\scW}{13.0326cm}
\providecommand{\scWnum}{13.0326}
% \scCanvas{<height>} pins the picture to the full text width, y running down
% from 0, so every figure shares one coordinate convention.  It *sets* the
% bounding box rather than contributing to it, so the emitted page is exactly
% the text width and every figure is included at 1:1 -- otherwise a stray
% invisible extent (pgfplots reserves a little around an axis) rescales one
% figure relative to the others and the same \scSmall prints at two sizes.
% Anything drawn outside the canvas is therefore cropped: keep content inside,
% and check with
%     \providecommand{\scCanvas}[1]{\path[draw=red] (0,0) rectangle (\scWnum,-#1);}
\providecommand{\scCanvas}[1]{%
  \path[use as bounding box] (0,0) rectangle (\scWnum,-#1);}

% --- palette --------------------------------------------------------------
% Canonical companion-paper values (../../../../Style/style.tex).
\definecolor{npgOrange}{HTML}{E69F00}
\definecolor{npgSky}{HTML}{56B4E9}
\definecolor{npgGreen}{HTML}{009E73}
\definecolor{npgYellow}{HTML}{F0E442}
\definecolor{npgBlue}{HTML}{0072B2}
\definecolor{npgVermillion}{HTML}{E64B35}
\definecolor{npgPurple}{HTML}{CC79A7}
\definecolor{figInk}{HTML}{1A1A1A}
\definecolor{figMuted}{HTML}{666666}
\definecolor{figHair}{HTML}{D9D9D9}
\definecolor{figWash}{HTML}{F7F7F7}

\colorlet{scInk}{figInk}              % primary text
\colorlet{scInkM}{figMuted}           % secondary text
\colorlet{scPaper}{figWash}           % quiet cards and note strips
\colorlet{scSpF}{npgSky!12}           % species fill
\colorlet{scSpE}{npgBlue}             % species edge
\colorlet{scRxF}{npgOrange!12}        % reaction fill
\colorlet{scRxE}{npgOrange!88!black}  % reaction edge
\colorlet{scReac}{npgBlue!78!black}   % reactant edge
\colorlet{scProd}{npgGreen!82!black}  % product edge
\colorlet{scHiF}{npgOrange!13}        % highlight fill
\colorlet{scHiE}{npgOrange!82!black}  % highlight edge
% Stage / accent tints.  stDE remains available for the fifth plot series.
\colorlet{stAF}{npgGreen!9}\colorlet{stAE}{npgGreen}
\colorlet{stBF}{npgOrange!11}\colorlet{stBE}{npgOrange}
\colorlet{stCF}{npgBlue!9}\colorlet{stCE}{npgBlue}
\colorlet{stDF}{npgPurple!10}\colorlet{stDE}{npgPurple!85!black}
\colorlet{stEF}{npgVermillion!9}\colorlet{stEE}{npgVermillion}

% Semantic aliases.  Figures should speak in these, not in stAE/stBE/...
\colorlet{scAccent}{stAE}\colorlet{scAccentF}{stAF}   % structure / normal
\colorlet{scMark}{stBE}\colorlet{scMarkF}{stBF}       % highlight / certificate
\colorlet{scRef}{stCE}\colorlet{scRefF}{stCF}         % secondary series
\colorlet{scAlert}{stEE}\colorlet{scAlertF}{stEF}     % failure / violation
\colorlet{scRule}{figHair}                            % card and axis hairlines
\colorlet{scRuleL}{figHair!55}                        % grid, dividers

% --- text styles ----------------------------------------------------------
% Three sizes, fixed in points so they do not drift with the document class.
% Arial is shared by labels, headings, panel badges, and chemistry labels.
\providecommand{\scTitle}{\scArial\fontsize{7}{8.4}\selectfont\bfseries\color{scInk}}
\providecommand{\scBody}{\scArial\fontsize{6.5}{7.8}\selectfont\color{scInkM}}
\providecommand{\scSmall}{\scArial\fontsize{6}{7.2}\selectfont\color{scInkM}}
\providecommand{\scSub}[1]{{\scSmall\itshape #1}}
\providecommand{\scPanel}[1]{{\scTitle(#1)}}
% Numerals that carry a result: ink, not colour, so the eye reads the value.
\providecommand{\scNum}{\scArial\fontsize{10}{11}\selectfont\bfseries\color{scInk}}

% --- shared TikZ styles ---------------------------------------------------
\tikzset{
  % text never hyphenates inside a figure
  every node/.append style={execute at begin node={%
      \hyphenpenalty=10000 \exhyphenpenalty=10000 \hbadness=10000 }},
  syn/.style={font=\scArial,color=scInk,
              execute at begin picture={\MTversion*{synarial}},
              >={Stealth[round,length=2.4mm,width=1.9mm]}},
  % containers ------------------------------------------------------------
  % One card chrome for the whole paper: white, 3 pt radius, one hairline.
  panel/.style={rounded corners=2.5pt,fill=none,draw=scRule,line width=0.45pt},
  card/.style={rounded corners=2.2pt,fill=scPaper,draw=scRule,line width=0.45pt,
               inner sep=5pt,align=left,anchor=north west},
  % a quiet strip for a panel's closing fact
  note/.style={rounded corners=2.5pt,fill=scPaper,draw=scRuleL,line width=0.4pt,
               inner xsep=6pt,inner ysep=3.5pt,align=left,font=\scSmall},
  % small labelled value pill
  chip/.style={rounded corners=2pt,fill=scPaper,draw=scRule,line width=0.4pt,
               inner xsep=4pt,inner ysep=2pt,font=\scSmall},
  % Circular letter badge: the same accessible panel marker used in Style/*.
  hbullet/.style={circle,fill=scInk,text=white,draw=none,
                  font=\scArial\fontsize{6}{6}\selectfont\bfseries,
                  minimum size=4.2mm,inner sep=0pt},
  ctag/.style={font=\scTitle,anchor=north west,text=scInk},
  % CRN primitives --------------------------------------------------------
  species/.style={rounded rectangle,fill=scSpF,draw=scSpE,line width=0.45pt,
                  inner xsep=2.4pt,inner ysep=1.6pt,font=\scSmall,
                  text=scInk,minimum height=3.6mm},
  cofactor/.style={species,fill=scRefF,draw=scRef},
  catalyst/.style={species,fill=scHiF,draw=scHiE},
  rxn/.style={rectangle,fill=scRxF,draw=scRxE,line width=0.45pt,
              inner sep=1.2pt,font=\scSmall,text=scInk,
              minimum width=3.0mm,minimum height=3.0mm},
  place/.style={circle,fill=scSpF,draw=scSpE,line width=0.45pt,
                minimum size=4.0mm,inner sep=0pt},
  trans/.style={rectangle,fill=scRxE,draw=scRxE,minimum width=1.1mm,
                minimum height=3.6mm,inner sep=0pt},
  token/.style={circle,fill=scInk,minimum size=1.1mm,inner sep=0pt},
  % edges -----------------------------------------------------------------
  reac/.style={->,line width=0.5pt,draw=scReac,shorten >=0.8pt,shorten <=0.8pt},
  prod/.style={->,line width=0.5pt,draw=scProd,shorten >=0.8pt,shorten <=0.8pt},
  cat/.style={->,line width=0.5pt,draw=scHiE,shorten >=0.8pt,shorten <=0.8pt},
  flow/.style={-{Stealth[round,length=2.2mm,width=1.8mm]},line width=1.1pt,
               draw=scInkM!45},
  flab/.style={font=\scSmall,text=scInkM,inner sep=2pt},
  stoi/.style={font=\scSmall,text=scInkM,inner sep=0.8pt},
  % verdict marks: shape carries the meaning, colour only reinforces it
  vpass/.style={circle,fill=scAccentF,draw=scAccent,text=scAccent,
                font=\fontsize{6}{6}\selectfont\bfseries,
                minimum size=3.2mm,inner sep=0pt},
  vfail/.style={circle,fill=scAlertF,draw=scAlert,text=scAlert,
                font=\fontsize{6}{6}\selectfont\bfseries,
                minimum size=3.2mm,inner sep=0pt},
  callout/.style={font=\scSmall,text=scInkM,align=left},
}

% \scHead[<badge colour>]{<anchor coord>}{<letter>}{<title>}{<subtitle or empty>}
% One header treatment for every panel in the paper.
\providecommand{\scHead}[5][scInk]{%
  \node[hbullet,fill=#1,anchor=north west] (scbadge#3) at ($#2+(0,-0.015)$)
    {\MakeUppercase{#3}};
  \node[ctag,anchor=west] at ($(scbadge#3.east)+(0.14,0)$) {#4};
  \ifx\relax#5\relax\else
    \node[anchor=north west,font=\scSmall,text=scInkM]
      at ($#2+(0.56,-0.245)$) {#5};
  \fi}

% --- shared pgfplots look (only if the figure loaded pgfplots) ------------
\ifdefined\pgfplotsset
  \pgfplotsset{compat=1.18}
  \pgfplotsset{
    synaxis/.style={
      scale only axis,
      tick align=outside, tick pos=left,
      axis line style={draw=scInkM!55,line width=0.4pt},
      % One grid, on the value axis only.  No minor grid: at these sizes it
      % turns the plot frame into a grey wash.
      ymajorgrids=true, xmajorgrids=false,
      major grid style={draw=scRuleL,line width=0.3pt},
      major tick length=1.6pt,
      label style={font=\scBody,color=scInk},
      tick label style={font=\scSmall,color=scInkM},
      legend style={font=\scSmall,draw=none,fill=none,inner sep=1pt,
                    row sep=0.4pt,legend columns=1},
      legend cell align=left,
      every axis plot/.append style={line width=0.7pt,mark size=1.2pt},
    },
  }
\fi
}{% ---------------------------------------------------------------------------
% synstyle.tex --- shared look for every TikZ figure of the SynKit CRN paper.
%
% Input it from a figure file:
%     \IfFileExists{synstyle.tex}{\input{synstyle}}{\input{Figure/synstyle}}
% so that the figure compiles both from Figure/ and from the paper root.
%
% The visual vocabulary follows ../../../../Style/style.tex: the
% colour-blind-safe Okabe--Ito accents, quiet neutral rules and wash, circular
% panel badges, and a compact technical type scale are shared by the companion
% papers.  Semantic aliases below keep the CRN meaning explicit.
%
% DESIGN CONTRACT -----------------------------------------------------------
% 1. Geometry.  Every figure is laid out on a canvas exactly \scW wide (the
%    sn-jnl text width, 372 pt) and is included unscaled, so one point of type
%    is the same physical size in all five figures.  Open each picture with
%    \scCanvas{<height in cm>} to pin the bounding box.
% 2. Colour is semantic, never decorative.  Four roles only:
%       scAccent (green)       structure, the normal/expected path
%       scMark   (orange)      catalyst, certificate, or highlight
%       scAlert  (vermillion)  failure, violation, or uncovered set
%       scRef    (blue)        secondary or reference series
%    A panel gets at most one accent.  Titles are always ink; panel identity
%    is also encoded by a letter inside a circular badge, not by colour alone.
% 3. Type scale.  Three sizes, no others, and nothing below 6 pt:
%       \scTitle (7 pt bold)  panel titles
%       \scBody  (6.5 pt)     labels, values, in-figure prose
%       \scSmall (6 pt)       units, footnotes, dense graph labels
% ---------------------------------------------------------------------------
% Load guard.  The paper inputs every figure source inline, so this file is
% reached once per figure; without the guard the second \newlength below
% would abort the run.
\ifdefined\SynKitStyleLoaded\endinput\fi
\newcommand{\SynKitStyleLoaded}{}

% Load in the document preamble. XeLaTeX/LuaLaTeX and an installed Arial
% family are required; fontspec reports an error instead of substituting.
\RequirePackage[no-math]{fontspec}
\newfontfamily\scArial{Arial}[NFSSFamily=scArial,Ligatures=TeX]
% Keep manuscript mathematics unchanged; select Arial letters and numerals
% only inside figures. TeX supplies specialist mathematical symbols.
\RequirePackage[subdued,italic]{mathastext}
\MTDeclareVersion{synarial}{TU}{scArial}{m}{n}
\SetMathAlphabet{\Mathsf}{synarial}{TU}{scArial}{m}{n}
\SetMathAlphabet{\Mathtt}{synarial}{TU}{scArial}{m}{n}

\usetikzlibrary{positioning,calc,fit,backgrounds,arrows.meta,shapes.geometric,
                shapes.misc,decorations.pathreplacing,patterns,math}

% --- canvas ---------------------------------------------------------------
% 372 pt text width minus the 1 pt standalone border on each side.
\newlength{\scW}\setlength{\scW}{13.0326cm}
\providecommand{\scWnum}{13.0326}
% \scCanvas{<height>} pins the picture to the full text width, y running down
% from 0, so every figure shares one coordinate convention.  It *sets* the
% bounding box rather than contributing to it, so the emitted page is exactly
% the text width and every figure is included at 1:1 -- otherwise a stray
% invisible extent (pgfplots reserves a little around an axis) rescales one
% figure relative to the others and the same \scSmall prints at two sizes.
% Anything drawn outside the canvas is therefore cropped: keep content inside,
% and check with
%     \providecommand{\scCanvas}[1]{\path[draw=red] (0,0) rectangle (\scWnum,-#1);}
\providecommand{\scCanvas}[1]{%
  \path[use as bounding box] (0,0) rectangle (\scWnum,-#1);}

% --- palette --------------------------------------------------------------
% Canonical companion-paper values (../../../../Style/style.tex).
\definecolor{npgOrange}{HTML}{E69F00}
\definecolor{npgSky}{HTML}{56B4E9}
\definecolor{npgGreen}{HTML}{009E73}
\definecolor{npgYellow}{HTML}{F0E442}
\definecolor{npgBlue}{HTML}{0072B2}
\definecolor{npgVermillion}{HTML}{E64B35}
\definecolor{npgPurple}{HTML}{CC79A7}
\definecolor{figInk}{HTML}{1A1A1A}
\definecolor{figMuted}{HTML}{666666}
\definecolor{figHair}{HTML}{D9D9D9}
\definecolor{figWash}{HTML}{F7F7F7}

\colorlet{scInk}{figInk}              % primary text
\colorlet{scInkM}{figMuted}           % secondary text
\colorlet{scPaper}{figWash}           % quiet cards and note strips
\colorlet{scSpF}{npgSky!12}           % species fill
\colorlet{scSpE}{npgBlue}             % species edge
\colorlet{scRxF}{npgOrange!12}        % reaction fill
\colorlet{scRxE}{npgOrange!88!black}  % reaction edge
\colorlet{scReac}{npgBlue!78!black}   % reactant edge
\colorlet{scProd}{npgGreen!82!black}  % product edge
\colorlet{scHiF}{npgOrange!13}        % highlight fill
\colorlet{scHiE}{npgOrange!82!black}  % highlight edge
% Stage / accent tints.  stDE remains available for the fifth plot series.
\colorlet{stAF}{npgGreen!9}\colorlet{stAE}{npgGreen}
\colorlet{stBF}{npgOrange!11}\colorlet{stBE}{npgOrange}
\colorlet{stCF}{npgBlue!9}\colorlet{stCE}{npgBlue}
\colorlet{stDF}{npgPurple!10}\colorlet{stDE}{npgPurple!85!black}
\colorlet{stEF}{npgVermillion!9}\colorlet{stEE}{npgVermillion}

% Semantic aliases.  Figures should speak in these, not in stAE/stBE/...
\colorlet{scAccent}{stAE}\colorlet{scAccentF}{stAF}   % structure / normal
\colorlet{scMark}{stBE}\colorlet{scMarkF}{stBF}       % highlight / certificate
\colorlet{scRef}{stCE}\colorlet{scRefF}{stCF}         % secondary series
\colorlet{scAlert}{stEE}\colorlet{scAlertF}{stEF}     % failure / violation
\colorlet{scRule}{figHair}                            % card and axis hairlines
\colorlet{scRuleL}{figHair!55}                        % grid, dividers

% --- text styles ----------------------------------------------------------
% Three sizes, fixed in points so they do not drift with the document class.
% Arial is shared by labels, headings, panel badges, and chemistry labels.
\providecommand{\scTitle}{\scArial\fontsize{7}{8.4}\selectfont\bfseries\color{scInk}}
\providecommand{\scBody}{\scArial\fontsize{6.5}{7.8}\selectfont\color{scInkM}}
\providecommand{\scSmall}{\scArial\fontsize{6}{7.2}\selectfont\color{scInkM}}
\providecommand{\scSub}[1]{{\scSmall\itshape #1}}
\providecommand{\scPanel}[1]{{\scTitle(#1)}}
% Numerals that carry a result: ink, not colour, so the eye reads the value.
\providecommand{\scNum}{\scArial\fontsize{10}{11}\selectfont\bfseries\color{scInk}}

% --- shared TikZ styles ---------------------------------------------------
\tikzset{
  % text never hyphenates inside a figure
  every node/.append style={execute at begin node={%
      \hyphenpenalty=10000 \exhyphenpenalty=10000 \hbadness=10000 }},
  syn/.style={font=\scArial,color=scInk,
              execute at begin picture={\MTversion*{synarial}},
              >={Stealth[round,length=2.4mm,width=1.9mm]}},
  % containers ------------------------------------------------------------
  % One card chrome for the whole paper: white, 3 pt radius, one hairline.
  panel/.style={rounded corners=2.5pt,fill=none,draw=scRule,line width=0.45pt},
  card/.style={rounded corners=2.2pt,fill=scPaper,draw=scRule,line width=0.45pt,
               inner sep=5pt,align=left,anchor=north west},
  % a quiet strip for a panel's closing fact
  note/.style={rounded corners=2.5pt,fill=scPaper,draw=scRuleL,line width=0.4pt,
               inner xsep=6pt,inner ysep=3.5pt,align=left,font=\scSmall},
  % small labelled value pill
  chip/.style={rounded corners=2pt,fill=scPaper,draw=scRule,line width=0.4pt,
               inner xsep=4pt,inner ysep=2pt,font=\scSmall},
  % Circular letter badge: the same accessible panel marker used in Style/*.
  hbullet/.style={circle,fill=scInk,text=white,draw=none,
                  font=\scArial\fontsize{6}{6}\selectfont\bfseries,
                  minimum size=4.2mm,inner sep=0pt},
  ctag/.style={font=\scTitle,anchor=north west,text=scInk},
  % CRN primitives --------------------------------------------------------
  species/.style={rounded rectangle,fill=scSpF,draw=scSpE,line width=0.45pt,
                  inner xsep=2.4pt,inner ysep=1.6pt,font=\scSmall,
                  text=scInk,minimum height=3.6mm},
  cofactor/.style={species,fill=scRefF,draw=scRef},
  catalyst/.style={species,fill=scHiF,draw=scHiE},
  rxn/.style={rectangle,fill=scRxF,draw=scRxE,line width=0.45pt,
              inner sep=1.2pt,font=\scSmall,text=scInk,
              minimum width=3.0mm,minimum height=3.0mm},
  place/.style={circle,fill=scSpF,draw=scSpE,line width=0.45pt,
                minimum size=4.0mm,inner sep=0pt},
  trans/.style={rectangle,fill=scRxE,draw=scRxE,minimum width=1.1mm,
                minimum height=3.6mm,inner sep=0pt},
  token/.style={circle,fill=scInk,minimum size=1.1mm,inner sep=0pt},
  % edges -----------------------------------------------------------------
  reac/.style={->,line width=0.5pt,draw=scReac,shorten >=0.8pt,shorten <=0.8pt},
  prod/.style={->,line width=0.5pt,draw=scProd,shorten >=0.8pt,shorten <=0.8pt},
  cat/.style={->,line width=0.5pt,draw=scHiE,shorten >=0.8pt,shorten <=0.8pt},
  flow/.style={-{Stealth[round,length=2.2mm,width=1.8mm]},line width=1.1pt,
               draw=scInkM!45},
  flab/.style={font=\scSmall,text=scInkM,inner sep=2pt},
  stoi/.style={font=\scSmall,text=scInkM,inner sep=0.8pt},
  % verdict marks: shape carries the meaning, colour only reinforces it
  vpass/.style={circle,fill=scAccentF,draw=scAccent,text=scAccent,
                font=\fontsize{6}{6}\selectfont\bfseries,
                minimum size=3.2mm,inner sep=0pt},
  vfail/.style={circle,fill=scAlertF,draw=scAlert,text=scAlert,
                font=\fontsize{6}{6}\selectfont\bfseries,
                minimum size=3.2mm,inner sep=0pt},
  callout/.style={font=\scSmall,text=scInkM,align=left},
}

% \scHead[<badge colour>]{<anchor coord>}{<letter>}{<title>}{<subtitle or empty>}
% One header treatment for every panel in the paper.
\providecommand{\scHead}[5][scInk]{%
  \node[hbullet,fill=#1,anchor=north west] (scbadge#3) at ($#2+(0,-0.015)$)
    {\MakeUppercase{#3}};
  \node[ctag,anchor=west] at ($(scbadge#3.east)+(0.14,0)$) {#4};
  \ifx\relax#5\relax\else
    \node[anchor=north west,font=\scSmall,text=scInkM]
      at ($#2+(0.56,-0.245)$) {#5};
  \fi}

% --- shared pgfplots look (only if the figure loaded pgfplots) ------------
\ifdefined\pgfplotsset
  \pgfplotsset{compat=1.18}
  \pgfplotsset{
    synaxis/.style={
      scale only axis,
      tick align=outside, tick pos=left,
      axis line style={draw=scInkM!55,line width=0.4pt},
      % One grid, on the value axis only.  No minor grid: at these sizes it
      % turns the plot frame into a grey wash.
      ymajorgrids=true, xmajorgrids=false,
      major grid style={draw=scRuleL,line width=0.3pt},
      major tick length=1.6pt,
      label style={font=\scBody,color=scInk},
      tick label style={font=\scSmall,color=scInkM},
      legend style={font=\scSmall,draw=none,fill=none,inner sep=1pt,
                    row sep=0.4pt,legend columns=1},
      legend cell align=left,
      every axis plot/.append style={line width=0.7pt,mark size=1.2pt},
    },
  }
\fi
}
% so that the figure compiles both from Figure/ and from the paper root.
%
% The visual vocabulary follows ../../../../Style/style.tex: the
% colour-blind-safe Okabe--Ito accents, quiet neutral rules and wash, circular
% panel badges, and a compact technical type scale are shared by the companion
% papers.  Semantic aliases below keep the CRN meaning explicit.
%
% DESIGN CONTRACT -----------------------------------------------------------
% 1. Geometry.  Every figure is laid out on a canvas exactly \scW wide (the
%    sn-jnl text width, 372 pt) and is included unscaled, so one point of type
%    is the same physical size in all five figures.  Open each picture with
%    \scCanvas{<height in cm>} to pin the bounding box.
% 2. Colour is semantic, never decorative.  Four roles only:
%       scAccent (green)       structure, the normal/expected path
%       scMark   (orange)      catalyst, certificate, or highlight
%       scAlert  (vermillion)  failure, violation, or uncovered set
%       scRef    (blue)        secondary or reference series
%    A panel gets at most one accent.  Titles are always ink; panel identity
%    is also encoded by a letter inside a circular badge, not by colour alone.
% 3. Type scale.  Three sizes, no others, and nothing below 6 pt:
%       \scTitle (7 pt bold)  panel titles
%       \scBody  (6.5 pt)     labels, values, in-figure prose
%       \scSmall (6 pt)       units, footnotes, dense graph labels
% ---------------------------------------------------------------------------
% Load guard.  The paper inputs every figure source inline, so this file is
% reached once per figure; without the guard the second \newlength below
% would abort the run.
\ifdefined\SynKitStyleLoaded\endinput\fi
\newcommand{\SynKitStyleLoaded}{}

% Load in the document preamble. XeLaTeX/LuaLaTeX and an installed Arial
% family are required; fontspec reports an error instead of substituting.
\RequirePackage[no-math]{fontspec}
\newfontfamily\scArial{Arial}[NFSSFamily=scArial,Ligatures=TeX]
% Keep manuscript mathematics unchanged; select Arial letters and numerals
% only inside figures. TeX supplies specialist mathematical symbols.
\RequirePackage[subdued,italic]{mathastext}
\MTDeclareVersion{synarial}{TU}{scArial}{m}{n}
\SetMathAlphabet{\Mathsf}{synarial}{TU}{scArial}{m}{n}
\SetMathAlphabet{\Mathtt}{synarial}{TU}{scArial}{m}{n}

\usetikzlibrary{positioning,calc,fit,backgrounds,arrows.meta,shapes.geometric,
                shapes.misc,decorations.pathreplacing,patterns,math}

% --- canvas ---------------------------------------------------------------
% 372 pt text width minus the 1 pt standalone border on each side.
\newlength{\scW}\setlength{\scW}{13.0326cm}
\providecommand{\scWnum}{13.0326}
% \scCanvas{<height>} pins the picture to the full text width, y running down
% from 0, so every figure shares one coordinate convention.  It *sets* the
% bounding box rather than contributing to it, so the emitted page is exactly
% the text width and every figure is included at 1:1 -- otherwise a stray
% invisible extent (pgfplots reserves a little around an axis) rescales one
% figure relative to the others and the same \scSmall prints at two sizes.
% Anything drawn outside the canvas is therefore cropped: keep content inside,
% and check with
%     \providecommand{\scCanvas}[1]{\path[draw=red] (0,0) rectangle (\scWnum,-#1);}
\providecommand{\scCanvas}[1]{%
  \path[use as bounding box] (0,0) rectangle (\scWnum,-#1);}

% --- palette --------------------------------------------------------------
% Canonical companion-paper values (../../../../Style/style.tex).
\definecolor{npgOrange}{HTML}{E69F00}
\definecolor{npgSky}{HTML}{56B4E9}
\definecolor{npgGreen}{HTML}{009E73}
\definecolor{npgYellow}{HTML}{F0E442}
\definecolor{npgBlue}{HTML}{0072B2}
\definecolor{npgVermillion}{HTML}{E64B35}
\definecolor{npgPurple}{HTML}{CC79A7}
\definecolor{figInk}{HTML}{1A1A1A}
\definecolor{figMuted}{HTML}{666666}
\definecolor{figHair}{HTML}{D9D9D9}
\definecolor{figWash}{HTML}{F7F7F7}

\colorlet{scInk}{figInk}              % primary text
\colorlet{scInkM}{figMuted}           % secondary text
\colorlet{scPaper}{figWash}           % quiet cards and note strips
\colorlet{scSpF}{npgSky!12}           % species fill
\colorlet{scSpE}{npgBlue}             % species edge
\colorlet{scRxF}{npgOrange!12}        % reaction fill
\colorlet{scRxE}{npgOrange!88!black}  % reaction edge
\colorlet{scReac}{npgBlue!78!black}   % reactant edge
\colorlet{scProd}{npgGreen!82!black}  % product edge
\colorlet{scHiF}{npgOrange!13}        % highlight fill
\colorlet{scHiE}{npgOrange!82!black}  % highlight edge
% Stage / accent tints.  stDE remains available for the fifth plot series.
\colorlet{stAF}{npgGreen!9}\colorlet{stAE}{npgGreen}
\colorlet{stBF}{npgOrange!11}\colorlet{stBE}{npgOrange}
\colorlet{stCF}{npgBlue!9}\colorlet{stCE}{npgBlue}
\colorlet{stDF}{npgPurple!10}\colorlet{stDE}{npgPurple!85!black}
\colorlet{stEF}{npgVermillion!9}\colorlet{stEE}{npgVermillion}

% Semantic aliases.  Figures should speak in these, not in stAE/stBE/...
\colorlet{scAccent}{stAE}\colorlet{scAccentF}{stAF}   % structure / normal
\colorlet{scMark}{stBE}\colorlet{scMarkF}{stBF}       % highlight / certificate
\colorlet{scRef}{stCE}\colorlet{scRefF}{stCF}         % secondary series
\colorlet{scAlert}{stEE}\colorlet{scAlertF}{stEF}     % failure / violation
\colorlet{scRule}{figHair}                            % card and axis hairlines
\colorlet{scRuleL}{figHair!55}                        % grid, dividers

% --- text styles ----------------------------------------------------------
% Three sizes, fixed in points so they do not drift with the document class.
% Arial is shared by labels, headings, panel badges, and chemistry labels.
\providecommand{\scTitle}{\scArial\fontsize{7}{8.4}\selectfont\bfseries\color{scInk}}
\providecommand{\scBody}{\scArial\fontsize{6.5}{7.8}\selectfont\color{scInkM}}
\providecommand{\scSmall}{\scArial\fontsize{6}{7.2}\selectfont\color{scInkM}}
\providecommand{\scSub}[1]{{\scSmall\itshape #1}}
\providecommand{\scPanel}[1]{{\scTitle(#1)}}
% Numerals that carry a result: ink, not colour, so the eye reads the value.
\providecommand{\scNum}{\scArial\fontsize{10}{11}\selectfont\bfseries\color{scInk}}

% --- shared TikZ styles ---------------------------------------------------
\tikzset{
  % text never hyphenates inside a figure
  every node/.append style={execute at begin node={%
      \hyphenpenalty=10000 \exhyphenpenalty=10000 \hbadness=10000 }},
  syn/.style={font=\scArial,color=scInk,
              execute at begin picture={\MTversion*{synarial}},
              >={Stealth[round,length=2.4mm,width=1.9mm]}},
  % containers ------------------------------------------------------------
  % One card chrome for the whole paper: white, 3 pt radius, one hairline.
  panel/.style={rounded corners=2.5pt,fill=none,draw=scRule,line width=0.45pt},
  card/.style={rounded corners=2.2pt,fill=scPaper,draw=scRule,line width=0.45pt,
               inner sep=5pt,align=left,anchor=north west},
  % a quiet strip for a panel's closing fact
  note/.style={rounded corners=2.5pt,fill=scPaper,draw=scRuleL,line width=0.4pt,
               inner xsep=6pt,inner ysep=3.5pt,align=left,font=\scSmall},
  % small labelled value pill
  chip/.style={rounded corners=2pt,fill=scPaper,draw=scRule,line width=0.4pt,
               inner xsep=4pt,inner ysep=2pt,font=\scSmall},
  % Circular letter badge: the same accessible panel marker used in Style/*.
  hbullet/.style={circle,fill=scInk,text=white,draw=none,
                  font=\scArial\fontsize{6}{6}\selectfont\bfseries,
                  minimum size=4.2mm,inner sep=0pt},
  ctag/.style={font=\scTitle,anchor=north west,text=scInk},
  % CRN primitives --------------------------------------------------------
  species/.style={rounded rectangle,fill=scSpF,draw=scSpE,line width=0.45pt,
                  inner xsep=2.4pt,inner ysep=1.6pt,font=\scSmall,
                  text=scInk,minimum height=3.6mm},
  cofactor/.style={species,fill=scRefF,draw=scRef},
  catalyst/.style={species,fill=scHiF,draw=scHiE},
  rxn/.style={rectangle,fill=scRxF,draw=scRxE,line width=0.45pt,
              inner sep=1.2pt,font=\scSmall,text=scInk,
              minimum width=3.0mm,minimum height=3.0mm},
  place/.style={circle,fill=scSpF,draw=scSpE,line width=0.45pt,
                minimum size=4.0mm,inner sep=0pt},
  trans/.style={rectangle,fill=scRxE,draw=scRxE,minimum width=1.1mm,
                minimum height=3.6mm,inner sep=0pt},
  token/.style={circle,fill=scInk,minimum size=1.1mm,inner sep=0pt},
  % edges -----------------------------------------------------------------
  reac/.style={->,line width=0.5pt,draw=scReac,shorten >=0.8pt,shorten <=0.8pt},
  prod/.style={->,line width=0.5pt,draw=scProd,shorten >=0.8pt,shorten <=0.8pt},
  cat/.style={->,line width=0.5pt,draw=scHiE,shorten >=0.8pt,shorten <=0.8pt},
  flow/.style={-{Stealth[round,length=2.2mm,width=1.8mm]},line width=1.1pt,
               draw=scInkM!45},
  flab/.style={font=\scSmall,text=scInkM,inner sep=2pt},
  stoi/.style={font=\scSmall,text=scInkM,inner sep=0.8pt},
  % verdict marks: shape carries the meaning, colour only reinforces it
  vpass/.style={circle,fill=scAccentF,draw=scAccent,text=scAccent,
                font=\fontsize{6}{6}\selectfont\bfseries,
                minimum size=3.2mm,inner sep=0pt},
  vfail/.style={circle,fill=scAlertF,draw=scAlert,text=scAlert,
                font=\fontsize{6}{6}\selectfont\bfseries,
                minimum size=3.2mm,inner sep=0pt},
  callout/.style={font=\scSmall,text=scInkM,align=left},
}

% \scHead[<badge colour>]{<anchor coord>}{<letter>}{<title>}{<subtitle or empty>}
% One header treatment for every panel in the paper.
\providecommand{\scHead}[5][scInk]{%
  \node[hbullet,fill=#1,anchor=north west] (scbadge#3) at ($#2+(0,-0.015)$)
    {\MakeUppercase{#3}};
  \node[ctag,anchor=west] at ($(scbadge#3.east)+(0.14,0)$) {#4};
  \ifx\relax#5\relax\else
    \node[anchor=north west,font=\scSmall,text=scInkM]
      at ($#2+(0.56,-0.245)$) {#5};
  \fi}

% --- shared pgfplots look (only if the figure loaded pgfplots) ------------
\ifdefined\pgfplotsset
  \pgfplotsset{compat=1.18}
  \pgfplotsset{
    synaxis/.style={
      scale only axis,
      tick align=outside, tick pos=left,
      axis line style={draw=scInkM!55,line width=0.4pt},
      % One grid, on the value axis only.  No minor grid: at these sizes it
      % turns the plot frame into a grey wash.
      ymajorgrids=true, xmajorgrids=false,
      major grid style={draw=scRuleL,line width=0.3pt},
      major tick length=1.6pt,
      label style={font=\scBody,color=scInk},
      tick label style={font=\scSmall,color=scInkM},
      legend style={font=\scSmall,draw=none,fill=none,inner sep=1pt,
                    row sep=0.4pt,legend columns=1},
      legend cell align=left,
      every axis plot/.append style={line width=0.7pt,mark size=1.2pt},
    },
  }
\fi
}
% so that the figure compiles both from Figure/ and from the paper root.
%
% The visual vocabulary follows ../../../../Style/style.tex: the
% colour-blind-safe Okabe--Ito accents, quiet neutral rules and wash, circular
% panel badges, and a compact technical type scale are shared by the companion
% papers.  Semantic aliases below keep the CRN meaning explicit.
%
% DESIGN CONTRACT -----------------------------------------------------------
% 1. Geometry.  Every figure is laid out on a canvas exactly \scW wide (the
%    sn-jnl text width, 372 pt) and is included unscaled, so one point of type
%    is the same physical size in all five figures.  Open each picture with
%    \scCanvas{<height in cm>} to pin the bounding box.
% 2. Colour is semantic, never decorative.  Four roles only:
%       scAccent (green)       structure, the normal/expected path
%       scMark   (orange)      catalyst, certificate, or highlight
%       scAlert  (vermillion)  failure, violation, or uncovered set
%       scRef    (blue)        secondary or reference series
%    A panel gets at most one accent.  Titles are always ink; panel identity
%    is also encoded by a letter inside a circular badge, not by colour alone.
% 3. Type scale.  Three sizes, no others, and nothing below 6 pt:
%       \scTitle (7 pt bold)  panel titles
%       \scBody  (6.5 pt)     labels, values, in-figure prose
%       \scSmall (6 pt)       units, footnotes, dense graph labels
% ---------------------------------------------------------------------------
% Load guard.  The paper inputs every figure source inline, so this file is
% reached once per figure; without the guard the second \newlength below
% would abort the run.
\ifdefined\SynKitStyleLoaded\endinput\fi
\newcommand{\SynKitStyleLoaded}{}

% Load in the document preamble. XeLaTeX/LuaLaTeX and an installed Arial
% family are required; fontspec reports an error instead of substituting.
\RequirePackage[no-math]{fontspec}
\newfontfamily\scArial{Arial}[NFSSFamily=scArial,Ligatures=TeX]
% Keep manuscript mathematics unchanged; select Arial letters and numerals
% only inside figures. TeX supplies specialist mathematical symbols.
\RequirePackage[subdued,italic]{mathastext}
\MTDeclareVersion{synarial}{TU}{scArial}{m}{n}
\SetMathAlphabet{\Mathsf}{synarial}{TU}{scArial}{m}{n}
\SetMathAlphabet{\Mathtt}{synarial}{TU}{scArial}{m}{n}

\usetikzlibrary{positioning,calc,fit,backgrounds,arrows.meta,shapes.geometric,
                shapes.misc,decorations.pathreplacing,patterns,math}

% --- canvas ---------------------------------------------------------------
% 372 pt text width minus the 1 pt standalone border on each side.
\newlength{\scW}\setlength{\scW}{13.0326cm}
\providecommand{\scWnum}{13.0326}
% \scCanvas{<height>} pins the picture to the full text width, y running down
% from 0, so every figure shares one coordinate convention.  It *sets* the
% bounding box rather than contributing to it, so the emitted page is exactly
% the text width and every figure is included at 1:1 -- otherwise a stray
% invisible extent (pgfplots reserves a little around an axis) rescales one
% figure relative to the others and the same \scSmall prints at two sizes.
% Anything drawn outside the canvas is therefore cropped: keep content inside,
% and check with
%     \providecommand{\scCanvas}[1]{\path[draw=red] (0,0) rectangle (\scWnum,-#1);}
\providecommand{\scCanvas}[1]{%
  \path[use as bounding box] (0,0) rectangle (\scWnum,-#1);}

% --- palette --------------------------------------------------------------
% Canonical companion-paper values (../../../../Style/style.tex).
\definecolor{npgOrange}{HTML}{E69F00}
\definecolor{npgSky}{HTML}{56B4E9}
\definecolor{npgGreen}{HTML}{009E73}
\definecolor{npgYellow}{HTML}{F0E442}
\definecolor{npgBlue}{HTML}{0072B2}
\definecolor{npgVermillion}{HTML}{E64B35}
\definecolor{npgPurple}{HTML}{CC79A7}
\definecolor{figInk}{HTML}{1A1A1A}
\definecolor{figMuted}{HTML}{666666}
\definecolor{figHair}{HTML}{D9D9D9}
\definecolor{figWash}{HTML}{F7F7F7}

\colorlet{scInk}{figInk}              % primary text
\colorlet{scInkM}{figMuted}           % secondary text
\colorlet{scPaper}{figWash}           % quiet cards and note strips
\colorlet{scSpF}{npgSky!12}           % species fill
\colorlet{scSpE}{npgBlue}             % species edge
\colorlet{scRxF}{npgOrange!12}        % reaction fill
\colorlet{scRxE}{npgOrange!88!black}  % reaction edge
\colorlet{scReac}{npgBlue!78!black}   % reactant edge
\colorlet{scProd}{npgGreen!82!black}  % product edge
\colorlet{scHiF}{npgOrange!13}        % highlight fill
\colorlet{scHiE}{npgOrange!82!black}  % highlight edge
% Stage / accent tints.  stDE remains available for the fifth plot series.
\colorlet{stAF}{npgGreen!9}\colorlet{stAE}{npgGreen}
\colorlet{stBF}{npgOrange!11}\colorlet{stBE}{npgOrange}
\colorlet{stCF}{npgBlue!9}\colorlet{stCE}{npgBlue}
\colorlet{stDF}{npgPurple!10}\colorlet{stDE}{npgPurple!85!black}
\colorlet{stEF}{npgVermillion!9}\colorlet{stEE}{npgVermillion}

% Semantic aliases.  Figures should speak in these, not in stAE/stBE/...
\colorlet{scAccent}{stAE}\colorlet{scAccentF}{stAF}   % structure / normal
\colorlet{scMark}{stBE}\colorlet{scMarkF}{stBF}       % highlight / certificate
\colorlet{scRef}{stCE}\colorlet{scRefF}{stCF}         % secondary series
\colorlet{scAlert}{stEE}\colorlet{scAlertF}{stEF}     % failure / violation
\colorlet{scRule}{figHair}                            % card and axis hairlines
\colorlet{scRuleL}{figHair!55}                        % grid, dividers

% --- text styles ----------------------------------------------------------
% Three sizes, fixed in points so they do not drift with the document class.
% Arial is shared by labels, headings, panel badges, and chemistry labels.
\providecommand{\scTitle}{\scArial\fontsize{7}{8.4}\selectfont\bfseries\color{scInk}}
\providecommand{\scBody}{\scArial\fontsize{6.5}{7.8}\selectfont\color{scInkM}}
\providecommand{\scSmall}{\scArial\fontsize{6}{7.2}\selectfont\color{scInkM}}
\providecommand{\scSub}[1]{{\scSmall\itshape #1}}
\providecommand{\scPanel}[1]{{\scTitle(#1)}}
% Numerals that carry a result: ink, not colour, so the eye reads the value.
\providecommand{\scNum}{\scArial\fontsize{10}{11}\selectfont\bfseries\color{scInk}}

% --- shared TikZ styles ---------------------------------------------------
\tikzset{
  % text never hyphenates inside a figure
  every node/.append style={execute at begin node={%
      \hyphenpenalty=10000 \exhyphenpenalty=10000 \hbadness=10000 }},
  syn/.style={font=\scArial,color=scInk,
              execute at begin picture={\MTversion*{synarial}},
              >={Stealth[round,length=2.4mm,width=1.9mm]}},
  % containers ------------------------------------------------------------
  % One card chrome for the whole paper: white, 3 pt radius, one hairline.
  panel/.style={rounded corners=2.5pt,fill=none,draw=scRule,line width=0.45pt},
  card/.style={rounded corners=2.2pt,fill=scPaper,draw=scRule,line width=0.45pt,
               inner sep=5pt,align=left,anchor=north west},
  % a quiet strip for a panel's closing fact
  note/.style={rounded corners=2.5pt,fill=scPaper,draw=scRuleL,line width=0.4pt,
               inner xsep=6pt,inner ysep=3.5pt,align=left,font=\scSmall},
  % small labelled value pill
  chip/.style={rounded corners=2pt,fill=scPaper,draw=scRule,line width=0.4pt,
               inner xsep=4pt,inner ysep=2pt,font=\scSmall},
  % Circular letter badge: the same accessible panel marker used in Style/*.
  hbullet/.style={circle,fill=scInk,text=white,draw=none,
                  font=\scArial\fontsize{6}{6}\selectfont\bfseries,
                  minimum size=4.2mm,inner sep=0pt},
  ctag/.style={font=\scTitle,anchor=north west,text=scInk},
  % CRN primitives --------------------------------------------------------
  species/.style={rounded rectangle,fill=scSpF,draw=scSpE,line width=0.45pt,
                  inner xsep=2.4pt,inner ysep=1.6pt,font=\scSmall,
                  text=scInk,minimum height=3.6mm},
  cofactor/.style={species,fill=scRefF,draw=scRef},
  catalyst/.style={species,fill=scHiF,draw=scHiE},
  rxn/.style={rectangle,fill=scRxF,draw=scRxE,line width=0.45pt,
              inner sep=1.2pt,font=\scSmall,text=scInk,
              minimum width=3.0mm,minimum height=3.0mm},
  place/.style={circle,fill=scSpF,draw=scSpE,line width=0.45pt,
                minimum size=4.0mm,inner sep=0pt},
  trans/.style={rectangle,fill=scRxE,draw=scRxE,minimum width=1.1mm,
                minimum height=3.6mm,inner sep=0pt},
  token/.style={circle,fill=scInk,minimum size=1.1mm,inner sep=0pt},
  % edges -----------------------------------------------------------------
  reac/.style={->,line width=0.5pt,draw=scReac,shorten >=0.8pt,shorten <=0.8pt},
  prod/.style={->,line width=0.5pt,draw=scProd,shorten >=0.8pt,shorten <=0.8pt},
  cat/.style={->,line width=0.5pt,draw=scHiE,shorten >=0.8pt,shorten <=0.8pt},
  flow/.style={-{Stealth[round,length=2.2mm,width=1.8mm]},line width=1.1pt,
               draw=scInkM!45},
  flab/.style={font=\scSmall,text=scInkM,inner sep=2pt},
  stoi/.style={font=\scSmall,text=scInkM,inner sep=0.8pt},
  % verdict marks: shape carries the meaning, colour only reinforces it
  vpass/.style={circle,fill=scAccentF,draw=scAccent,text=scAccent,
                font=\fontsize{6}{6}\selectfont\bfseries,
                minimum size=3.2mm,inner sep=0pt},
  vfail/.style={circle,fill=scAlertF,draw=scAlert,text=scAlert,
                font=\fontsize{6}{6}\selectfont\bfseries,
                minimum size=3.2mm,inner sep=0pt},
  callout/.style={font=\scSmall,text=scInkM,align=left},
}

% \scHead[<badge colour>]{<anchor coord>}{<letter>}{<title>}{<subtitle or empty>}
% One header treatment for every panel in the paper.
\providecommand{\scHead}[5][scInk]{%
  \node[hbullet,fill=#1,anchor=north west] (scbadge#3) at ($#2+(0,-0.015)$)
    {\MakeUppercase{#3}};
  \node[ctag,anchor=west] at ($(scbadge#3.east)+(0.14,0)$) {#4};
  \ifx\relax#5\relax\else
    \node[anchor=north west,font=\scSmall,text=scInkM]
      at ($#2+(0.56,-0.245)$) {#5};
  \fi}

% --- shared pgfplots look (only if the figure loaded pgfplots) ------------
\ifdefined\pgfplotsset
  \pgfplotsset{compat=1.18}
  \pgfplotsset{
    synaxis/.style={
      scale only axis,
      tick align=outside, tick pos=left,
      axis line style={draw=scInkM!55,line width=0.4pt},
      % One grid, on the value axis only.  No minor grid: at these sizes it
      % turns the plot frame into a grey wash.
      ymajorgrids=true, xmajorgrids=false,
      major grid style={draw=scRuleL,line width=0.3pt},
      major tick length=1.6pt,
      label style={font=\scBody,color=scInk},
      tick label style={font=\scSmall,color=scInkM},
      legend style={font=\scSmall,draw=none,fill=none,inner sep=1pt,
                    row sep=0.4pt,legend columns=1},
      legend cell align=left,
      every axis plot/.append style={line width=0.7pt,mark size=1.2pt},
    },
  }
\fi
}{%
  \IfFileExists{Figure/synstyle.tex}{% ---------------------------------------------------------------------------
% synstyle.tex --- shared look for every TikZ figure of the SynKit CRN paper.
%
% Input it from a figure file:
%     \IfFileExists{synstyle.tex}{% ---------------------------------------------------------------------------
% synstyle.tex --- shared look for every TikZ figure of the SynKit CRN paper.
%
% Input it from a figure file:
%     \IfFileExists{synstyle.tex}{% ---------------------------------------------------------------------------
% synstyle.tex --- shared look for every TikZ figure of the SynKit CRN paper.
%
% Input it from a figure file:
%     \IfFileExists{synstyle.tex}{\input{synstyle}}{\input{Figure/synstyle}}
% so that the figure compiles both from Figure/ and from the paper root.
%
% The visual vocabulary follows ../../../../Style/style.tex: the
% colour-blind-safe Okabe--Ito accents, quiet neutral rules and wash, circular
% panel badges, and a compact technical type scale are shared by the companion
% papers.  Semantic aliases below keep the CRN meaning explicit.
%
% DESIGN CONTRACT -----------------------------------------------------------
% 1. Geometry.  Every figure is laid out on a canvas exactly \scW wide (the
%    sn-jnl text width, 372 pt) and is included unscaled, so one point of type
%    is the same physical size in all five figures.  Open each picture with
%    \scCanvas{<height in cm>} to pin the bounding box.
% 2. Colour is semantic, never decorative.  Four roles only:
%       scAccent (green)       structure, the normal/expected path
%       scMark   (orange)      catalyst, certificate, or highlight
%       scAlert  (vermillion)  failure, violation, or uncovered set
%       scRef    (blue)        secondary or reference series
%    A panel gets at most one accent.  Titles are always ink; panel identity
%    is also encoded by a letter inside a circular badge, not by colour alone.
% 3. Type scale.  Three sizes, no others, and nothing below 6 pt:
%       \scTitle (7 pt bold)  panel titles
%       \scBody  (6.5 pt)     labels, values, in-figure prose
%       \scSmall (6 pt)       units, footnotes, dense graph labels
% ---------------------------------------------------------------------------
% Load guard.  The paper inputs every figure source inline, so this file is
% reached once per figure; without the guard the second \newlength below
% would abort the run.
\ifdefined\SynKitStyleLoaded\endinput\fi
\newcommand{\SynKitStyleLoaded}{}

% Load in the document preamble. XeLaTeX/LuaLaTeX and an installed Arial
% family are required; fontspec reports an error instead of substituting.
\RequirePackage[no-math]{fontspec}
\newfontfamily\scArial{Arial}[NFSSFamily=scArial,Ligatures=TeX]
% Keep manuscript mathematics unchanged; select Arial letters and numerals
% only inside figures. TeX supplies specialist mathematical symbols.
\RequirePackage[subdued,italic]{mathastext}
\MTDeclareVersion{synarial}{TU}{scArial}{m}{n}
\SetMathAlphabet{\Mathsf}{synarial}{TU}{scArial}{m}{n}
\SetMathAlphabet{\Mathtt}{synarial}{TU}{scArial}{m}{n}

\usetikzlibrary{positioning,calc,fit,backgrounds,arrows.meta,shapes.geometric,
                shapes.misc,decorations.pathreplacing,patterns,math}

% --- canvas ---------------------------------------------------------------
% 372 pt text width minus the 1 pt standalone border on each side.
\newlength{\scW}\setlength{\scW}{13.0326cm}
\providecommand{\scWnum}{13.0326}
% \scCanvas{<height>} pins the picture to the full text width, y running down
% from 0, so every figure shares one coordinate convention.  It *sets* the
% bounding box rather than contributing to it, so the emitted page is exactly
% the text width and every figure is included at 1:1 -- otherwise a stray
% invisible extent (pgfplots reserves a little around an axis) rescales one
% figure relative to the others and the same \scSmall prints at two sizes.
% Anything drawn outside the canvas is therefore cropped: keep content inside,
% and check with
%     \providecommand{\scCanvas}[1]{\path[draw=red] (0,0) rectangle (\scWnum,-#1);}
\providecommand{\scCanvas}[1]{%
  \path[use as bounding box] (0,0) rectangle (\scWnum,-#1);}

% --- palette --------------------------------------------------------------
% Canonical companion-paper values (../../../../Style/style.tex).
\definecolor{npgOrange}{HTML}{E69F00}
\definecolor{npgSky}{HTML}{56B4E9}
\definecolor{npgGreen}{HTML}{009E73}
\definecolor{npgYellow}{HTML}{F0E442}
\definecolor{npgBlue}{HTML}{0072B2}
\definecolor{npgVermillion}{HTML}{E64B35}
\definecolor{npgPurple}{HTML}{CC79A7}
\definecolor{figInk}{HTML}{1A1A1A}
\definecolor{figMuted}{HTML}{666666}
\definecolor{figHair}{HTML}{D9D9D9}
\definecolor{figWash}{HTML}{F7F7F7}

\colorlet{scInk}{figInk}              % primary text
\colorlet{scInkM}{figMuted}           % secondary text
\colorlet{scPaper}{figWash}           % quiet cards and note strips
\colorlet{scSpF}{npgSky!12}           % species fill
\colorlet{scSpE}{npgBlue}             % species edge
\colorlet{scRxF}{npgOrange!12}        % reaction fill
\colorlet{scRxE}{npgOrange!88!black}  % reaction edge
\colorlet{scReac}{npgBlue!78!black}   % reactant edge
\colorlet{scProd}{npgGreen!82!black}  % product edge
\colorlet{scHiF}{npgOrange!13}        % highlight fill
\colorlet{scHiE}{npgOrange!82!black}  % highlight edge
% Stage / accent tints.  stDE remains available for the fifth plot series.
\colorlet{stAF}{npgGreen!9}\colorlet{stAE}{npgGreen}
\colorlet{stBF}{npgOrange!11}\colorlet{stBE}{npgOrange}
\colorlet{stCF}{npgBlue!9}\colorlet{stCE}{npgBlue}
\colorlet{stDF}{npgPurple!10}\colorlet{stDE}{npgPurple!85!black}
\colorlet{stEF}{npgVermillion!9}\colorlet{stEE}{npgVermillion}

% Semantic aliases.  Figures should speak in these, not in stAE/stBE/...
\colorlet{scAccent}{stAE}\colorlet{scAccentF}{stAF}   % structure / normal
\colorlet{scMark}{stBE}\colorlet{scMarkF}{stBF}       % highlight / certificate
\colorlet{scRef}{stCE}\colorlet{scRefF}{stCF}         % secondary series
\colorlet{scAlert}{stEE}\colorlet{scAlertF}{stEF}     % failure / violation
\colorlet{scRule}{figHair}                            % card and axis hairlines
\colorlet{scRuleL}{figHair!55}                        % grid, dividers

% --- text styles ----------------------------------------------------------
% Three sizes, fixed in points so they do not drift with the document class.
% Arial is shared by labels, headings, panel badges, and chemistry labels.
\providecommand{\scTitle}{\scArial\fontsize{7}{8.4}\selectfont\bfseries\color{scInk}}
\providecommand{\scBody}{\scArial\fontsize{6.5}{7.8}\selectfont\color{scInkM}}
\providecommand{\scSmall}{\scArial\fontsize{6}{7.2}\selectfont\color{scInkM}}
\providecommand{\scSub}[1]{{\scSmall\itshape #1}}
\providecommand{\scPanel}[1]{{\scTitle(#1)}}
% Numerals that carry a result: ink, not colour, so the eye reads the value.
\providecommand{\scNum}{\scArial\fontsize{10}{11}\selectfont\bfseries\color{scInk}}

% --- shared TikZ styles ---------------------------------------------------
\tikzset{
  % text never hyphenates inside a figure
  every node/.append style={execute at begin node={%
      \hyphenpenalty=10000 \exhyphenpenalty=10000 \hbadness=10000 }},
  syn/.style={font=\scArial,color=scInk,
              execute at begin picture={\MTversion*{synarial}},
              >={Stealth[round,length=2.4mm,width=1.9mm]}},
  % containers ------------------------------------------------------------
  % One card chrome for the whole paper: white, 3 pt radius, one hairline.
  panel/.style={rounded corners=2.5pt,fill=none,draw=scRule,line width=0.45pt},
  card/.style={rounded corners=2.2pt,fill=scPaper,draw=scRule,line width=0.45pt,
               inner sep=5pt,align=left,anchor=north west},
  % a quiet strip for a panel's closing fact
  note/.style={rounded corners=2.5pt,fill=scPaper,draw=scRuleL,line width=0.4pt,
               inner xsep=6pt,inner ysep=3.5pt,align=left,font=\scSmall},
  % small labelled value pill
  chip/.style={rounded corners=2pt,fill=scPaper,draw=scRule,line width=0.4pt,
               inner xsep=4pt,inner ysep=2pt,font=\scSmall},
  % Circular letter badge: the same accessible panel marker used in Style/*.
  hbullet/.style={circle,fill=scInk,text=white,draw=none,
                  font=\scArial\fontsize{6}{6}\selectfont\bfseries,
                  minimum size=4.2mm,inner sep=0pt},
  ctag/.style={font=\scTitle,anchor=north west,text=scInk},
  % CRN primitives --------------------------------------------------------
  species/.style={rounded rectangle,fill=scSpF,draw=scSpE,line width=0.45pt,
                  inner xsep=2.4pt,inner ysep=1.6pt,font=\scSmall,
                  text=scInk,minimum height=3.6mm},
  cofactor/.style={species,fill=scRefF,draw=scRef},
  catalyst/.style={species,fill=scHiF,draw=scHiE},
  rxn/.style={rectangle,fill=scRxF,draw=scRxE,line width=0.45pt,
              inner sep=1.2pt,font=\scSmall,text=scInk,
              minimum width=3.0mm,minimum height=3.0mm},
  place/.style={circle,fill=scSpF,draw=scSpE,line width=0.45pt,
                minimum size=4.0mm,inner sep=0pt},
  trans/.style={rectangle,fill=scRxE,draw=scRxE,minimum width=1.1mm,
                minimum height=3.6mm,inner sep=0pt},
  token/.style={circle,fill=scInk,minimum size=1.1mm,inner sep=0pt},
  % edges -----------------------------------------------------------------
  reac/.style={->,line width=0.5pt,draw=scReac,shorten >=0.8pt,shorten <=0.8pt},
  prod/.style={->,line width=0.5pt,draw=scProd,shorten >=0.8pt,shorten <=0.8pt},
  cat/.style={->,line width=0.5pt,draw=scHiE,shorten >=0.8pt,shorten <=0.8pt},
  flow/.style={-{Stealth[round,length=2.2mm,width=1.8mm]},line width=1.1pt,
               draw=scInkM!45},
  flab/.style={font=\scSmall,text=scInkM,inner sep=2pt},
  stoi/.style={font=\scSmall,text=scInkM,inner sep=0.8pt},
  % verdict marks: shape carries the meaning, colour only reinforces it
  vpass/.style={circle,fill=scAccentF,draw=scAccent,text=scAccent,
                font=\fontsize{6}{6}\selectfont\bfseries,
                minimum size=3.2mm,inner sep=0pt},
  vfail/.style={circle,fill=scAlertF,draw=scAlert,text=scAlert,
                font=\fontsize{6}{6}\selectfont\bfseries,
                minimum size=3.2mm,inner sep=0pt},
  callout/.style={font=\scSmall,text=scInkM,align=left},
}

% \scHead[<badge colour>]{<anchor coord>}{<letter>}{<title>}{<subtitle or empty>}
% One header treatment for every panel in the paper.
\providecommand{\scHead}[5][scInk]{%
  \node[hbullet,fill=#1,anchor=north west] (scbadge#3) at ($#2+(0,-0.015)$)
    {\MakeUppercase{#3}};
  \node[ctag,anchor=west] at ($(scbadge#3.east)+(0.14,0)$) {#4};
  \ifx\relax#5\relax\else
    \node[anchor=north west,font=\scSmall,text=scInkM]
      at ($#2+(0.56,-0.245)$) {#5};
  \fi}

% --- shared pgfplots look (only if the figure loaded pgfplots) ------------
\ifdefined\pgfplotsset
  \pgfplotsset{compat=1.18}
  \pgfplotsset{
    synaxis/.style={
      scale only axis,
      tick align=outside, tick pos=left,
      axis line style={draw=scInkM!55,line width=0.4pt},
      % One grid, on the value axis only.  No minor grid: at these sizes it
      % turns the plot frame into a grey wash.
      ymajorgrids=true, xmajorgrids=false,
      major grid style={draw=scRuleL,line width=0.3pt},
      major tick length=1.6pt,
      label style={font=\scBody,color=scInk},
      tick label style={font=\scSmall,color=scInkM},
      legend style={font=\scSmall,draw=none,fill=none,inner sep=1pt,
                    row sep=0.4pt,legend columns=1},
      legend cell align=left,
      every axis plot/.append style={line width=0.7pt,mark size=1.2pt},
    },
  }
\fi
}{% ---------------------------------------------------------------------------
% synstyle.tex --- shared look for every TikZ figure of the SynKit CRN paper.
%
% Input it from a figure file:
%     \IfFileExists{synstyle.tex}{\input{synstyle}}{\input{Figure/synstyle}}
% so that the figure compiles both from Figure/ and from the paper root.
%
% The visual vocabulary follows ../../../../Style/style.tex: the
% colour-blind-safe Okabe--Ito accents, quiet neutral rules and wash, circular
% panel badges, and a compact technical type scale are shared by the companion
% papers.  Semantic aliases below keep the CRN meaning explicit.
%
% DESIGN CONTRACT -----------------------------------------------------------
% 1. Geometry.  Every figure is laid out on a canvas exactly \scW wide (the
%    sn-jnl text width, 372 pt) and is included unscaled, so one point of type
%    is the same physical size in all five figures.  Open each picture with
%    \scCanvas{<height in cm>} to pin the bounding box.
% 2. Colour is semantic, never decorative.  Four roles only:
%       scAccent (green)       structure, the normal/expected path
%       scMark   (orange)      catalyst, certificate, or highlight
%       scAlert  (vermillion)  failure, violation, or uncovered set
%       scRef    (blue)        secondary or reference series
%    A panel gets at most one accent.  Titles are always ink; panel identity
%    is also encoded by a letter inside a circular badge, not by colour alone.
% 3. Type scale.  Three sizes, no others, and nothing below 6 pt:
%       \scTitle (7 pt bold)  panel titles
%       \scBody  (6.5 pt)     labels, values, in-figure prose
%       \scSmall (6 pt)       units, footnotes, dense graph labels
% ---------------------------------------------------------------------------
% Load guard.  The paper inputs every figure source inline, so this file is
% reached once per figure; without the guard the second \newlength below
% would abort the run.
\ifdefined\SynKitStyleLoaded\endinput\fi
\newcommand{\SynKitStyleLoaded}{}

% Load in the document preamble. XeLaTeX/LuaLaTeX and an installed Arial
% family are required; fontspec reports an error instead of substituting.
\RequirePackage[no-math]{fontspec}
\newfontfamily\scArial{Arial}[NFSSFamily=scArial,Ligatures=TeX]
% Keep manuscript mathematics unchanged; select Arial letters and numerals
% only inside figures. TeX supplies specialist mathematical symbols.
\RequirePackage[subdued,italic]{mathastext}
\MTDeclareVersion{synarial}{TU}{scArial}{m}{n}
\SetMathAlphabet{\Mathsf}{synarial}{TU}{scArial}{m}{n}
\SetMathAlphabet{\Mathtt}{synarial}{TU}{scArial}{m}{n}

\usetikzlibrary{positioning,calc,fit,backgrounds,arrows.meta,shapes.geometric,
                shapes.misc,decorations.pathreplacing,patterns,math}

% --- canvas ---------------------------------------------------------------
% 372 pt text width minus the 1 pt standalone border on each side.
\newlength{\scW}\setlength{\scW}{13.0326cm}
\providecommand{\scWnum}{13.0326}
% \scCanvas{<height>} pins the picture to the full text width, y running down
% from 0, so every figure shares one coordinate convention.  It *sets* the
% bounding box rather than contributing to it, so the emitted page is exactly
% the text width and every figure is included at 1:1 -- otherwise a stray
% invisible extent (pgfplots reserves a little around an axis) rescales one
% figure relative to the others and the same \scSmall prints at two sizes.
% Anything drawn outside the canvas is therefore cropped: keep content inside,
% and check with
%     \providecommand{\scCanvas}[1]{\path[draw=red] (0,0) rectangle (\scWnum,-#1);}
\providecommand{\scCanvas}[1]{%
  \path[use as bounding box] (0,0) rectangle (\scWnum,-#1);}

% --- palette --------------------------------------------------------------
% Canonical companion-paper values (../../../../Style/style.tex).
\definecolor{npgOrange}{HTML}{E69F00}
\definecolor{npgSky}{HTML}{56B4E9}
\definecolor{npgGreen}{HTML}{009E73}
\definecolor{npgYellow}{HTML}{F0E442}
\definecolor{npgBlue}{HTML}{0072B2}
\definecolor{npgVermillion}{HTML}{E64B35}
\definecolor{npgPurple}{HTML}{CC79A7}
\definecolor{figInk}{HTML}{1A1A1A}
\definecolor{figMuted}{HTML}{666666}
\definecolor{figHair}{HTML}{D9D9D9}
\definecolor{figWash}{HTML}{F7F7F7}

\colorlet{scInk}{figInk}              % primary text
\colorlet{scInkM}{figMuted}           % secondary text
\colorlet{scPaper}{figWash}           % quiet cards and note strips
\colorlet{scSpF}{npgSky!12}           % species fill
\colorlet{scSpE}{npgBlue}             % species edge
\colorlet{scRxF}{npgOrange!12}        % reaction fill
\colorlet{scRxE}{npgOrange!88!black}  % reaction edge
\colorlet{scReac}{npgBlue!78!black}   % reactant edge
\colorlet{scProd}{npgGreen!82!black}  % product edge
\colorlet{scHiF}{npgOrange!13}        % highlight fill
\colorlet{scHiE}{npgOrange!82!black}  % highlight edge
% Stage / accent tints.  stDE remains available for the fifth plot series.
\colorlet{stAF}{npgGreen!9}\colorlet{stAE}{npgGreen}
\colorlet{stBF}{npgOrange!11}\colorlet{stBE}{npgOrange}
\colorlet{stCF}{npgBlue!9}\colorlet{stCE}{npgBlue}
\colorlet{stDF}{npgPurple!10}\colorlet{stDE}{npgPurple!85!black}
\colorlet{stEF}{npgVermillion!9}\colorlet{stEE}{npgVermillion}

% Semantic aliases.  Figures should speak in these, not in stAE/stBE/...
\colorlet{scAccent}{stAE}\colorlet{scAccentF}{stAF}   % structure / normal
\colorlet{scMark}{stBE}\colorlet{scMarkF}{stBF}       % highlight / certificate
\colorlet{scRef}{stCE}\colorlet{scRefF}{stCF}         % secondary series
\colorlet{scAlert}{stEE}\colorlet{scAlertF}{stEF}     % failure / violation
\colorlet{scRule}{figHair}                            % card and axis hairlines
\colorlet{scRuleL}{figHair!55}                        % grid, dividers

% --- text styles ----------------------------------------------------------
% Three sizes, fixed in points so they do not drift with the document class.
% Arial is shared by labels, headings, panel badges, and chemistry labels.
\providecommand{\scTitle}{\scArial\fontsize{7}{8.4}\selectfont\bfseries\color{scInk}}
\providecommand{\scBody}{\scArial\fontsize{6.5}{7.8}\selectfont\color{scInkM}}
\providecommand{\scSmall}{\scArial\fontsize{6}{7.2}\selectfont\color{scInkM}}
\providecommand{\scSub}[1]{{\scSmall\itshape #1}}
\providecommand{\scPanel}[1]{{\scTitle(#1)}}
% Numerals that carry a result: ink, not colour, so the eye reads the value.
\providecommand{\scNum}{\scArial\fontsize{10}{11}\selectfont\bfseries\color{scInk}}

% --- shared TikZ styles ---------------------------------------------------
\tikzset{
  % text never hyphenates inside a figure
  every node/.append style={execute at begin node={%
      \hyphenpenalty=10000 \exhyphenpenalty=10000 \hbadness=10000 }},
  syn/.style={font=\scArial,color=scInk,
              execute at begin picture={\MTversion*{synarial}},
              >={Stealth[round,length=2.4mm,width=1.9mm]}},
  % containers ------------------------------------------------------------
  % One card chrome for the whole paper: white, 3 pt radius, one hairline.
  panel/.style={rounded corners=2.5pt,fill=none,draw=scRule,line width=0.45pt},
  card/.style={rounded corners=2.2pt,fill=scPaper,draw=scRule,line width=0.45pt,
               inner sep=5pt,align=left,anchor=north west},
  % a quiet strip for a panel's closing fact
  note/.style={rounded corners=2.5pt,fill=scPaper,draw=scRuleL,line width=0.4pt,
               inner xsep=6pt,inner ysep=3.5pt,align=left,font=\scSmall},
  % small labelled value pill
  chip/.style={rounded corners=2pt,fill=scPaper,draw=scRule,line width=0.4pt,
               inner xsep=4pt,inner ysep=2pt,font=\scSmall},
  % Circular letter badge: the same accessible panel marker used in Style/*.
  hbullet/.style={circle,fill=scInk,text=white,draw=none,
                  font=\scArial\fontsize{6}{6}\selectfont\bfseries,
                  minimum size=4.2mm,inner sep=0pt},
  ctag/.style={font=\scTitle,anchor=north west,text=scInk},
  % CRN primitives --------------------------------------------------------
  species/.style={rounded rectangle,fill=scSpF,draw=scSpE,line width=0.45pt,
                  inner xsep=2.4pt,inner ysep=1.6pt,font=\scSmall,
                  text=scInk,minimum height=3.6mm},
  cofactor/.style={species,fill=scRefF,draw=scRef},
  catalyst/.style={species,fill=scHiF,draw=scHiE},
  rxn/.style={rectangle,fill=scRxF,draw=scRxE,line width=0.45pt,
              inner sep=1.2pt,font=\scSmall,text=scInk,
              minimum width=3.0mm,minimum height=3.0mm},
  place/.style={circle,fill=scSpF,draw=scSpE,line width=0.45pt,
                minimum size=4.0mm,inner sep=0pt},
  trans/.style={rectangle,fill=scRxE,draw=scRxE,minimum width=1.1mm,
                minimum height=3.6mm,inner sep=0pt},
  token/.style={circle,fill=scInk,minimum size=1.1mm,inner sep=0pt},
  % edges -----------------------------------------------------------------
  reac/.style={->,line width=0.5pt,draw=scReac,shorten >=0.8pt,shorten <=0.8pt},
  prod/.style={->,line width=0.5pt,draw=scProd,shorten >=0.8pt,shorten <=0.8pt},
  cat/.style={->,line width=0.5pt,draw=scHiE,shorten >=0.8pt,shorten <=0.8pt},
  flow/.style={-{Stealth[round,length=2.2mm,width=1.8mm]},line width=1.1pt,
               draw=scInkM!45},
  flab/.style={font=\scSmall,text=scInkM,inner sep=2pt},
  stoi/.style={font=\scSmall,text=scInkM,inner sep=0.8pt},
  % verdict marks: shape carries the meaning, colour only reinforces it
  vpass/.style={circle,fill=scAccentF,draw=scAccent,text=scAccent,
                font=\fontsize{6}{6}\selectfont\bfseries,
                minimum size=3.2mm,inner sep=0pt},
  vfail/.style={circle,fill=scAlertF,draw=scAlert,text=scAlert,
                font=\fontsize{6}{6}\selectfont\bfseries,
                minimum size=3.2mm,inner sep=0pt},
  callout/.style={font=\scSmall,text=scInkM,align=left},
}

% \scHead[<badge colour>]{<anchor coord>}{<letter>}{<title>}{<subtitle or empty>}
% One header treatment for every panel in the paper.
\providecommand{\scHead}[5][scInk]{%
  \node[hbullet,fill=#1,anchor=north west] (scbadge#3) at ($#2+(0,-0.015)$)
    {\MakeUppercase{#3}};
  \node[ctag,anchor=west] at ($(scbadge#3.east)+(0.14,0)$) {#4};
  \ifx\relax#5\relax\else
    \node[anchor=north west,font=\scSmall,text=scInkM]
      at ($#2+(0.56,-0.245)$) {#5};
  \fi}

% --- shared pgfplots look (only if the figure loaded pgfplots) ------------
\ifdefined\pgfplotsset
  \pgfplotsset{compat=1.18}
  \pgfplotsset{
    synaxis/.style={
      scale only axis,
      tick align=outside, tick pos=left,
      axis line style={draw=scInkM!55,line width=0.4pt},
      % One grid, on the value axis only.  No minor grid: at these sizes it
      % turns the plot frame into a grey wash.
      ymajorgrids=true, xmajorgrids=false,
      major grid style={draw=scRuleL,line width=0.3pt},
      major tick length=1.6pt,
      label style={font=\scBody,color=scInk},
      tick label style={font=\scSmall,color=scInkM},
      legend style={font=\scSmall,draw=none,fill=none,inner sep=1pt,
                    row sep=0.4pt,legend columns=1},
      legend cell align=left,
      every axis plot/.append style={line width=0.7pt,mark size=1.2pt},
    },
  }
\fi
}
% so that the figure compiles both from Figure/ and from the paper root.
%
% The visual vocabulary follows ../../../../Style/style.tex: the
% colour-blind-safe Okabe--Ito accents, quiet neutral rules and wash, circular
% panel badges, and a compact technical type scale are shared by the companion
% papers.  Semantic aliases below keep the CRN meaning explicit.
%
% DESIGN CONTRACT -----------------------------------------------------------
% 1. Geometry.  Every figure is laid out on a canvas exactly \scW wide (the
%    sn-jnl text width, 372 pt) and is included unscaled, so one point of type
%    is the same physical size in all five figures.  Open each picture with
%    \scCanvas{<height in cm>} to pin the bounding box.
% 2. Colour is semantic, never decorative.  Four roles only:
%       scAccent (green)       structure, the normal/expected path
%       scMark   (orange)      catalyst, certificate, or highlight
%       scAlert  (vermillion)  failure, violation, or uncovered set
%       scRef    (blue)        secondary or reference series
%    A panel gets at most one accent.  Titles are always ink; panel identity
%    is also encoded by a letter inside a circular badge, not by colour alone.
% 3. Type scale.  Three sizes, no others, and nothing below 6 pt:
%       \scTitle (7 pt bold)  panel titles
%       \scBody  (6.5 pt)     labels, values, in-figure prose
%       \scSmall (6 pt)       units, footnotes, dense graph labels
% ---------------------------------------------------------------------------
% Load guard.  The paper inputs every figure source inline, so this file is
% reached once per figure; without the guard the second \newlength below
% would abort the run.
\ifdefined\SynKitStyleLoaded\endinput\fi
\newcommand{\SynKitStyleLoaded}{}

% Load in the document preamble. XeLaTeX/LuaLaTeX and an installed Arial
% family are required; fontspec reports an error instead of substituting.
\RequirePackage[no-math]{fontspec}
\newfontfamily\scArial{Arial}[NFSSFamily=scArial,Ligatures=TeX]
% Keep manuscript mathematics unchanged; select Arial letters and numerals
% only inside figures. TeX supplies specialist mathematical symbols.
\RequirePackage[subdued,italic]{mathastext}
\MTDeclareVersion{synarial}{TU}{scArial}{m}{n}
\SetMathAlphabet{\Mathsf}{synarial}{TU}{scArial}{m}{n}
\SetMathAlphabet{\Mathtt}{synarial}{TU}{scArial}{m}{n}

\usetikzlibrary{positioning,calc,fit,backgrounds,arrows.meta,shapes.geometric,
                shapes.misc,decorations.pathreplacing,patterns,math}

% --- canvas ---------------------------------------------------------------
% 372 pt text width minus the 1 pt standalone border on each side.
\newlength{\scW}\setlength{\scW}{13.0326cm}
\providecommand{\scWnum}{13.0326}
% \scCanvas{<height>} pins the picture to the full text width, y running down
% from 0, so every figure shares one coordinate convention.  It *sets* the
% bounding box rather than contributing to it, so the emitted page is exactly
% the text width and every figure is included at 1:1 -- otherwise a stray
% invisible extent (pgfplots reserves a little around an axis) rescales one
% figure relative to the others and the same \scSmall prints at two sizes.
% Anything drawn outside the canvas is therefore cropped: keep content inside,
% and check with
%     \providecommand{\scCanvas}[1]{\path[draw=red] (0,0) rectangle (\scWnum,-#1);}
\providecommand{\scCanvas}[1]{%
  \path[use as bounding box] (0,0) rectangle (\scWnum,-#1);}

% --- palette --------------------------------------------------------------
% Canonical companion-paper values (../../../../Style/style.tex).
\definecolor{npgOrange}{HTML}{E69F00}
\definecolor{npgSky}{HTML}{56B4E9}
\definecolor{npgGreen}{HTML}{009E73}
\definecolor{npgYellow}{HTML}{F0E442}
\definecolor{npgBlue}{HTML}{0072B2}
\definecolor{npgVermillion}{HTML}{E64B35}
\definecolor{npgPurple}{HTML}{CC79A7}
\definecolor{figInk}{HTML}{1A1A1A}
\definecolor{figMuted}{HTML}{666666}
\definecolor{figHair}{HTML}{D9D9D9}
\definecolor{figWash}{HTML}{F7F7F7}

\colorlet{scInk}{figInk}              % primary text
\colorlet{scInkM}{figMuted}           % secondary text
\colorlet{scPaper}{figWash}           % quiet cards and note strips
\colorlet{scSpF}{npgSky!12}           % species fill
\colorlet{scSpE}{npgBlue}             % species edge
\colorlet{scRxF}{npgOrange!12}        % reaction fill
\colorlet{scRxE}{npgOrange!88!black}  % reaction edge
\colorlet{scReac}{npgBlue!78!black}   % reactant edge
\colorlet{scProd}{npgGreen!82!black}  % product edge
\colorlet{scHiF}{npgOrange!13}        % highlight fill
\colorlet{scHiE}{npgOrange!82!black}  % highlight edge
% Stage / accent tints.  stDE remains available for the fifth plot series.
\colorlet{stAF}{npgGreen!9}\colorlet{stAE}{npgGreen}
\colorlet{stBF}{npgOrange!11}\colorlet{stBE}{npgOrange}
\colorlet{stCF}{npgBlue!9}\colorlet{stCE}{npgBlue}
\colorlet{stDF}{npgPurple!10}\colorlet{stDE}{npgPurple!85!black}
\colorlet{stEF}{npgVermillion!9}\colorlet{stEE}{npgVermillion}

% Semantic aliases.  Figures should speak in these, not in stAE/stBE/...
\colorlet{scAccent}{stAE}\colorlet{scAccentF}{stAF}   % structure / normal
\colorlet{scMark}{stBE}\colorlet{scMarkF}{stBF}       % highlight / certificate
\colorlet{scRef}{stCE}\colorlet{scRefF}{stCF}         % secondary series
\colorlet{scAlert}{stEE}\colorlet{scAlertF}{stEF}     % failure / violation
\colorlet{scRule}{figHair}                            % card and axis hairlines
\colorlet{scRuleL}{figHair!55}                        % grid, dividers

% --- text styles ----------------------------------------------------------
% Three sizes, fixed in points so they do not drift with the document class.
% Arial is shared by labels, headings, panel badges, and chemistry labels.
\providecommand{\scTitle}{\scArial\fontsize{7}{8.4}\selectfont\bfseries\color{scInk}}
\providecommand{\scBody}{\scArial\fontsize{6.5}{7.8}\selectfont\color{scInkM}}
\providecommand{\scSmall}{\scArial\fontsize{6}{7.2}\selectfont\color{scInkM}}
\providecommand{\scSub}[1]{{\scSmall\itshape #1}}
\providecommand{\scPanel}[1]{{\scTitle(#1)}}
% Numerals that carry a result: ink, not colour, so the eye reads the value.
\providecommand{\scNum}{\scArial\fontsize{10}{11}\selectfont\bfseries\color{scInk}}

% --- shared TikZ styles ---------------------------------------------------
\tikzset{
  % text never hyphenates inside a figure
  every node/.append style={execute at begin node={%
      \hyphenpenalty=10000 \exhyphenpenalty=10000 \hbadness=10000 }},
  syn/.style={font=\scArial,color=scInk,
              execute at begin picture={\MTversion*{synarial}},
              >={Stealth[round,length=2.4mm,width=1.9mm]}},
  % containers ------------------------------------------------------------
  % One card chrome for the whole paper: white, 3 pt radius, one hairline.
  panel/.style={rounded corners=2.5pt,fill=none,draw=scRule,line width=0.45pt},
  card/.style={rounded corners=2.2pt,fill=scPaper,draw=scRule,line width=0.45pt,
               inner sep=5pt,align=left,anchor=north west},
  % a quiet strip for a panel's closing fact
  note/.style={rounded corners=2.5pt,fill=scPaper,draw=scRuleL,line width=0.4pt,
               inner xsep=6pt,inner ysep=3.5pt,align=left,font=\scSmall},
  % small labelled value pill
  chip/.style={rounded corners=2pt,fill=scPaper,draw=scRule,line width=0.4pt,
               inner xsep=4pt,inner ysep=2pt,font=\scSmall},
  % Circular letter badge: the same accessible panel marker used in Style/*.
  hbullet/.style={circle,fill=scInk,text=white,draw=none,
                  font=\scArial\fontsize{6}{6}\selectfont\bfseries,
                  minimum size=4.2mm,inner sep=0pt},
  ctag/.style={font=\scTitle,anchor=north west,text=scInk},
  % CRN primitives --------------------------------------------------------
  species/.style={rounded rectangle,fill=scSpF,draw=scSpE,line width=0.45pt,
                  inner xsep=2.4pt,inner ysep=1.6pt,font=\scSmall,
                  text=scInk,minimum height=3.6mm},
  cofactor/.style={species,fill=scRefF,draw=scRef},
  catalyst/.style={species,fill=scHiF,draw=scHiE},
  rxn/.style={rectangle,fill=scRxF,draw=scRxE,line width=0.45pt,
              inner sep=1.2pt,font=\scSmall,text=scInk,
              minimum width=3.0mm,minimum height=3.0mm},
  place/.style={circle,fill=scSpF,draw=scSpE,line width=0.45pt,
                minimum size=4.0mm,inner sep=0pt},
  trans/.style={rectangle,fill=scRxE,draw=scRxE,minimum width=1.1mm,
                minimum height=3.6mm,inner sep=0pt},
  token/.style={circle,fill=scInk,minimum size=1.1mm,inner sep=0pt},
  % edges -----------------------------------------------------------------
  reac/.style={->,line width=0.5pt,draw=scReac,shorten >=0.8pt,shorten <=0.8pt},
  prod/.style={->,line width=0.5pt,draw=scProd,shorten >=0.8pt,shorten <=0.8pt},
  cat/.style={->,line width=0.5pt,draw=scHiE,shorten >=0.8pt,shorten <=0.8pt},
  flow/.style={-{Stealth[round,length=2.2mm,width=1.8mm]},line width=1.1pt,
               draw=scInkM!45},
  flab/.style={font=\scSmall,text=scInkM,inner sep=2pt},
  stoi/.style={font=\scSmall,text=scInkM,inner sep=0.8pt},
  % verdict marks: shape carries the meaning, colour only reinforces it
  vpass/.style={circle,fill=scAccentF,draw=scAccent,text=scAccent,
                font=\fontsize{6}{6}\selectfont\bfseries,
                minimum size=3.2mm,inner sep=0pt},
  vfail/.style={circle,fill=scAlertF,draw=scAlert,text=scAlert,
                font=\fontsize{6}{6}\selectfont\bfseries,
                minimum size=3.2mm,inner sep=0pt},
  callout/.style={font=\scSmall,text=scInkM,align=left},
}

% \scHead[<badge colour>]{<anchor coord>}{<letter>}{<title>}{<subtitle or empty>}
% One header treatment for every panel in the paper.
\providecommand{\scHead}[5][scInk]{%
  \node[hbullet,fill=#1,anchor=north west] (scbadge#3) at ($#2+(0,-0.015)$)
    {\MakeUppercase{#3}};
  \node[ctag,anchor=west] at ($(scbadge#3.east)+(0.14,0)$) {#4};
  \ifx\relax#5\relax\else
    \node[anchor=north west,font=\scSmall,text=scInkM]
      at ($#2+(0.56,-0.245)$) {#5};
  \fi}

% --- shared pgfplots look (only if the figure loaded pgfplots) ------------
\ifdefined\pgfplotsset
  \pgfplotsset{compat=1.18}
  \pgfplotsset{
    synaxis/.style={
      scale only axis,
      tick align=outside, tick pos=left,
      axis line style={draw=scInkM!55,line width=0.4pt},
      % One grid, on the value axis only.  No minor grid: at these sizes it
      % turns the plot frame into a grey wash.
      ymajorgrids=true, xmajorgrids=false,
      major grid style={draw=scRuleL,line width=0.3pt},
      major tick length=1.6pt,
      label style={font=\scBody,color=scInk},
      tick label style={font=\scSmall,color=scInkM},
      legend style={font=\scSmall,draw=none,fill=none,inner sep=1pt,
                    row sep=0.4pt,legend columns=1},
      legend cell align=left,
      every axis plot/.append style={line width=0.7pt,mark size=1.2pt},
    },
  }
\fi
}{% ---------------------------------------------------------------------------
% synstyle.tex --- shared look for every TikZ figure of the SynKit CRN paper.
%
% Input it from a figure file:
%     \IfFileExists{synstyle.tex}{% ---------------------------------------------------------------------------
% synstyle.tex --- shared look for every TikZ figure of the SynKit CRN paper.
%
% Input it from a figure file:
%     \IfFileExists{synstyle.tex}{\input{synstyle}}{\input{Figure/synstyle}}
% so that the figure compiles both from Figure/ and from the paper root.
%
% The visual vocabulary follows ../../../../Style/style.tex: the
% colour-blind-safe Okabe--Ito accents, quiet neutral rules and wash, circular
% panel badges, and a compact technical type scale are shared by the companion
% papers.  Semantic aliases below keep the CRN meaning explicit.
%
% DESIGN CONTRACT -----------------------------------------------------------
% 1. Geometry.  Every figure is laid out on a canvas exactly \scW wide (the
%    sn-jnl text width, 372 pt) and is included unscaled, so one point of type
%    is the same physical size in all five figures.  Open each picture with
%    \scCanvas{<height in cm>} to pin the bounding box.
% 2. Colour is semantic, never decorative.  Four roles only:
%       scAccent (green)       structure, the normal/expected path
%       scMark   (orange)      catalyst, certificate, or highlight
%       scAlert  (vermillion)  failure, violation, or uncovered set
%       scRef    (blue)        secondary or reference series
%    A panel gets at most one accent.  Titles are always ink; panel identity
%    is also encoded by a letter inside a circular badge, not by colour alone.
% 3. Type scale.  Three sizes, no others, and nothing below 6 pt:
%       \scTitle (7 pt bold)  panel titles
%       \scBody  (6.5 pt)     labels, values, in-figure prose
%       \scSmall (6 pt)       units, footnotes, dense graph labels
% ---------------------------------------------------------------------------
% Load guard.  The paper inputs every figure source inline, so this file is
% reached once per figure; without the guard the second \newlength below
% would abort the run.
\ifdefined\SynKitStyleLoaded\endinput\fi
\newcommand{\SynKitStyleLoaded}{}

% Load in the document preamble. XeLaTeX/LuaLaTeX and an installed Arial
% family are required; fontspec reports an error instead of substituting.
\RequirePackage[no-math]{fontspec}
\newfontfamily\scArial{Arial}[NFSSFamily=scArial,Ligatures=TeX]
% Keep manuscript mathematics unchanged; select Arial letters and numerals
% only inside figures. TeX supplies specialist mathematical symbols.
\RequirePackage[subdued,italic]{mathastext}
\MTDeclareVersion{synarial}{TU}{scArial}{m}{n}
\SetMathAlphabet{\Mathsf}{synarial}{TU}{scArial}{m}{n}
\SetMathAlphabet{\Mathtt}{synarial}{TU}{scArial}{m}{n}

\usetikzlibrary{positioning,calc,fit,backgrounds,arrows.meta,shapes.geometric,
                shapes.misc,decorations.pathreplacing,patterns,math}

% --- canvas ---------------------------------------------------------------
% 372 pt text width minus the 1 pt standalone border on each side.
\newlength{\scW}\setlength{\scW}{13.0326cm}
\providecommand{\scWnum}{13.0326}
% \scCanvas{<height>} pins the picture to the full text width, y running down
% from 0, so every figure shares one coordinate convention.  It *sets* the
% bounding box rather than contributing to it, so the emitted page is exactly
% the text width and every figure is included at 1:1 -- otherwise a stray
% invisible extent (pgfplots reserves a little around an axis) rescales one
% figure relative to the others and the same \scSmall prints at two sizes.
% Anything drawn outside the canvas is therefore cropped: keep content inside,
% and check with
%     \providecommand{\scCanvas}[1]{\path[draw=red] (0,0) rectangle (\scWnum,-#1);}
\providecommand{\scCanvas}[1]{%
  \path[use as bounding box] (0,0) rectangle (\scWnum,-#1);}

% --- palette --------------------------------------------------------------
% Canonical companion-paper values (../../../../Style/style.tex).
\definecolor{npgOrange}{HTML}{E69F00}
\definecolor{npgSky}{HTML}{56B4E9}
\definecolor{npgGreen}{HTML}{009E73}
\definecolor{npgYellow}{HTML}{F0E442}
\definecolor{npgBlue}{HTML}{0072B2}
\definecolor{npgVermillion}{HTML}{E64B35}
\definecolor{npgPurple}{HTML}{CC79A7}
\definecolor{figInk}{HTML}{1A1A1A}
\definecolor{figMuted}{HTML}{666666}
\definecolor{figHair}{HTML}{D9D9D9}
\definecolor{figWash}{HTML}{F7F7F7}

\colorlet{scInk}{figInk}              % primary text
\colorlet{scInkM}{figMuted}           % secondary text
\colorlet{scPaper}{figWash}           % quiet cards and note strips
\colorlet{scSpF}{npgSky!12}           % species fill
\colorlet{scSpE}{npgBlue}             % species edge
\colorlet{scRxF}{npgOrange!12}        % reaction fill
\colorlet{scRxE}{npgOrange!88!black}  % reaction edge
\colorlet{scReac}{npgBlue!78!black}   % reactant edge
\colorlet{scProd}{npgGreen!82!black}  % product edge
\colorlet{scHiF}{npgOrange!13}        % highlight fill
\colorlet{scHiE}{npgOrange!82!black}  % highlight edge
% Stage / accent tints.  stDE remains available for the fifth plot series.
\colorlet{stAF}{npgGreen!9}\colorlet{stAE}{npgGreen}
\colorlet{stBF}{npgOrange!11}\colorlet{stBE}{npgOrange}
\colorlet{stCF}{npgBlue!9}\colorlet{stCE}{npgBlue}
\colorlet{stDF}{npgPurple!10}\colorlet{stDE}{npgPurple!85!black}
\colorlet{stEF}{npgVermillion!9}\colorlet{stEE}{npgVermillion}

% Semantic aliases.  Figures should speak in these, not in stAE/stBE/...
\colorlet{scAccent}{stAE}\colorlet{scAccentF}{stAF}   % structure / normal
\colorlet{scMark}{stBE}\colorlet{scMarkF}{stBF}       % highlight / certificate
\colorlet{scRef}{stCE}\colorlet{scRefF}{stCF}         % secondary series
\colorlet{scAlert}{stEE}\colorlet{scAlertF}{stEF}     % failure / violation
\colorlet{scRule}{figHair}                            % card and axis hairlines
\colorlet{scRuleL}{figHair!55}                        % grid, dividers

% --- text styles ----------------------------------------------------------
% Three sizes, fixed in points so they do not drift with the document class.
% Arial is shared by labels, headings, panel badges, and chemistry labels.
\providecommand{\scTitle}{\scArial\fontsize{7}{8.4}\selectfont\bfseries\color{scInk}}
\providecommand{\scBody}{\scArial\fontsize{6.5}{7.8}\selectfont\color{scInkM}}
\providecommand{\scSmall}{\scArial\fontsize{6}{7.2}\selectfont\color{scInkM}}
\providecommand{\scSub}[1]{{\scSmall\itshape #1}}
\providecommand{\scPanel}[1]{{\scTitle(#1)}}
% Numerals that carry a result: ink, not colour, so the eye reads the value.
\providecommand{\scNum}{\scArial\fontsize{10}{11}\selectfont\bfseries\color{scInk}}

% --- shared TikZ styles ---------------------------------------------------
\tikzset{
  % text never hyphenates inside a figure
  every node/.append style={execute at begin node={%
      \hyphenpenalty=10000 \exhyphenpenalty=10000 \hbadness=10000 }},
  syn/.style={font=\scArial,color=scInk,
              execute at begin picture={\MTversion*{synarial}},
              >={Stealth[round,length=2.4mm,width=1.9mm]}},
  % containers ------------------------------------------------------------
  % One card chrome for the whole paper: white, 3 pt radius, one hairline.
  panel/.style={rounded corners=2.5pt,fill=none,draw=scRule,line width=0.45pt},
  card/.style={rounded corners=2.2pt,fill=scPaper,draw=scRule,line width=0.45pt,
               inner sep=5pt,align=left,anchor=north west},
  % a quiet strip for a panel's closing fact
  note/.style={rounded corners=2.5pt,fill=scPaper,draw=scRuleL,line width=0.4pt,
               inner xsep=6pt,inner ysep=3.5pt,align=left,font=\scSmall},
  % small labelled value pill
  chip/.style={rounded corners=2pt,fill=scPaper,draw=scRule,line width=0.4pt,
               inner xsep=4pt,inner ysep=2pt,font=\scSmall},
  % Circular letter badge: the same accessible panel marker used in Style/*.
  hbullet/.style={circle,fill=scInk,text=white,draw=none,
                  font=\scArial\fontsize{6}{6}\selectfont\bfseries,
                  minimum size=4.2mm,inner sep=0pt},
  ctag/.style={font=\scTitle,anchor=north west,text=scInk},
  % CRN primitives --------------------------------------------------------
  species/.style={rounded rectangle,fill=scSpF,draw=scSpE,line width=0.45pt,
                  inner xsep=2.4pt,inner ysep=1.6pt,font=\scSmall,
                  text=scInk,minimum height=3.6mm},
  cofactor/.style={species,fill=scRefF,draw=scRef},
  catalyst/.style={species,fill=scHiF,draw=scHiE},
  rxn/.style={rectangle,fill=scRxF,draw=scRxE,line width=0.45pt,
              inner sep=1.2pt,font=\scSmall,text=scInk,
              minimum width=3.0mm,minimum height=3.0mm},
  place/.style={circle,fill=scSpF,draw=scSpE,line width=0.45pt,
                minimum size=4.0mm,inner sep=0pt},
  trans/.style={rectangle,fill=scRxE,draw=scRxE,minimum width=1.1mm,
                minimum height=3.6mm,inner sep=0pt},
  token/.style={circle,fill=scInk,minimum size=1.1mm,inner sep=0pt},
  % edges -----------------------------------------------------------------
  reac/.style={->,line width=0.5pt,draw=scReac,shorten >=0.8pt,shorten <=0.8pt},
  prod/.style={->,line width=0.5pt,draw=scProd,shorten >=0.8pt,shorten <=0.8pt},
  cat/.style={->,line width=0.5pt,draw=scHiE,shorten >=0.8pt,shorten <=0.8pt},
  flow/.style={-{Stealth[round,length=2.2mm,width=1.8mm]},line width=1.1pt,
               draw=scInkM!45},
  flab/.style={font=\scSmall,text=scInkM,inner sep=2pt},
  stoi/.style={font=\scSmall,text=scInkM,inner sep=0.8pt},
  % verdict marks: shape carries the meaning, colour only reinforces it
  vpass/.style={circle,fill=scAccentF,draw=scAccent,text=scAccent,
                font=\fontsize{6}{6}\selectfont\bfseries,
                minimum size=3.2mm,inner sep=0pt},
  vfail/.style={circle,fill=scAlertF,draw=scAlert,text=scAlert,
                font=\fontsize{6}{6}\selectfont\bfseries,
                minimum size=3.2mm,inner sep=0pt},
  callout/.style={font=\scSmall,text=scInkM,align=left},
}

% \scHead[<badge colour>]{<anchor coord>}{<letter>}{<title>}{<subtitle or empty>}
% One header treatment for every panel in the paper.
\providecommand{\scHead}[5][scInk]{%
  \node[hbullet,fill=#1,anchor=north west] (scbadge#3) at ($#2+(0,-0.015)$)
    {\MakeUppercase{#3}};
  \node[ctag,anchor=west] at ($(scbadge#3.east)+(0.14,0)$) {#4};
  \ifx\relax#5\relax\else
    \node[anchor=north west,font=\scSmall,text=scInkM]
      at ($#2+(0.56,-0.245)$) {#5};
  \fi}

% --- shared pgfplots look (only if the figure loaded pgfplots) ------------
\ifdefined\pgfplotsset
  \pgfplotsset{compat=1.18}
  \pgfplotsset{
    synaxis/.style={
      scale only axis,
      tick align=outside, tick pos=left,
      axis line style={draw=scInkM!55,line width=0.4pt},
      % One grid, on the value axis only.  No minor grid: at these sizes it
      % turns the plot frame into a grey wash.
      ymajorgrids=true, xmajorgrids=false,
      major grid style={draw=scRuleL,line width=0.3pt},
      major tick length=1.6pt,
      label style={font=\scBody,color=scInk},
      tick label style={font=\scSmall,color=scInkM},
      legend style={font=\scSmall,draw=none,fill=none,inner sep=1pt,
                    row sep=0.4pt,legend columns=1},
      legend cell align=left,
      every axis plot/.append style={line width=0.7pt,mark size=1.2pt},
    },
  }
\fi
}{% ---------------------------------------------------------------------------
% synstyle.tex --- shared look for every TikZ figure of the SynKit CRN paper.
%
% Input it from a figure file:
%     \IfFileExists{synstyle.tex}{\input{synstyle}}{\input{Figure/synstyle}}
% so that the figure compiles both from Figure/ and from the paper root.
%
% The visual vocabulary follows ../../../../Style/style.tex: the
% colour-blind-safe Okabe--Ito accents, quiet neutral rules and wash, circular
% panel badges, and a compact technical type scale are shared by the companion
% papers.  Semantic aliases below keep the CRN meaning explicit.
%
% DESIGN CONTRACT -----------------------------------------------------------
% 1. Geometry.  Every figure is laid out on a canvas exactly \scW wide (the
%    sn-jnl text width, 372 pt) and is included unscaled, so one point of type
%    is the same physical size in all five figures.  Open each picture with
%    \scCanvas{<height in cm>} to pin the bounding box.
% 2. Colour is semantic, never decorative.  Four roles only:
%       scAccent (green)       structure, the normal/expected path
%       scMark   (orange)      catalyst, certificate, or highlight
%       scAlert  (vermillion)  failure, violation, or uncovered set
%       scRef    (blue)        secondary or reference series
%    A panel gets at most one accent.  Titles are always ink; panel identity
%    is also encoded by a letter inside a circular badge, not by colour alone.
% 3. Type scale.  Three sizes, no others, and nothing below 6 pt:
%       \scTitle (7 pt bold)  panel titles
%       \scBody  (6.5 pt)     labels, values, in-figure prose
%       \scSmall (6 pt)       units, footnotes, dense graph labels
% ---------------------------------------------------------------------------
% Load guard.  The paper inputs every figure source inline, so this file is
% reached once per figure; without the guard the second \newlength below
% would abort the run.
\ifdefined\SynKitStyleLoaded\endinput\fi
\newcommand{\SynKitStyleLoaded}{}

% Load in the document preamble. XeLaTeX/LuaLaTeX and an installed Arial
% family are required; fontspec reports an error instead of substituting.
\RequirePackage[no-math]{fontspec}
\newfontfamily\scArial{Arial}[NFSSFamily=scArial,Ligatures=TeX]
% Keep manuscript mathematics unchanged; select Arial letters and numerals
% only inside figures. TeX supplies specialist mathematical symbols.
\RequirePackage[subdued,italic]{mathastext}
\MTDeclareVersion{synarial}{TU}{scArial}{m}{n}
\SetMathAlphabet{\Mathsf}{synarial}{TU}{scArial}{m}{n}
\SetMathAlphabet{\Mathtt}{synarial}{TU}{scArial}{m}{n}

\usetikzlibrary{positioning,calc,fit,backgrounds,arrows.meta,shapes.geometric,
                shapes.misc,decorations.pathreplacing,patterns,math}

% --- canvas ---------------------------------------------------------------
% 372 pt text width minus the 1 pt standalone border on each side.
\newlength{\scW}\setlength{\scW}{13.0326cm}
\providecommand{\scWnum}{13.0326}
% \scCanvas{<height>} pins the picture to the full text width, y running down
% from 0, so every figure shares one coordinate convention.  It *sets* the
% bounding box rather than contributing to it, so the emitted page is exactly
% the text width and every figure is included at 1:1 -- otherwise a stray
% invisible extent (pgfplots reserves a little around an axis) rescales one
% figure relative to the others and the same \scSmall prints at two sizes.
% Anything drawn outside the canvas is therefore cropped: keep content inside,
% and check with
%     \providecommand{\scCanvas}[1]{\path[draw=red] (0,0) rectangle (\scWnum,-#1);}
\providecommand{\scCanvas}[1]{%
  \path[use as bounding box] (0,0) rectangle (\scWnum,-#1);}

% --- palette --------------------------------------------------------------
% Canonical companion-paper values (../../../../Style/style.tex).
\definecolor{npgOrange}{HTML}{E69F00}
\definecolor{npgSky}{HTML}{56B4E9}
\definecolor{npgGreen}{HTML}{009E73}
\definecolor{npgYellow}{HTML}{F0E442}
\definecolor{npgBlue}{HTML}{0072B2}
\definecolor{npgVermillion}{HTML}{E64B35}
\definecolor{npgPurple}{HTML}{CC79A7}
\definecolor{figInk}{HTML}{1A1A1A}
\definecolor{figMuted}{HTML}{666666}
\definecolor{figHair}{HTML}{D9D9D9}
\definecolor{figWash}{HTML}{F7F7F7}

\colorlet{scInk}{figInk}              % primary text
\colorlet{scInkM}{figMuted}           % secondary text
\colorlet{scPaper}{figWash}           % quiet cards and note strips
\colorlet{scSpF}{npgSky!12}           % species fill
\colorlet{scSpE}{npgBlue}             % species edge
\colorlet{scRxF}{npgOrange!12}        % reaction fill
\colorlet{scRxE}{npgOrange!88!black}  % reaction edge
\colorlet{scReac}{npgBlue!78!black}   % reactant edge
\colorlet{scProd}{npgGreen!82!black}  % product edge
\colorlet{scHiF}{npgOrange!13}        % highlight fill
\colorlet{scHiE}{npgOrange!82!black}  % highlight edge
% Stage / accent tints.  stDE remains available for the fifth plot series.
\colorlet{stAF}{npgGreen!9}\colorlet{stAE}{npgGreen}
\colorlet{stBF}{npgOrange!11}\colorlet{stBE}{npgOrange}
\colorlet{stCF}{npgBlue!9}\colorlet{stCE}{npgBlue}
\colorlet{stDF}{npgPurple!10}\colorlet{stDE}{npgPurple!85!black}
\colorlet{stEF}{npgVermillion!9}\colorlet{stEE}{npgVermillion}

% Semantic aliases.  Figures should speak in these, not in stAE/stBE/...
\colorlet{scAccent}{stAE}\colorlet{scAccentF}{stAF}   % structure / normal
\colorlet{scMark}{stBE}\colorlet{scMarkF}{stBF}       % highlight / certificate
\colorlet{scRef}{stCE}\colorlet{scRefF}{stCF}         % secondary series
\colorlet{scAlert}{stEE}\colorlet{scAlertF}{stEF}     % failure / violation
\colorlet{scRule}{figHair}                            % card and axis hairlines
\colorlet{scRuleL}{figHair!55}                        % grid, dividers

% --- text styles ----------------------------------------------------------
% Three sizes, fixed in points so they do not drift with the document class.
% Arial is shared by labels, headings, panel badges, and chemistry labels.
\providecommand{\scTitle}{\scArial\fontsize{7}{8.4}\selectfont\bfseries\color{scInk}}
\providecommand{\scBody}{\scArial\fontsize{6.5}{7.8}\selectfont\color{scInkM}}
\providecommand{\scSmall}{\scArial\fontsize{6}{7.2}\selectfont\color{scInkM}}
\providecommand{\scSub}[1]{{\scSmall\itshape #1}}
\providecommand{\scPanel}[1]{{\scTitle(#1)}}
% Numerals that carry a result: ink, not colour, so the eye reads the value.
\providecommand{\scNum}{\scArial\fontsize{10}{11}\selectfont\bfseries\color{scInk}}

% --- shared TikZ styles ---------------------------------------------------
\tikzset{
  % text never hyphenates inside a figure
  every node/.append style={execute at begin node={%
      \hyphenpenalty=10000 \exhyphenpenalty=10000 \hbadness=10000 }},
  syn/.style={font=\scArial,color=scInk,
              execute at begin picture={\MTversion*{synarial}},
              >={Stealth[round,length=2.4mm,width=1.9mm]}},
  % containers ------------------------------------------------------------
  % One card chrome for the whole paper: white, 3 pt radius, one hairline.
  panel/.style={rounded corners=2.5pt,fill=none,draw=scRule,line width=0.45pt},
  card/.style={rounded corners=2.2pt,fill=scPaper,draw=scRule,line width=0.45pt,
               inner sep=5pt,align=left,anchor=north west},
  % a quiet strip for a panel's closing fact
  note/.style={rounded corners=2.5pt,fill=scPaper,draw=scRuleL,line width=0.4pt,
               inner xsep=6pt,inner ysep=3.5pt,align=left,font=\scSmall},
  % small labelled value pill
  chip/.style={rounded corners=2pt,fill=scPaper,draw=scRule,line width=0.4pt,
               inner xsep=4pt,inner ysep=2pt,font=\scSmall},
  % Circular letter badge: the same accessible panel marker used in Style/*.
  hbullet/.style={circle,fill=scInk,text=white,draw=none,
                  font=\scArial\fontsize{6}{6}\selectfont\bfseries,
                  minimum size=4.2mm,inner sep=0pt},
  ctag/.style={font=\scTitle,anchor=north west,text=scInk},
  % CRN primitives --------------------------------------------------------
  species/.style={rounded rectangle,fill=scSpF,draw=scSpE,line width=0.45pt,
                  inner xsep=2.4pt,inner ysep=1.6pt,font=\scSmall,
                  text=scInk,minimum height=3.6mm},
  cofactor/.style={species,fill=scRefF,draw=scRef},
  catalyst/.style={species,fill=scHiF,draw=scHiE},
  rxn/.style={rectangle,fill=scRxF,draw=scRxE,line width=0.45pt,
              inner sep=1.2pt,font=\scSmall,text=scInk,
              minimum width=3.0mm,minimum height=3.0mm},
  place/.style={circle,fill=scSpF,draw=scSpE,line width=0.45pt,
                minimum size=4.0mm,inner sep=0pt},
  trans/.style={rectangle,fill=scRxE,draw=scRxE,minimum width=1.1mm,
                minimum height=3.6mm,inner sep=0pt},
  token/.style={circle,fill=scInk,minimum size=1.1mm,inner sep=0pt},
  % edges -----------------------------------------------------------------
  reac/.style={->,line width=0.5pt,draw=scReac,shorten >=0.8pt,shorten <=0.8pt},
  prod/.style={->,line width=0.5pt,draw=scProd,shorten >=0.8pt,shorten <=0.8pt},
  cat/.style={->,line width=0.5pt,draw=scHiE,shorten >=0.8pt,shorten <=0.8pt},
  flow/.style={-{Stealth[round,length=2.2mm,width=1.8mm]},line width=1.1pt,
               draw=scInkM!45},
  flab/.style={font=\scSmall,text=scInkM,inner sep=2pt},
  stoi/.style={font=\scSmall,text=scInkM,inner sep=0.8pt},
  % verdict marks: shape carries the meaning, colour only reinforces it
  vpass/.style={circle,fill=scAccentF,draw=scAccent,text=scAccent,
                font=\fontsize{6}{6}\selectfont\bfseries,
                minimum size=3.2mm,inner sep=0pt},
  vfail/.style={circle,fill=scAlertF,draw=scAlert,text=scAlert,
                font=\fontsize{6}{6}\selectfont\bfseries,
                minimum size=3.2mm,inner sep=0pt},
  callout/.style={font=\scSmall,text=scInkM,align=left},
}

% \scHead[<badge colour>]{<anchor coord>}{<letter>}{<title>}{<subtitle or empty>}
% One header treatment for every panel in the paper.
\providecommand{\scHead}[5][scInk]{%
  \node[hbullet,fill=#1,anchor=north west] (scbadge#3) at ($#2+(0,-0.015)$)
    {\MakeUppercase{#3}};
  \node[ctag,anchor=west] at ($(scbadge#3.east)+(0.14,0)$) {#4};
  \ifx\relax#5\relax\else
    \node[anchor=north west,font=\scSmall,text=scInkM]
      at ($#2+(0.56,-0.245)$) {#5};
  \fi}

% --- shared pgfplots look (only if the figure loaded pgfplots) ------------
\ifdefined\pgfplotsset
  \pgfplotsset{compat=1.18}
  \pgfplotsset{
    synaxis/.style={
      scale only axis,
      tick align=outside, tick pos=left,
      axis line style={draw=scInkM!55,line width=0.4pt},
      % One grid, on the value axis only.  No minor grid: at these sizes it
      % turns the plot frame into a grey wash.
      ymajorgrids=true, xmajorgrids=false,
      major grid style={draw=scRuleL,line width=0.3pt},
      major tick length=1.6pt,
      label style={font=\scBody,color=scInk},
      tick label style={font=\scSmall,color=scInkM},
      legend style={font=\scSmall,draw=none,fill=none,inner sep=1pt,
                    row sep=0.4pt,legend columns=1},
      legend cell align=left,
      every axis plot/.append style={line width=0.7pt,mark size=1.2pt},
    },
  }
\fi
}
% so that the figure compiles both from Figure/ and from the paper root.
%
% The visual vocabulary follows ../../../../Style/style.tex: the
% colour-blind-safe Okabe--Ito accents, quiet neutral rules and wash, circular
% panel badges, and a compact technical type scale are shared by the companion
% papers.  Semantic aliases below keep the CRN meaning explicit.
%
% DESIGN CONTRACT -----------------------------------------------------------
% 1. Geometry.  Every figure is laid out on a canvas exactly \scW wide (the
%    sn-jnl text width, 372 pt) and is included unscaled, so one point of type
%    is the same physical size in all five figures.  Open each picture with
%    \scCanvas{<height in cm>} to pin the bounding box.
% 2. Colour is semantic, never decorative.  Four roles only:
%       scAccent (green)       structure, the normal/expected path
%       scMark   (orange)      catalyst, certificate, or highlight
%       scAlert  (vermillion)  failure, violation, or uncovered set
%       scRef    (blue)        secondary or reference series
%    A panel gets at most one accent.  Titles are always ink; panel identity
%    is also encoded by a letter inside a circular badge, not by colour alone.
% 3. Type scale.  Three sizes, no others, and nothing below 6 pt:
%       \scTitle (7 pt bold)  panel titles
%       \scBody  (6.5 pt)     labels, values, in-figure prose
%       \scSmall (6 pt)       units, footnotes, dense graph labels
% ---------------------------------------------------------------------------
% Load guard.  The paper inputs every figure source inline, so this file is
% reached once per figure; without the guard the second \newlength below
% would abort the run.
\ifdefined\SynKitStyleLoaded\endinput\fi
\newcommand{\SynKitStyleLoaded}{}

% Load in the document preamble. XeLaTeX/LuaLaTeX and an installed Arial
% family are required; fontspec reports an error instead of substituting.
\RequirePackage[no-math]{fontspec}
\newfontfamily\scArial{Arial}[NFSSFamily=scArial,Ligatures=TeX]
% Keep manuscript mathematics unchanged; select Arial letters and numerals
% only inside figures. TeX supplies specialist mathematical symbols.
\RequirePackage[subdued,italic]{mathastext}
\MTDeclareVersion{synarial}{TU}{scArial}{m}{n}
\SetMathAlphabet{\Mathsf}{synarial}{TU}{scArial}{m}{n}
\SetMathAlphabet{\Mathtt}{synarial}{TU}{scArial}{m}{n}

\usetikzlibrary{positioning,calc,fit,backgrounds,arrows.meta,shapes.geometric,
                shapes.misc,decorations.pathreplacing,patterns,math}

% --- canvas ---------------------------------------------------------------
% 372 pt text width minus the 1 pt standalone border on each side.
\newlength{\scW}\setlength{\scW}{13.0326cm}
\providecommand{\scWnum}{13.0326}
% \scCanvas{<height>} pins the picture to the full text width, y running down
% from 0, so every figure shares one coordinate convention.  It *sets* the
% bounding box rather than contributing to it, so the emitted page is exactly
% the text width and every figure is included at 1:1 -- otherwise a stray
% invisible extent (pgfplots reserves a little around an axis) rescales one
% figure relative to the others and the same \scSmall prints at two sizes.
% Anything drawn outside the canvas is therefore cropped: keep content inside,
% and check with
%     \providecommand{\scCanvas}[1]{\path[draw=red] (0,0) rectangle (\scWnum,-#1);}
\providecommand{\scCanvas}[1]{%
  \path[use as bounding box] (0,0) rectangle (\scWnum,-#1);}

% --- palette --------------------------------------------------------------
% Canonical companion-paper values (../../../../Style/style.tex).
\definecolor{npgOrange}{HTML}{E69F00}
\definecolor{npgSky}{HTML}{56B4E9}
\definecolor{npgGreen}{HTML}{009E73}
\definecolor{npgYellow}{HTML}{F0E442}
\definecolor{npgBlue}{HTML}{0072B2}
\definecolor{npgVermillion}{HTML}{E64B35}
\definecolor{npgPurple}{HTML}{CC79A7}
\definecolor{figInk}{HTML}{1A1A1A}
\definecolor{figMuted}{HTML}{666666}
\definecolor{figHair}{HTML}{D9D9D9}
\definecolor{figWash}{HTML}{F7F7F7}

\colorlet{scInk}{figInk}              % primary text
\colorlet{scInkM}{figMuted}           % secondary text
\colorlet{scPaper}{figWash}           % quiet cards and note strips
\colorlet{scSpF}{npgSky!12}           % species fill
\colorlet{scSpE}{npgBlue}             % species edge
\colorlet{scRxF}{npgOrange!12}        % reaction fill
\colorlet{scRxE}{npgOrange!88!black}  % reaction edge
\colorlet{scReac}{npgBlue!78!black}   % reactant edge
\colorlet{scProd}{npgGreen!82!black}  % product edge
\colorlet{scHiF}{npgOrange!13}        % highlight fill
\colorlet{scHiE}{npgOrange!82!black}  % highlight edge
% Stage / accent tints.  stDE remains available for the fifth plot series.
\colorlet{stAF}{npgGreen!9}\colorlet{stAE}{npgGreen}
\colorlet{stBF}{npgOrange!11}\colorlet{stBE}{npgOrange}
\colorlet{stCF}{npgBlue!9}\colorlet{stCE}{npgBlue}
\colorlet{stDF}{npgPurple!10}\colorlet{stDE}{npgPurple!85!black}
\colorlet{stEF}{npgVermillion!9}\colorlet{stEE}{npgVermillion}

% Semantic aliases.  Figures should speak in these, not in stAE/stBE/...
\colorlet{scAccent}{stAE}\colorlet{scAccentF}{stAF}   % structure / normal
\colorlet{scMark}{stBE}\colorlet{scMarkF}{stBF}       % highlight / certificate
\colorlet{scRef}{stCE}\colorlet{scRefF}{stCF}         % secondary series
\colorlet{scAlert}{stEE}\colorlet{scAlertF}{stEF}     % failure / violation
\colorlet{scRule}{figHair}                            % card and axis hairlines
\colorlet{scRuleL}{figHair!55}                        % grid, dividers

% --- text styles ----------------------------------------------------------
% Three sizes, fixed in points so they do not drift with the document class.
% Arial is shared by labels, headings, panel badges, and chemistry labels.
\providecommand{\scTitle}{\scArial\fontsize{7}{8.4}\selectfont\bfseries\color{scInk}}
\providecommand{\scBody}{\scArial\fontsize{6.5}{7.8}\selectfont\color{scInkM}}
\providecommand{\scSmall}{\scArial\fontsize{6}{7.2}\selectfont\color{scInkM}}
\providecommand{\scSub}[1]{{\scSmall\itshape #1}}
\providecommand{\scPanel}[1]{{\scTitle(#1)}}
% Numerals that carry a result: ink, not colour, so the eye reads the value.
\providecommand{\scNum}{\scArial\fontsize{10}{11}\selectfont\bfseries\color{scInk}}

% --- shared TikZ styles ---------------------------------------------------
\tikzset{
  % text never hyphenates inside a figure
  every node/.append style={execute at begin node={%
      \hyphenpenalty=10000 \exhyphenpenalty=10000 \hbadness=10000 }},
  syn/.style={font=\scArial,color=scInk,
              execute at begin picture={\MTversion*{synarial}},
              >={Stealth[round,length=2.4mm,width=1.9mm]}},
  % containers ------------------------------------------------------------
  % One card chrome for the whole paper: white, 3 pt radius, one hairline.
  panel/.style={rounded corners=2.5pt,fill=none,draw=scRule,line width=0.45pt},
  card/.style={rounded corners=2.2pt,fill=scPaper,draw=scRule,line width=0.45pt,
               inner sep=5pt,align=left,anchor=north west},
  % a quiet strip for a panel's closing fact
  note/.style={rounded corners=2.5pt,fill=scPaper,draw=scRuleL,line width=0.4pt,
               inner xsep=6pt,inner ysep=3.5pt,align=left,font=\scSmall},
  % small labelled value pill
  chip/.style={rounded corners=2pt,fill=scPaper,draw=scRule,line width=0.4pt,
               inner xsep=4pt,inner ysep=2pt,font=\scSmall},
  % Circular letter badge: the same accessible panel marker used in Style/*.
  hbullet/.style={circle,fill=scInk,text=white,draw=none,
                  font=\scArial\fontsize{6}{6}\selectfont\bfseries,
                  minimum size=4.2mm,inner sep=0pt},
  ctag/.style={font=\scTitle,anchor=north west,text=scInk},
  % CRN primitives --------------------------------------------------------
  species/.style={rounded rectangle,fill=scSpF,draw=scSpE,line width=0.45pt,
                  inner xsep=2.4pt,inner ysep=1.6pt,font=\scSmall,
                  text=scInk,minimum height=3.6mm},
  cofactor/.style={species,fill=scRefF,draw=scRef},
  catalyst/.style={species,fill=scHiF,draw=scHiE},
  rxn/.style={rectangle,fill=scRxF,draw=scRxE,line width=0.45pt,
              inner sep=1.2pt,font=\scSmall,text=scInk,
              minimum width=3.0mm,minimum height=3.0mm},
  place/.style={circle,fill=scSpF,draw=scSpE,line width=0.45pt,
                minimum size=4.0mm,inner sep=0pt},
  trans/.style={rectangle,fill=scRxE,draw=scRxE,minimum width=1.1mm,
                minimum height=3.6mm,inner sep=0pt},
  token/.style={circle,fill=scInk,minimum size=1.1mm,inner sep=0pt},
  % edges -----------------------------------------------------------------
  reac/.style={->,line width=0.5pt,draw=scReac,shorten >=0.8pt,shorten <=0.8pt},
  prod/.style={->,line width=0.5pt,draw=scProd,shorten >=0.8pt,shorten <=0.8pt},
  cat/.style={->,line width=0.5pt,draw=scHiE,shorten >=0.8pt,shorten <=0.8pt},
  flow/.style={-{Stealth[round,length=2.2mm,width=1.8mm]},line width=1.1pt,
               draw=scInkM!45},
  flab/.style={font=\scSmall,text=scInkM,inner sep=2pt},
  stoi/.style={font=\scSmall,text=scInkM,inner sep=0.8pt},
  % verdict marks: shape carries the meaning, colour only reinforces it
  vpass/.style={circle,fill=scAccentF,draw=scAccent,text=scAccent,
                font=\fontsize{6}{6}\selectfont\bfseries,
                minimum size=3.2mm,inner sep=0pt},
  vfail/.style={circle,fill=scAlertF,draw=scAlert,text=scAlert,
                font=\fontsize{6}{6}\selectfont\bfseries,
                minimum size=3.2mm,inner sep=0pt},
  callout/.style={font=\scSmall,text=scInkM,align=left},
}

% \scHead[<badge colour>]{<anchor coord>}{<letter>}{<title>}{<subtitle or empty>}
% One header treatment for every panel in the paper.
\providecommand{\scHead}[5][scInk]{%
  \node[hbullet,fill=#1,anchor=north west] (scbadge#3) at ($#2+(0,-0.015)$)
    {\MakeUppercase{#3}};
  \node[ctag,anchor=west] at ($(scbadge#3.east)+(0.14,0)$) {#4};
  \ifx\relax#5\relax\else
    \node[anchor=north west,font=\scSmall,text=scInkM]
      at ($#2+(0.56,-0.245)$) {#5};
  \fi}

% --- shared pgfplots look (only if the figure loaded pgfplots) ------------
\ifdefined\pgfplotsset
  \pgfplotsset{compat=1.18}
  \pgfplotsset{
    synaxis/.style={
      scale only axis,
      tick align=outside, tick pos=left,
      axis line style={draw=scInkM!55,line width=0.4pt},
      % One grid, on the value axis only.  No minor grid: at these sizes it
      % turns the plot frame into a grey wash.
      ymajorgrids=true, xmajorgrids=false,
      major grid style={draw=scRuleL,line width=0.3pt},
      major tick length=1.6pt,
      label style={font=\scBody,color=scInk},
      tick label style={font=\scSmall,color=scInkM},
      legend style={font=\scSmall,draw=none,fill=none,inner sep=1pt,
                    row sep=0.4pt,legend columns=1},
      legend cell align=left,
      every axis plot/.append style={line width=0.7pt,mark size=1.2pt},
    },
  }
\fi
}
% so that the figure compiles both from Figure/ and from the paper root.
%
% The visual vocabulary follows ../../../../Style/style.tex: the
% colour-blind-safe Okabe--Ito accents, quiet neutral rules and wash, circular
% panel badges, and a compact technical type scale are shared by the companion
% papers.  Semantic aliases below keep the CRN meaning explicit.
%
% DESIGN CONTRACT -----------------------------------------------------------
% 1. Geometry.  Every figure is laid out on a canvas exactly \scW wide (the
%    sn-jnl text width, 372 pt) and is included unscaled, so one point of type
%    is the same physical size in all five figures.  Open each picture with
%    \scCanvas{<height in cm>} to pin the bounding box.
% 2. Colour is semantic, never decorative.  Four roles only:
%       scAccent (green)       structure, the normal/expected path
%       scMark   (orange)      catalyst, certificate, or highlight
%       scAlert  (vermillion)  failure, violation, or uncovered set
%       scRef    (blue)        secondary or reference series
%    A panel gets at most one accent.  Titles are always ink; panel identity
%    is also encoded by a letter inside a circular badge, not by colour alone.
% 3. Type scale.  Three sizes, no others, and nothing below 6 pt:
%       \scTitle (7 pt bold)  panel titles
%       \scBody  (6.5 pt)     labels, values, in-figure prose
%       \scSmall (6 pt)       units, footnotes, dense graph labels
% ---------------------------------------------------------------------------
% Load guard.  The paper inputs every figure source inline, so this file is
% reached once per figure; without the guard the second \newlength below
% would abort the run.
\ifdefined\SynKitStyleLoaded\endinput\fi
\newcommand{\SynKitStyleLoaded}{}

% Load in the document preamble. XeLaTeX/LuaLaTeX and an installed Arial
% family are required; fontspec reports an error instead of substituting.
\RequirePackage[no-math]{fontspec}
\newfontfamily\scArial{Arial}[NFSSFamily=scArial,Ligatures=TeX]
% Keep manuscript mathematics unchanged; select Arial letters and numerals
% only inside figures. TeX supplies specialist mathematical symbols.
\RequirePackage[subdued,italic]{mathastext}
\MTDeclareVersion{synarial}{TU}{scArial}{m}{n}
\SetMathAlphabet{\Mathsf}{synarial}{TU}{scArial}{m}{n}
\SetMathAlphabet{\Mathtt}{synarial}{TU}{scArial}{m}{n}

\usetikzlibrary{positioning,calc,fit,backgrounds,arrows.meta,shapes.geometric,
                shapes.misc,decorations.pathreplacing,patterns,math}

% --- canvas ---------------------------------------------------------------
% 372 pt text width minus the 1 pt standalone border on each side.
\newlength{\scW}\setlength{\scW}{13.0326cm}
\providecommand{\scWnum}{13.0326}
% \scCanvas{<height>} pins the picture to the full text width, y running down
% from 0, so every figure shares one coordinate convention.  It *sets* the
% bounding box rather than contributing to it, so the emitted page is exactly
% the text width and every figure is included at 1:1 -- otherwise a stray
% invisible extent (pgfplots reserves a little around an axis) rescales one
% figure relative to the others and the same \scSmall prints at two sizes.
% Anything drawn outside the canvas is therefore cropped: keep content inside,
% and check with
%     \providecommand{\scCanvas}[1]{\path[draw=red] (0,0) rectangle (\scWnum,-#1);}
\providecommand{\scCanvas}[1]{%
  \path[use as bounding box] (0,0) rectangle (\scWnum,-#1);}

% --- palette --------------------------------------------------------------
% Canonical companion-paper values (../../../../Style/style.tex).
\definecolor{npgOrange}{HTML}{E69F00}
\definecolor{npgSky}{HTML}{56B4E9}
\definecolor{npgGreen}{HTML}{009E73}
\definecolor{npgYellow}{HTML}{F0E442}
\definecolor{npgBlue}{HTML}{0072B2}
\definecolor{npgVermillion}{HTML}{E64B35}
\definecolor{npgPurple}{HTML}{CC79A7}
\definecolor{figInk}{HTML}{1A1A1A}
\definecolor{figMuted}{HTML}{666666}
\definecolor{figHair}{HTML}{D9D9D9}
\definecolor{figWash}{HTML}{F7F7F7}

\colorlet{scInk}{figInk}              % primary text
\colorlet{scInkM}{figMuted}           % secondary text
\colorlet{scPaper}{figWash}           % quiet cards and note strips
\colorlet{scSpF}{npgSky!12}           % species fill
\colorlet{scSpE}{npgBlue}             % species edge
\colorlet{scRxF}{npgOrange!12}        % reaction fill
\colorlet{scRxE}{npgOrange!88!black}  % reaction edge
\colorlet{scReac}{npgBlue!78!black}   % reactant edge
\colorlet{scProd}{npgGreen!82!black}  % product edge
\colorlet{scHiF}{npgOrange!13}        % highlight fill
\colorlet{scHiE}{npgOrange!82!black}  % highlight edge
% Stage / accent tints.  stDE remains available for the fifth plot series.
\colorlet{stAF}{npgGreen!9}\colorlet{stAE}{npgGreen}
\colorlet{stBF}{npgOrange!11}\colorlet{stBE}{npgOrange}
\colorlet{stCF}{npgBlue!9}\colorlet{stCE}{npgBlue}
\colorlet{stDF}{npgPurple!10}\colorlet{stDE}{npgPurple!85!black}
\colorlet{stEF}{npgVermillion!9}\colorlet{stEE}{npgVermillion}

% Semantic aliases.  Figures should speak in these, not in stAE/stBE/...
\colorlet{scAccent}{stAE}\colorlet{scAccentF}{stAF}   % structure / normal
\colorlet{scMark}{stBE}\colorlet{scMarkF}{stBF}       % highlight / certificate
\colorlet{scRef}{stCE}\colorlet{scRefF}{stCF}         % secondary series
\colorlet{scAlert}{stEE}\colorlet{scAlertF}{stEF}     % failure / violation
\colorlet{scRule}{figHair}                            % card and axis hairlines
\colorlet{scRuleL}{figHair!55}                        % grid, dividers

% --- text styles ----------------------------------------------------------
% Three sizes, fixed in points so they do not drift with the document class.
% Arial is shared by labels, headings, panel badges, and chemistry labels.
\providecommand{\scTitle}{\scArial\fontsize{7}{8.4}\selectfont\bfseries\color{scInk}}
\providecommand{\scBody}{\scArial\fontsize{6.5}{7.8}\selectfont\color{scInkM}}
\providecommand{\scSmall}{\scArial\fontsize{6}{7.2}\selectfont\color{scInkM}}
\providecommand{\scSub}[1]{{\scSmall\itshape #1}}
\providecommand{\scPanel}[1]{{\scTitle(#1)}}
% Numerals that carry a result: ink, not colour, so the eye reads the value.
\providecommand{\scNum}{\scArial\fontsize{10}{11}\selectfont\bfseries\color{scInk}}

% --- shared TikZ styles ---------------------------------------------------
\tikzset{
  % text never hyphenates inside a figure
  every node/.append style={execute at begin node={%
      \hyphenpenalty=10000 \exhyphenpenalty=10000 \hbadness=10000 }},
  syn/.style={font=\scArial,color=scInk,
              execute at begin picture={\MTversion*{synarial}},
              >={Stealth[round,length=2.4mm,width=1.9mm]}},
  % containers ------------------------------------------------------------
  % One card chrome for the whole paper: white, 3 pt radius, one hairline.
  panel/.style={rounded corners=2.5pt,fill=none,draw=scRule,line width=0.45pt},
  card/.style={rounded corners=2.2pt,fill=scPaper,draw=scRule,line width=0.45pt,
               inner sep=5pt,align=left,anchor=north west},
  % a quiet strip for a panel's closing fact
  note/.style={rounded corners=2.5pt,fill=scPaper,draw=scRuleL,line width=0.4pt,
               inner xsep=6pt,inner ysep=3.5pt,align=left,font=\scSmall},
  % small labelled value pill
  chip/.style={rounded corners=2pt,fill=scPaper,draw=scRule,line width=0.4pt,
               inner xsep=4pt,inner ysep=2pt,font=\scSmall},
  % Circular letter badge: the same accessible panel marker used in Style/*.
  hbullet/.style={circle,fill=scInk,text=white,draw=none,
                  font=\scArial\fontsize{6}{6}\selectfont\bfseries,
                  minimum size=4.2mm,inner sep=0pt},
  ctag/.style={font=\scTitle,anchor=north west,text=scInk},
  % CRN primitives --------------------------------------------------------
  species/.style={rounded rectangle,fill=scSpF,draw=scSpE,line width=0.45pt,
                  inner xsep=2.4pt,inner ysep=1.6pt,font=\scSmall,
                  text=scInk,minimum height=3.6mm},
  cofactor/.style={species,fill=scRefF,draw=scRef},
  catalyst/.style={species,fill=scHiF,draw=scHiE},
  rxn/.style={rectangle,fill=scRxF,draw=scRxE,line width=0.45pt,
              inner sep=1.2pt,font=\scSmall,text=scInk,
              minimum width=3.0mm,minimum height=3.0mm},
  place/.style={circle,fill=scSpF,draw=scSpE,line width=0.45pt,
                minimum size=4.0mm,inner sep=0pt},
  trans/.style={rectangle,fill=scRxE,draw=scRxE,minimum width=1.1mm,
                minimum height=3.6mm,inner sep=0pt},
  token/.style={circle,fill=scInk,minimum size=1.1mm,inner sep=0pt},
  % edges -----------------------------------------------------------------
  reac/.style={->,line width=0.5pt,draw=scReac,shorten >=0.8pt,shorten <=0.8pt},
  prod/.style={->,line width=0.5pt,draw=scProd,shorten >=0.8pt,shorten <=0.8pt},
  cat/.style={->,line width=0.5pt,draw=scHiE,shorten >=0.8pt,shorten <=0.8pt},
  flow/.style={-{Stealth[round,length=2.2mm,width=1.8mm]},line width=1.1pt,
               draw=scInkM!45},
  flab/.style={font=\scSmall,text=scInkM,inner sep=2pt},
  stoi/.style={font=\scSmall,text=scInkM,inner sep=0.8pt},
  % verdict marks: shape carries the meaning, colour only reinforces it
  vpass/.style={circle,fill=scAccentF,draw=scAccent,text=scAccent,
                font=\fontsize{6}{6}\selectfont\bfseries,
                minimum size=3.2mm,inner sep=0pt},
  vfail/.style={circle,fill=scAlertF,draw=scAlert,text=scAlert,
                font=\fontsize{6}{6}\selectfont\bfseries,
                minimum size=3.2mm,inner sep=0pt},
  callout/.style={font=\scSmall,text=scInkM,align=left},
}

% \scHead[<badge colour>]{<anchor coord>}{<letter>}{<title>}{<subtitle or empty>}
% One header treatment for every panel in the paper.
\providecommand{\scHead}[5][scInk]{%
  \node[hbullet,fill=#1,anchor=north west] (scbadge#3) at ($#2+(0,-0.015)$)
    {\MakeUppercase{#3}};
  \node[ctag,anchor=west] at ($(scbadge#3.east)+(0.14,0)$) {#4};
  \ifx\relax#5\relax\else
    \node[anchor=north west,font=\scSmall,text=scInkM]
      at ($#2+(0.56,-0.245)$) {#5};
  \fi}

% --- shared pgfplots look (only if the figure loaded pgfplots) ------------
\ifdefined\pgfplotsset
  \pgfplotsset{compat=1.18}
  \pgfplotsset{
    synaxis/.style={
      scale only axis,
      tick align=outside, tick pos=left,
      axis line style={draw=scInkM!55,line width=0.4pt},
      % One grid, on the value axis only.  No minor grid: at these sizes it
      % turns the plot frame into a grey wash.
      ymajorgrids=true, xmajorgrids=false,
      major grid style={draw=scRuleL,line width=0.3pt},
      major tick length=1.6pt,
      label style={font=\scBody,color=scInk},
      tick label style={font=\scSmall,color=scInkM},
      legend style={font=\scSmall,draw=none,fill=none,inner sep=1pt,
                    row sep=0.4pt,legend columns=1},
      legend cell align=left,
      every axis plot/.append style={line width=0.7pt,mark size=1.2pt},
    },
  }
\fi
}{%
    % ---------------------------------------------------------------------------
% synstyle.tex --- shared look for every TikZ figure of the SynKit CRN paper.
%
% Input it from a figure file:
%     \IfFileExists{synstyle.tex}{% ---------------------------------------------------------------------------
% synstyle.tex --- shared look for every TikZ figure of the SynKit CRN paper.
%
% Input it from a figure file:
%     \IfFileExists{synstyle.tex}{% ---------------------------------------------------------------------------
% synstyle.tex --- shared look for every TikZ figure of the SynKit CRN paper.
%
% Input it from a figure file:
%     \IfFileExists{synstyle.tex}{\input{synstyle}}{\input{Figure/synstyle}}
% so that the figure compiles both from Figure/ and from the paper root.
%
% The visual vocabulary follows ../../../../Style/style.tex: the
% colour-blind-safe Okabe--Ito accents, quiet neutral rules and wash, circular
% panel badges, and a compact technical type scale are shared by the companion
% papers.  Semantic aliases below keep the CRN meaning explicit.
%
% DESIGN CONTRACT -----------------------------------------------------------
% 1. Geometry.  Every figure is laid out on a canvas exactly \scW wide (the
%    sn-jnl text width, 372 pt) and is included unscaled, so one point of type
%    is the same physical size in all five figures.  Open each picture with
%    \scCanvas{<height in cm>} to pin the bounding box.
% 2. Colour is semantic, never decorative.  Four roles only:
%       scAccent (green)       structure, the normal/expected path
%       scMark   (orange)      catalyst, certificate, or highlight
%       scAlert  (vermillion)  failure, violation, or uncovered set
%       scRef    (blue)        secondary or reference series
%    A panel gets at most one accent.  Titles are always ink; panel identity
%    is also encoded by a letter inside a circular badge, not by colour alone.
% 3. Type scale.  Three sizes, no others, and nothing below 6 pt:
%       \scTitle (7 pt bold)  panel titles
%       \scBody  (6.5 pt)     labels, values, in-figure prose
%       \scSmall (6 pt)       units, footnotes, dense graph labels
% ---------------------------------------------------------------------------
% Load guard.  The paper inputs every figure source inline, so this file is
% reached once per figure; without the guard the second \newlength below
% would abort the run.
\ifdefined\SynKitStyleLoaded\endinput\fi
\newcommand{\SynKitStyleLoaded}{}

% Load in the document preamble. XeLaTeX/LuaLaTeX and an installed Arial
% family are required; fontspec reports an error instead of substituting.
\RequirePackage[no-math]{fontspec}
\newfontfamily\scArial{Arial}[NFSSFamily=scArial,Ligatures=TeX]
% Keep manuscript mathematics unchanged; select Arial letters and numerals
% only inside figures. TeX supplies specialist mathematical symbols.
\RequirePackage[subdued,italic]{mathastext}
\MTDeclareVersion{synarial}{TU}{scArial}{m}{n}
\SetMathAlphabet{\Mathsf}{synarial}{TU}{scArial}{m}{n}
\SetMathAlphabet{\Mathtt}{synarial}{TU}{scArial}{m}{n}

\usetikzlibrary{positioning,calc,fit,backgrounds,arrows.meta,shapes.geometric,
                shapes.misc,decorations.pathreplacing,patterns,math}

% --- canvas ---------------------------------------------------------------
% 372 pt text width minus the 1 pt standalone border on each side.
\newlength{\scW}\setlength{\scW}{13.0326cm}
\providecommand{\scWnum}{13.0326}
% \scCanvas{<height>} pins the picture to the full text width, y running down
% from 0, so every figure shares one coordinate convention.  It *sets* the
% bounding box rather than contributing to it, so the emitted page is exactly
% the text width and every figure is included at 1:1 -- otherwise a stray
% invisible extent (pgfplots reserves a little around an axis) rescales one
% figure relative to the others and the same \scSmall prints at two sizes.
% Anything drawn outside the canvas is therefore cropped: keep content inside,
% and check with
%     \providecommand{\scCanvas}[1]{\path[draw=red] (0,0) rectangle (\scWnum,-#1);}
\providecommand{\scCanvas}[1]{%
  \path[use as bounding box] (0,0) rectangle (\scWnum,-#1);}

% --- palette --------------------------------------------------------------
% Canonical companion-paper values (../../../../Style/style.tex).
\definecolor{npgOrange}{HTML}{E69F00}
\definecolor{npgSky}{HTML}{56B4E9}
\definecolor{npgGreen}{HTML}{009E73}
\definecolor{npgYellow}{HTML}{F0E442}
\definecolor{npgBlue}{HTML}{0072B2}
\definecolor{npgVermillion}{HTML}{E64B35}
\definecolor{npgPurple}{HTML}{CC79A7}
\definecolor{figInk}{HTML}{1A1A1A}
\definecolor{figMuted}{HTML}{666666}
\definecolor{figHair}{HTML}{D9D9D9}
\definecolor{figWash}{HTML}{F7F7F7}

\colorlet{scInk}{figInk}              % primary text
\colorlet{scInkM}{figMuted}           % secondary text
\colorlet{scPaper}{figWash}           % quiet cards and note strips
\colorlet{scSpF}{npgSky!12}           % species fill
\colorlet{scSpE}{npgBlue}             % species edge
\colorlet{scRxF}{npgOrange!12}        % reaction fill
\colorlet{scRxE}{npgOrange!88!black}  % reaction edge
\colorlet{scReac}{npgBlue!78!black}   % reactant edge
\colorlet{scProd}{npgGreen!82!black}  % product edge
\colorlet{scHiF}{npgOrange!13}        % highlight fill
\colorlet{scHiE}{npgOrange!82!black}  % highlight edge
% Stage / accent tints.  stDE remains available for the fifth plot series.
\colorlet{stAF}{npgGreen!9}\colorlet{stAE}{npgGreen}
\colorlet{stBF}{npgOrange!11}\colorlet{stBE}{npgOrange}
\colorlet{stCF}{npgBlue!9}\colorlet{stCE}{npgBlue}
\colorlet{stDF}{npgPurple!10}\colorlet{stDE}{npgPurple!85!black}
\colorlet{stEF}{npgVermillion!9}\colorlet{stEE}{npgVermillion}

% Semantic aliases.  Figures should speak in these, not in stAE/stBE/...
\colorlet{scAccent}{stAE}\colorlet{scAccentF}{stAF}   % structure / normal
\colorlet{scMark}{stBE}\colorlet{scMarkF}{stBF}       % highlight / certificate
\colorlet{scRef}{stCE}\colorlet{scRefF}{stCF}         % secondary series
\colorlet{scAlert}{stEE}\colorlet{scAlertF}{stEF}     % failure / violation
\colorlet{scRule}{figHair}                            % card and axis hairlines
\colorlet{scRuleL}{figHair!55}                        % grid, dividers

% --- text styles ----------------------------------------------------------
% Three sizes, fixed in points so they do not drift with the document class.
% Arial is shared by labels, headings, panel badges, and chemistry labels.
\providecommand{\scTitle}{\scArial\fontsize{7}{8.4}\selectfont\bfseries\color{scInk}}
\providecommand{\scBody}{\scArial\fontsize{6.5}{7.8}\selectfont\color{scInkM}}
\providecommand{\scSmall}{\scArial\fontsize{6}{7.2}\selectfont\color{scInkM}}
\providecommand{\scSub}[1]{{\scSmall\itshape #1}}
\providecommand{\scPanel}[1]{{\scTitle(#1)}}
% Numerals that carry a result: ink, not colour, so the eye reads the value.
\providecommand{\scNum}{\scArial\fontsize{10}{11}\selectfont\bfseries\color{scInk}}

% --- shared TikZ styles ---------------------------------------------------
\tikzset{
  % text never hyphenates inside a figure
  every node/.append style={execute at begin node={%
      \hyphenpenalty=10000 \exhyphenpenalty=10000 \hbadness=10000 }},
  syn/.style={font=\scArial,color=scInk,
              execute at begin picture={\MTversion*{synarial}},
              >={Stealth[round,length=2.4mm,width=1.9mm]}},
  % containers ------------------------------------------------------------
  % One card chrome for the whole paper: white, 3 pt radius, one hairline.
  panel/.style={rounded corners=2.5pt,fill=none,draw=scRule,line width=0.45pt},
  card/.style={rounded corners=2.2pt,fill=scPaper,draw=scRule,line width=0.45pt,
               inner sep=5pt,align=left,anchor=north west},
  % a quiet strip for a panel's closing fact
  note/.style={rounded corners=2.5pt,fill=scPaper,draw=scRuleL,line width=0.4pt,
               inner xsep=6pt,inner ysep=3.5pt,align=left,font=\scSmall},
  % small labelled value pill
  chip/.style={rounded corners=2pt,fill=scPaper,draw=scRule,line width=0.4pt,
               inner xsep=4pt,inner ysep=2pt,font=\scSmall},
  % Circular letter badge: the same accessible panel marker used in Style/*.
  hbullet/.style={circle,fill=scInk,text=white,draw=none,
                  font=\scArial\fontsize{6}{6}\selectfont\bfseries,
                  minimum size=4.2mm,inner sep=0pt},
  ctag/.style={font=\scTitle,anchor=north west,text=scInk},
  % CRN primitives --------------------------------------------------------
  species/.style={rounded rectangle,fill=scSpF,draw=scSpE,line width=0.45pt,
                  inner xsep=2.4pt,inner ysep=1.6pt,font=\scSmall,
                  text=scInk,minimum height=3.6mm},
  cofactor/.style={species,fill=scRefF,draw=scRef},
  catalyst/.style={species,fill=scHiF,draw=scHiE},
  rxn/.style={rectangle,fill=scRxF,draw=scRxE,line width=0.45pt,
              inner sep=1.2pt,font=\scSmall,text=scInk,
              minimum width=3.0mm,minimum height=3.0mm},
  place/.style={circle,fill=scSpF,draw=scSpE,line width=0.45pt,
                minimum size=4.0mm,inner sep=0pt},
  trans/.style={rectangle,fill=scRxE,draw=scRxE,minimum width=1.1mm,
                minimum height=3.6mm,inner sep=0pt},
  token/.style={circle,fill=scInk,minimum size=1.1mm,inner sep=0pt},
  % edges -----------------------------------------------------------------
  reac/.style={->,line width=0.5pt,draw=scReac,shorten >=0.8pt,shorten <=0.8pt},
  prod/.style={->,line width=0.5pt,draw=scProd,shorten >=0.8pt,shorten <=0.8pt},
  cat/.style={->,line width=0.5pt,draw=scHiE,shorten >=0.8pt,shorten <=0.8pt},
  flow/.style={-{Stealth[round,length=2.2mm,width=1.8mm]},line width=1.1pt,
               draw=scInkM!45},
  flab/.style={font=\scSmall,text=scInkM,inner sep=2pt},
  stoi/.style={font=\scSmall,text=scInkM,inner sep=0.8pt},
  % verdict marks: shape carries the meaning, colour only reinforces it
  vpass/.style={circle,fill=scAccentF,draw=scAccent,text=scAccent,
                font=\fontsize{6}{6}\selectfont\bfseries,
                minimum size=3.2mm,inner sep=0pt},
  vfail/.style={circle,fill=scAlertF,draw=scAlert,text=scAlert,
                font=\fontsize{6}{6}\selectfont\bfseries,
                minimum size=3.2mm,inner sep=0pt},
  callout/.style={font=\scSmall,text=scInkM,align=left},
}

% \scHead[<badge colour>]{<anchor coord>}{<letter>}{<title>}{<subtitle or empty>}
% One header treatment for every panel in the paper.
\providecommand{\scHead}[5][scInk]{%
  \node[hbullet,fill=#1,anchor=north west] (scbadge#3) at ($#2+(0,-0.015)$)
    {\MakeUppercase{#3}};
  \node[ctag,anchor=west] at ($(scbadge#3.east)+(0.14,0)$) {#4};
  \ifx\relax#5\relax\else
    \node[anchor=north west,font=\scSmall,text=scInkM]
      at ($#2+(0.56,-0.245)$) {#5};
  \fi}

% --- shared pgfplots look (only if the figure loaded pgfplots) ------------
\ifdefined\pgfplotsset
  \pgfplotsset{compat=1.18}
  \pgfplotsset{
    synaxis/.style={
      scale only axis,
      tick align=outside, tick pos=left,
      axis line style={draw=scInkM!55,line width=0.4pt},
      % One grid, on the value axis only.  No minor grid: at these sizes it
      % turns the plot frame into a grey wash.
      ymajorgrids=true, xmajorgrids=false,
      major grid style={draw=scRuleL,line width=0.3pt},
      major tick length=1.6pt,
      label style={font=\scBody,color=scInk},
      tick label style={font=\scSmall,color=scInkM},
      legend style={font=\scSmall,draw=none,fill=none,inner sep=1pt,
                    row sep=0.4pt,legend columns=1},
      legend cell align=left,
      every axis plot/.append style={line width=0.7pt,mark size=1.2pt},
    },
  }
\fi
}{% ---------------------------------------------------------------------------
% synstyle.tex --- shared look for every TikZ figure of the SynKit CRN paper.
%
% Input it from a figure file:
%     \IfFileExists{synstyle.tex}{\input{synstyle}}{\input{Figure/synstyle}}
% so that the figure compiles both from Figure/ and from the paper root.
%
% The visual vocabulary follows ../../../../Style/style.tex: the
% colour-blind-safe Okabe--Ito accents, quiet neutral rules and wash, circular
% panel badges, and a compact technical type scale are shared by the companion
% papers.  Semantic aliases below keep the CRN meaning explicit.
%
% DESIGN CONTRACT -----------------------------------------------------------
% 1. Geometry.  Every figure is laid out on a canvas exactly \scW wide (the
%    sn-jnl text width, 372 pt) and is included unscaled, so one point of type
%    is the same physical size in all five figures.  Open each picture with
%    \scCanvas{<height in cm>} to pin the bounding box.
% 2. Colour is semantic, never decorative.  Four roles only:
%       scAccent (green)       structure, the normal/expected path
%       scMark   (orange)      catalyst, certificate, or highlight
%       scAlert  (vermillion)  failure, violation, or uncovered set
%       scRef    (blue)        secondary or reference series
%    A panel gets at most one accent.  Titles are always ink; panel identity
%    is also encoded by a letter inside a circular badge, not by colour alone.
% 3. Type scale.  Three sizes, no others, and nothing below 6 pt:
%       \scTitle (7 pt bold)  panel titles
%       \scBody  (6.5 pt)     labels, values, in-figure prose
%       \scSmall (6 pt)       units, footnotes, dense graph labels
% ---------------------------------------------------------------------------
% Load guard.  The paper inputs every figure source inline, so this file is
% reached once per figure; without the guard the second \newlength below
% would abort the run.
\ifdefined\SynKitStyleLoaded\endinput\fi
\newcommand{\SynKitStyleLoaded}{}

% Load in the document preamble. XeLaTeX/LuaLaTeX and an installed Arial
% family are required; fontspec reports an error instead of substituting.
\RequirePackage[no-math]{fontspec}
\newfontfamily\scArial{Arial}[NFSSFamily=scArial,Ligatures=TeX]
% Keep manuscript mathematics unchanged; select Arial letters and numerals
% only inside figures. TeX supplies specialist mathematical symbols.
\RequirePackage[subdued,italic]{mathastext}
\MTDeclareVersion{synarial}{TU}{scArial}{m}{n}
\SetMathAlphabet{\Mathsf}{synarial}{TU}{scArial}{m}{n}
\SetMathAlphabet{\Mathtt}{synarial}{TU}{scArial}{m}{n}

\usetikzlibrary{positioning,calc,fit,backgrounds,arrows.meta,shapes.geometric,
                shapes.misc,decorations.pathreplacing,patterns,math}

% --- canvas ---------------------------------------------------------------
% 372 pt text width minus the 1 pt standalone border on each side.
\newlength{\scW}\setlength{\scW}{13.0326cm}
\providecommand{\scWnum}{13.0326}
% \scCanvas{<height>} pins the picture to the full text width, y running down
% from 0, so every figure shares one coordinate convention.  It *sets* the
% bounding box rather than contributing to it, so the emitted page is exactly
% the text width and every figure is included at 1:1 -- otherwise a stray
% invisible extent (pgfplots reserves a little around an axis) rescales one
% figure relative to the others and the same \scSmall prints at two sizes.
% Anything drawn outside the canvas is therefore cropped: keep content inside,
% and check with
%     \providecommand{\scCanvas}[1]{\path[draw=red] (0,0) rectangle (\scWnum,-#1);}
\providecommand{\scCanvas}[1]{%
  \path[use as bounding box] (0,0) rectangle (\scWnum,-#1);}

% --- palette --------------------------------------------------------------
% Canonical companion-paper values (../../../../Style/style.tex).
\definecolor{npgOrange}{HTML}{E69F00}
\definecolor{npgSky}{HTML}{56B4E9}
\definecolor{npgGreen}{HTML}{009E73}
\definecolor{npgYellow}{HTML}{F0E442}
\definecolor{npgBlue}{HTML}{0072B2}
\definecolor{npgVermillion}{HTML}{E64B35}
\definecolor{npgPurple}{HTML}{CC79A7}
\definecolor{figInk}{HTML}{1A1A1A}
\definecolor{figMuted}{HTML}{666666}
\definecolor{figHair}{HTML}{D9D9D9}
\definecolor{figWash}{HTML}{F7F7F7}

\colorlet{scInk}{figInk}              % primary text
\colorlet{scInkM}{figMuted}           % secondary text
\colorlet{scPaper}{figWash}           % quiet cards and note strips
\colorlet{scSpF}{npgSky!12}           % species fill
\colorlet{scSpE}{npgBlue}             % species edge
\colorlet{scRxF}{npgOrange!12}        % reaction fill
\colorlet{scRxE}{npgOrange!88!black}  % reaction edge
\colorlet{scReac}{npgBlue!78!black}   % reactant edge
\colorlet{scProd}{npgGreen!82!black}  % product edge
\colorlet{scHiF}{npgOrange!13}        % highlight fill
\colorlet{scHiE}{npgOrange!82!black}  % highlight edge
% Stage / accent tints.  stDE remains available for the fifth plot series.
\colorlet{stAF}{npgGreen!9}\colorlet{stAE}{npgGreen}
\colorlet{stBF}{npgOrange!11}\colorlet{stBE}{npgOrange}
\colorlet{stCF}{npgBlue!9}\colorlet{stCE}{npgBlue}
\colorlet{stDF}{npgPurple!10}\colorlet{stDE}{npgPurple!85!black}
\colorlet{stEF}{npgVermillion!9}\colorlet{stEE}{npgVermillion}

% Semantic aliases.  Figures should speak in these, not in stAE/stBE/...
\colorlet{scAccent}{stAE}\colorlet{scAccentF}{stAF}   % structure / normal
\colorlet{scMark}{stBE}\colorlet{scMarkF}{stBF}       % highlight / certificate
\colorlet{scRef}{stCE}\colorlet{scRefF}{stCF}         % secondary series
\colorlet{scAlert}{stEE}\colorlet{scAlertF}{stEF}     % failure / violation
\colorlet{scRule}{figHair}                            % card and axis hairlines
\colorlet{scRuleL}{figHair!55}                        % grid, dividers

% --- text styles ----------------------------------------------------------
% Three sizes, fixed in points so they do not drift with the document class.
% Arial is shared by labels, headings, panel badges, and chemistry labels.
\providecommand{\scTitle}{\scArial\fontsize{7}{8.4}\selectfont\bfseries\color{scInk}}
\providecommand{\scBody}{\scArial\fontsize{6.5}{7.8}\selectfont\color{scInkM}}
\providecommand{\scSmall}{\scArial\fontsize{6}{7.2}\selectfont\color{scInkM}}
\providecommand{\scSub}[1]{{\scSmall\itshape #1}}
\providecommand{\scPanel}[1]{{\scTitle(#1)}}
% Numerals that carry a result: ink, not colour, so the eye reads the value.
\providecommand{\scNum}{\scArial\fontsize{10}{11}\selectfont\bfseries\color{scInk}}

% --- shared TikZ styles ---------------------------------------------------
\tikzset{
  % text never hyphenates inside a figure
  every node/.append style={execute at begin node={%
      \hyphenpenalty=10000 \exhyphenpenalty=10000 \hbadness=10000 }},
  syn/.style={font=\scArial,color=scInk,
              execute at begin picture={\MTversion*{synarial}},
              >={Stealth[round,length=2.4mm,width=1.9mm]}},
  % containers ------------------------------------------------------------
  % One card chrome for the whole paper: white, 3 pt radius, one hairline.
  panel/.style={rounded corners=2.5pt,fill=none,draw=scRule,line width=0.45pt},
  card/.style={rounded corners=2.2pt,fill=scPaper,draw=scRule,line width=0.45pt,
               inner sep=5pt,align=left,anchor=north west},
  % a quiet strip for a panel's closing fact
  note/.style={rounded corners=2.5pt,fill=scPaper,draw=scRuleL,line width=0.4pt,
               inner xsep=6pt,inner ysep=3.5pt,align=left,font=\scSmall},
  % small labelled value pill
  chip/.style={rounded corners=2pt,fill=scPaper,draw=scRule,line width=0.4pt,
               inner xsep=4pt,inner ysep=2pt,font=\scSmall},
  % Circular letter badge: the same accessible panel marker used in Style/*.
  hbullet/.style={circle,fill=scInk,text=white,draw=none,
                  font=\scArial\fontsize{6}{6}\selectfont\bfseries,
                  minimum size=4.2mm,inner sep=0pt},
  ctag/.style={font=\scTitle,anchor=north west,text=scInk},
  % CRN primitives --------------------------------------------------------
  species/.style={rounded rectangle,fill=scSpF,draw=scSpE,line width=0.45pt,
                  inner xsep=2.4pt,inner ysep=1.6pt,font=\scSmall,
                  text=scInk,minimum height=3.6mm},
  cofactor/.style={species,fill=scRefF,draw=scRef},
  catalyst/.style={species,fill=scHiF,draw=scHiE},
  rxn/.style={rectangle,fill=scRxF,draw=scRxE,line width=0.45pt,
              inner sep=1.2pt,font=\scSmall,text=scInk,
              minimum width=3.0mm,minimum height=3.0mm},
  place/.style={circle,fill=scSpF,draw=scSpE,line width=0.45pt,
                minimum size=4.0mm,inner sep=0pt},
  trans/.style={rectangle,fill=scRxE,draw=scRxE,minimum width=1.1mm,
                minimum height=3.6mm,inner sep=0pt},
  token/.style={circle,fill=scInk,minimum size=1.1mm,inner sep=0pt},
  % edges -----------------------------------------------------------------
  reac/.style={->,line width=0.5pt,draw=scReac,shorten >=0.8pt,shorten <=0.8pt},
  prod/.style={->,line width=0.5pt,draw=scProd,shorten >=0.8pt,shorten <=0.8pt},
  cat/.style={->,line width=0.5pt,draw=scHiE,shorten >=0.8pt,shorten <=0.8pt},
  flow/.style={-{Stealth[round,length=2.2mm,width=1.8mm]},line width=1.1pt,
               draw=scInkM!45},
  flab/.style={font=\scSmall,text=scInkM,inner sep=2pt},
  stoi/.style={font=\scSmall,text=scInkM,inner sep=0.8pt},
  % verdict marks: shape carries the meaning, colour only reinforces it
  vpass/.style={circle,fill=scAccentF,draw=scAccent,text=scAccent,
                font=\fontsize{6}{6}\selectfont\bfseries,
                minimum size=3.2mm,inner sep=0pt},
  vfail/.style={circle,fill=scAlertF,draw=scAlert,text=scAlert,
                font=\fontsize{6}{6}\selectfont\bfseries,
                minimum size=3.2mm,inner sep=0pt},
  callout/.style={font=\scSmall,text=scInkM,align=left},
}

% \scHead[<badge colour>]{<anchor coord>}{<letter>}{<title>}{<subtitle or empty>}
% One header treatment for every panel in the paper.
\providecommand{\scHead}[5][scInk]{%
  \node[hbullet,fill=#1,anchor=north west] (scbadge#3) at ($#2+(0,-0.015)$)
    {\MakeUppercase{#3}};
  \node[ctag,anchor=west] at ($(scbadge#3.east)+(0.14,0)$) {#4};
  \ifx\relax#5\relax\else
    \node[anchor=north west,font=\scSmall,text=scInkM]
      at ($#2+(0.56,-0.245)$) {#5};
  \fi}

% --- shared pgfplots look (only if the figure loaded pgfplots) ------------
\ifdefined\pgfplotsset
  \pgfplotsset{compat=1.18}
  \pgfplotsset{
    synaxis/.style={
      scale only axis,
      tick align=outside, tick pos=left,
      axis line style={draw=scInkM!55,line width=0.4pt},
      % One grid, on the value axis only.  No minor grid: at these sizes it
      % turns the plot frame into a grey wash.
      ymajorgrids=true, xmajorgrids=false,
      major grid style={draw=scRuleL,line width=0.3pt},
      major tick length=1.6pt,
      label style={font=\scBody,color=scInk},
      tick label style={font=\scSmall,color=scInkM},
      legend style={font=\scSmall,draw=none,fill=none,inner sep=1pt,
                    row sep=0.4pt,legend columns=1},
      legend cell align=left,
      every axis plot/.append style={line width=0.7pt,mark size=1.2pt},
    },
  }
\fi
}
% so that the figure compiles both from Figure/ and from the paper root.
%
% The visual vocabulary follows ../../../../Style/style.tex: the
% colour-blind-safe Okabe--Ito accents, quiet neutral rules and wash, circular
% panel badges, and a compact technical type scale are shared by the companion
% papers.  Semantic aliases below keep the CRN meaning explicit.
%
% DESIGN CONTRACT -----------------------------------------------------------
% 1. Geometry.  Every figure is laid out on a canvas exactly \scW wide (the
%    sn-jnl text width, 372 pt) and is included unscaled, so one point of type
%    is the same physical size in all five figures.  Open each picture with
%    \scCanvas{<height in cm>} to pin the bounding box.
% 2. Colour is semantic, never decorative.  Four roles only:
%       scAccent (green)       structure, the normal/expected path
%       scMark   (orange)      catalyst, certificate, or highlight
%       scAlert  (vermillion)  failure, violation, or uncovered set
%       scRef    (blue)        secondary or reference series
%    A panel gets at most one accent.  Titles are always ink; panel identity
%    is also encoded by a letter inside a circular badge, not by colour alone.
% 3. Type scale.  Three sizes, no others, and nothing below 6 pt:
%       \scTitle (7 pt bold)  panel titles
%       \scBody  (6.5 pt)     labels, values, in-figure prose
%       \scSmall (6 pt)       units, footnotes, dense graph labels
% ---------------------------------------------------------------------------
% Load guard.  The paper inputs every figure source inline, so this file is
% reached once per figure; without the guard the second \newlength below
% would abort the run.
\ifdefined\SynKitStyleLoaded\endinput\fi
\newcommand{\SynKitStyleLoaded}{}

% Load in the document preamble. XeLaTeX/LuaLaTeX and an installed Arial
% family are required; fontspec reports an error instead of substituting.
\RequirePackage[no-math]{fontspec}
\newfontfamily\scArial{Arial}[NFSSFamily=scArial,Ligatures=TeX]
% Keep manuscript mathematics unchanged; select Arial letters and numerals
% only inside figures. TeX supplies specialist mathematical symbols.
\RequirePackage[subdued,italic]{mathastext}
\MTDeclareVersion{synarial}{TU}{scArial}{m}{n}
\SetMathAlphabet{\Mathsf}{synarial}{TU}{scArial}{m}{n}
\SetMathAlphabet{\Mathtt}{synarial}{TU}{scArial}{m}{n}

\usetikzlibrary{positioning,calc,fit,backgrounds,arrows.meta,shapes.geometric,
                shapes.misc,decorations.pathreplacing,patterns,math}

% --- canvas ---------------------------------------------------------------
% 372 pt text width minus the 1 pt standalone border on each side.
\newlength{\scW}\setlength{\scW}{13.0326cm}
\providecommand{\scWnum}{13.0326}
% \scCanvas{<height>} pins the picture to the full text width, y running down
% from 0, so every figure shares one coordinate convention.  It *sets* the
% bounding box rather than contributing to it, so the emitted page is exactly
% the text width and every figure is included at 1:1 -- otherwise a stray
% invisible extent (pgfplots reserves a little around an axis) rescales one
% figure relative to the others and the same \scSmall prints at two sizes.
% Anything drawn outside the canvas is therefore cropped: keep content inside,
% and check with
%     \providecommand{\scCanvas}[1]{\path[draw=red] (0,0) rectangle (\scWnum,-#1);}
\providecommand{\scCanvas}[1]{%
  \path[use as bounding box] (0,0) rectangle (\scWnum,-#1);}

% --- palette --------------------------------------------------------------
% Canonical companion-paper values (../../../../Style/style.tex).
\definecolor{npgOrange}{HTML}{E69F00}
\definecolor{npgSky}{HTML}{56B4E9}
\definecolor{npgGreen}{HTML}{009E73}
\definecolor{npgYellow}{HTML}{F0E442}
\definecolor{npgBlue}{HTML}{0072B2}
\definecolor{npgVermillion}{HTML}{E64B35}
\definecolor{npgPurple}{HTML}{CC79A7}
\definecolor{figInk}{HTML}{1A1A1A}
\definecolor{figMuted}{HTML}{666666}
\definecolor{figHair}{HTML}{D9D9D9}
\definecolor{figWash}{HTML}{F7F7F7}

\colorlet{scInk}{figInk}              % primary text
\colorlet{scInkM}{figMuted}           % secondary text
\colorlet{scPaper}{figWash}           % quiet cards and note strips
\colorlet{scSpF}{npgSky!12}           % species fill
\colorlet{scSpE}{npgBlue}             % species edge
\colorlet{scRxF}{npgOrange!12}        % reaction fill
\colorlet{scRxE}{npgOrange!88!black}  % reaction edge
\colorlet{scReac}{npgBlue!78!black}   % reactant edge
\colorlet{scProd}{npgGreen!82!black}  % product edge
\colorlet{scHiF}{npgOrange!13}        % highlight fill
\colorlet{scHiE}{npgOrange!82!black}  % highlight edge
% Stage / accent tints.  stDE remains available for the fifth plot series.
\colorlet{stAF}{npgGreen!9}\colorlet{stAE}{npgGreen}
\colorlet{stBF}{npgOrange!11}\colorlet{stBE}{npgOrange}
\colorlet{stCF}{npgBlue!9}\colorlet{stCE}{npgBlue}
\colorlet{stDF}{npgPurple!10}\colorlet{stDE}{npgPurple!85!black}
\colorlet{stEF}{npgVermillion!9}\colorlet{stEE}{npgVermillion}

% Semantic aliases.  Figures should speak in these, not in stAE/stBE/...
\colorlet{scAccent}{stAE}\colorlet{scAccentF}{stAF}   % structure / normal
\colorlet{scMark}{stBE}\colorlet{scMarkF}{stBF}       % highlight / certificate
\colorlet{scRef}{stCE}\colorlet{scRefF}{stCF}         % secondary series
\colorlet{scAlert}{stEE}\colorlet{scAlertF}{stEF}     % failure / violation
\colorlet{scRule}{figHair}                            % card and axis hairlines
\colorlet{scRuleL}{figHair!55}                        % grid, dividers

% --- text styles ----------------------------------------------------------
% Three sizes, fixed in points so they do not drift with the document class.
% Arial is shared by labels, headings, panel badges, and chemistry labels.
\providecommand{\scTitle}{\scArial\fontsize{7}{8.4}\selectfont\bfseries\color{scInk}}
\providecommand{\scBody}{\scArial\fontsize{6.5}{7.8}\selectfont\color{scInkM}}
\providecommand{\scSmall}{\scArial\fontsize{6}{7.2}\selectfont\color{scInkM}}
\providecommand{\scSub}[1]{{\scSmall\itshape #1}}
\providecommand{\scPanel}[1]{{\scTitle(#1)}}
% Numerals that carry a result: ink, not colour, so the eye reads the value.
\providecommand{\scNum}{\scArial\fontsize{10}{11}\selectfont\bfseries\color{scInk}}

% --- shared TikZ styles ---------------------------------------------------
\tikzset{
  % text never hyphenates inside a figure
  every node/.append style={execute at begin node={%
      \hyphenpenalty=10000 \exhyphenpenalty=10000 \hbadness=10000 }},
  syn/.style={font=\scArial,color=scInk,
              execute at begin picture={\MTversion*{synarial}},
              >={Stealth[round,length=2.4mm,width=1.9mm]}},
  % containers ------------------------------------------------------------
  % One card chrome for the whole paper: white, 3 pt radius, one hairline.
  panel/.style={rounded corners=2.5pt,fill=none,draw=scRule,line width=0.45pt},
  card/.style={rounded corners=2.2pt,fill=scPaper,draw=scRule,line width=0.45pt,
               inner sep=5pt,align=left,anchor=north west},
  % a quiet strip for a panel's closing fact
  note/.style={rounded corners=2.5pt,fill=scPaper,draw=scRuleL,line width=0.4pt,
               inner xsep=6pt,inner ysep=3.5pt,align=left,font=\scSmall},
  % small labelled value pill
  chip/.style={rounded corners=2pt,fill=scPaper,draw=scRule,line width=0.4pt,
               inner xsep=4pt,inner ysep=2pt,font=\scSmall},
  % Circular letter badge: the same accessible panel marker used in Style/*.
  hbullet/.style={circle,fill=scInk,text=white,draw=none,
                  font=\scArial\fontsize{6}{6}\selectfont\bfseries,
                  minimum size=4.2mm,inner sep=0pt},
  ctag/.style={font=\scTitle,anchor=north west,text=scInk},
  % CRN primitives --------------------------------------------------------
  species/.style={rounded rectangle,fill=scSpF,draw=scSpE,line width=0.45pt,
                  inner xsep=2.4pt,inner ysep=1.6pt,font=\scSmall,
                  text=scInk,minimum height=3.6mm},
  cofactor/.style={species,fill=scRefF,draw=scRef},
  catalyst/.style={species,fill=scHiF,draw=scHiE},
  rxn/.style={rectangle,fill=scRxF,draw=scRxE,line width=0.45pt,
              inner sep=1.2pt,font=\scSmall,text=scInk,
              minimum width=3.0mm,minimum height=3.0mm},
  place/.style={circle,fill=scSpF,draw=scSpE,line width=0.45pt,
                minimum size=4.0mm,inner sep=0pt},
  trans/.style={rectangle,fill=scRxE,draw=scRxE,minimum width=1.1mm,
                minimum height=3.6mm,inner sep=0pt},
  token/.style={circle,fill=scInk,minimum size=1.1mm,inner sep=0pt},
  % edges -----------------------------------------------------------------
  reac/.style={->,line width=0.5pt,draw=scReac,shorten >=0.8pt,shorten <=0.8pt},
  prod/.style={->,line width=0.5pt,draw=scProd,shorten >=0.8pt,shorten <=0.8pt},
  cat/.style={->,line width=0.5pt,draw=scHiE,shorten >=0.8pt,shorten <=0.8pt},
  flow/.style={-{Stealth[round,length=2.2mm,width=1.8mm]},line width=1.1pt,
               draw=scInkM!45},
  flab/.style={font=\scSmall,text=scInkM,inner sep=2pt},
  stoi/.style={font=\scSmall,text=scInkM,inner sep=0.8pt},
  % verdict marks: shape carries the meaning, colour only reinforces it
  vpass/.style={circle,fill=scAccentF,draw=scAccent,text=scAccent,
                font=\fontsize{6}{6}\selectfont\bfseries,
                minimum size=3.2mm,inner sep=0pt},
  vfail/.style={circle,fill=scAlertF,draw=scAlert,text=scAlert,
                font=\fontsize{6}{6}\selectfont\bfseries,
                minimum size=3.2mm,inner sep=0pt},
  callout/.style={font=\scSmall,text=scInkM,align=left},
}

% \scHead[<badge colour>]{<anchor coord>}{<letter>}{<title>}{<subtitle or empty>}
% One header treatment for every panel in the paper.
\providecommand{\scHead}[5][scInk]{%
  \node[hbullet,fill=#1,anchor=north west] (scbadge#3) at ($#2+(0,-0.015)$)
    {\MakeUppercase{#3}};
  \node[ctag,anchor=west] at ($(scbadge#3.east)+(0.14,0)$) {#4};
  \ifx\relax#5\relax\else
    \node[anchor=north west,font=\scSmall,text=scInkM]
      at ($#2+(0.56,-0.245)$) {#5};
  \fi}

% --- shared pgfplots look (only if the figure loaded pgfplots) ------------
\ifdefined\pgfplotsset
  \pgfplotsset{compat=1.18}
  \pgfplotsset{
    synaxis/.style={
      scale only axis,
      tick align=outside, tick pos=left,
      axis line style={draw=scInkM!55,line width=0.4pt},
      % One grid, on the value axis only.  No minor grid: at these sizes it
      % turns the plot frame into a grey wash.
      ymajorgrids=true, xmajorgrids=false,
      major grid style={draw=scRuleL,line width=0.3pt},
      major tick length=1.6pt,
      label style={font=\scBody,color=scInk},
      tick label style={font=\scSmall,color=scInkM},
      legend style={font=\scSmall,draw=none,fill=none,inner sep=1pt,
                    row sep=0.4pt,legend columns=1},
      legend cell align=left,
      every axis plot/.append style={line width=0.7pt,mark size=1.2pt},
    },
  }
\fi
}{% ---------------------------------------------------------------------------
% synstyle.tex --- shared look for every TikZ figure of the SynKit CRN paper.
%
% Input it from a figure file:
%     \IfFileExists{synstyle.tex}{% ---------------------------------------------------------------------------
% synstyle.tex --- shared look for every TikZ figure of the SynKit CRN paper.
%
% Input it from a figure file:
%     \IfFileExists{synstyle.tex}{\input{synstyle}}{\input{Figure/synstyle}}
% so that the figure compiles both from Figure/ and from the paper root.
%
% The visual vocabulary follows ../../../../Style/style.tex: the
% colour-blind-safe Okabe--Ito accents, quiet neutral rules and wash, circular
% panel badges, and a compact technical type scale are shared by the companion
% papers.  Semantic aliases below keep the CRN meaning explicit.
%
% DESIGN CONTRACT -----------------------------------------------------------
% 1. Geometry.  Every figure is laid out on a canvas exactly \scW wide (the
%    sn-jnl text width, 372 pt) and is included unscaled, so one point of type
%    is the same physical size in all five figures.  Open each picture with
%    \scCanvas{<height in cm>} to pin the bounding box.
% 2. Colour is semantic, never decorative.  Four roles only:
%       scAccent (green)       structure, the normal/expected path
%       scMark   (orange)      catalyst, certificate, or highlight
%       scAlert  (vermillion)  failure, violation, or uncovered set
%       scRef    (blue)        secondary or reference series
%    A panel gets at most one accent.  Titles are always ink; panel identity
%    is also encoded by a letter inside a circular badge, not by colour alone.
% 3. Type scale.  Three sizes, no others, and nothing below 6 pt:
%       \scTitle (7 pt bold)  panel titles
%       \scBody  (6.5 pt)     labels, values, in-figure prose
%       \scSmall (6 pt)       units, footnotes, dense graph labels
% ---------------------------------------------------------------------------
% Load guard.  The paper inputs every figure source inline, so this file is
% reached once per figure; without the guard the second \newlength below
% would abort the run.
\ifdefined\SynKitStyleLoaded\endinput\fi
\newcommand{\SynKitStyleLoaded}{}

% Load in the document preamble. XeLaTeX/LuaLaTeX and an installed Arial
% family are required; fontspec reports an error instead of substituting.
\RequirePackage[no-math]{fontspec}
\newfontfamily\scArial{Arial}[NFSSFamily=scArial,Ligatures=TeX]
% Keep manuscript mathematics unchanged; select Arial letters and numerals
% only inside figures. TeX supplies specialist mathematical symbols.
\RequirePackage[subdued,italic]{mathastext}
\MTDeclareVersion{synarial}{TU}{scArial}{m}{n}
\SetMathAlphabet{\Mathsf}{synarial}{TU}{scArial}{m}{n}
\SetMathAlphabet{\Mathtt}{synarial}{TU}{scArial}{m}{n}

\usetikzlibrary{positioning,calc,fit,backgrounds,arrows.meta,shapes.geometric,
                shapes.misc,decorations.pathreplacing,patterns,math}

% --- canvas ---------------------------------------------------------------
% 372 pt text width minus the 1 pt standalone border on each side.
\newlength{\scW}\setlength{\scW}{13.0326cm}
\providecommand{\scWnum}{13.0326}
% \scCanvas{<height>} pins the picture to the full text width, y running down
% from 0, so every figure shares one coordinate convention.  It *sets* the
% bounding box rather than contributing to it, so the emitted page is exactly
% the text width and every figure is included at 1:1 -- otherwise a stray
% invisible extent (pgfplots reserves a little around an axis) rescales one
% figure relative to the others and the same \scSmall prints at two sizes.
% Anything drawn outside the canvas is therefore cropped: keep content inside,
% and check with
%     \providecommand{\scCanvas}[1]{\path[draw=red] (0,0) rectangle (\scWnum,-#1);}
\providecommand{\scCanvas}[1]{%
  \path[use as bounding box] (0,0) rectangle (\scWnum,-#1);}

% --- palette --------------------------------------------------------------
% Canonical companion-paper values (../../../../Style/style.tex).
\definecolor{npgOrange}{HTML}{E69F00}
\definecolor{npgSky}{HTML}{56B4E9}
\definecolor{npgGreen}{HTML}{009E73}
\definecolor{npgYellow}{HTML}{F0E442}
\definecolor{npgBlue}{HTML}{0072B2}
\definecolor{npgVermillion}{HTML}{E64B35}
\definecolor{npgPurple}{HTML}{CC79A7}
\definecolor{figInk}{HTML}{1A1A1A}
\definecolor{figMuted}{HTML}{666666}
\definecolor{figHair}{HTML}{D9D9D9}
\definecolor{figWash}{HTML}{F7F7F7}

\colorlet{scInk}{figInk}              % primary text
\colorlet{scInkM}{figMuted}           % secondary text
\colorlet{scPaper}{figWash}           % quiet cards and note strips
\colorlet{scSpF}{npgSky!12}           % species fill
\colorlet{scSpE}{npgBlue}             % species edge
\colorlet{scRxF}{npgOrange!12}        % reaction fill
\colorlet{scRxE}{npgOrange!88!black}  % reaction edge
\colorlet{scReac}{npgBlue!78!black}   % reactant edge
\colorlet{scProd}{npgGreen!82!black}  % product edge
\colorlet{scHiF}{npgOrange!13}        % highlight fill
\colorlet{scHiE}{npgOrange!82!black}  % highlight edge
% Stage / accent tints.  stDE remains available for the fifth plot series.
\colorlet{stAF}{npgGreen!9}\colorlet{stAE}{npgGreen}
\colorlet{stBF}{npgOrange!11}\colorlet{stBE}{npgOrange}
\colorlet{stCF}{npgBlue!9}\colorlet{stCE}{npgBlue}
\colorlet{stDF}{npgPurple!10}\colorlet{stDE}{npgPurple!85!black}
\colorlet{stEF}{npgVermillion!9}\colorlet{stEE}{npgVermillion}

% Semantic aliases.  Figures should speak in these, not in stAE/stBE/...
\colorlet{scAccent}{stAE}\colorlet{scAccentF}{stAF}   % structure / normal
\colorlet{scMark}{stBE}\colorlet{scMarkF}{stBF}       % highlight / certificate
\colorlet{scRef}{stCE}\colorlet{scRefF}{stCF}         % secondary series
\colorlet{scAlert}{stEE}\colorlet{scAlertF}{stEF}     % failure / violation
\colorlet{scRule}{figHair}                            % card and axis hairlines
\colorlet{scRuleL}{figHair!55}                        % grid, dividers

% --- text styles ----------------------------------------------------------
% Three sizes, fixed in points so they do not drift with the document class.
% Arial is shared by labels, headings, panel badges, and chemistry labels.
\providecommand{\scTitle}{\scArial\fontsize{7}{8.4}\selectfont\bfseries\color{scInk}}
\providecommand{\scBody}{\scArial\fontsize{6.5}{7.8}\selectfont\color{scInkM}}
\providecommand{\scSmall}{\scArial\fontsize{6}{7.2}\selectfont\color{scInkM}}
\providecommand{\scSub}[1]{{\scSmall\itshape #1}}
\providecommand{\scPanel}[1]{{\scTitle(#1)}}
% Numerals that carry a result: ink, not colour, so the eye reads the value.
\providecommand{\scNum}{\scArial\fontsize{10}{11}\selectfont\bfseries\color{scInk}}

% --- shared TikZ styles ---------------------------------------------------
\tikzset{
  % text never hyphenates inside a figure
  every node/.append style={execute at begin node={%
      \hyphenpenalty=10000 \exhyphenpenalty=10000 \hbadness=10000 }},
  syn/.style={font=\scArial,color=scInk,
              execute at begin picture={\MTversion*{synarial}},
              >={Stealth[round,length=2.4mm,width=1.9mm]}},
  % containers ------------------------------------------------------------
  % One card chrome for the whole paper: white, 3 pt radius, one hairline.
  panel/.style={rounded corners=2.5pt,fill=none,draw=scRule,line width=0.45pt},
  card/.style={rounded corners=2.2pt,fill=scPaper,draw=scRule,line width=0.45pt,
               inner sep=5pt,align=left,anchor=north west},
  % a quiet strip for a panel's closing fact
  note/.style={rounded corners=2.5pt,fill=scPaper,draw=scRuleL,line width=0.4pt,
               inner xsep=6pt,inner ysep=3.5pt,align=left,font=\scSmall},
  % small labelled value pill
  chip/.style={rounded corners=2pt,fill=scPaper,draw=scRule,line width=0.4pt,
               inner xsep=4pt,inner ysep=2pt,font=\scSmall},
  % Circular letter badge: the same accessible panel marker used in Style/*.
  hbullet/.style={circle,fill=scInk,text=white,draw=none,
                  font=\scArial\fontsize{6}{6}\selectfont\bfseries,
                  minimum size=4.2mm,inner sep=0pt},
  ctag/.style={font=\scTitle,anchor=north west,text=scInk},
  % CRN primitives --------------------------------------------------------
  species/.style={rounded rectangle,fill=scSpF,draw=scSpE,line width=0.45pt,
                  inner xsep=2.4pt,inner ysep=1.6pt,font=\scSmall,
                  text=scInk,minimum height=3.6mm},
  cofactor/.style={species,fill=scRefF,draw=scRef},
  catalyst/.style={species,fill=scHiF,draw=scHiE},
  rxn/.style={rectangle,fill=scRxF,draw=scRxE,line width=0.45pt,
              inner sep=1.2pt,font=\scSmall,text=scInk,
              minimum width=3.0mm,minimum height=3.0mm},
  place/.style={circle,fill=scSpF,draw=scSpE,line width=0.45pt,
                minimum size=4.0mm,inner sep=0pt},
  trans/.style={rectangle,fill=scRxE,draw=scRxE,minimum width=1.1mm,
                minimum height=3.6mm,inner sep=0pt},
  token/.style={circle,fill=scInk,minimum size=1.1mm,inner sep=0pt},
  % edges -----------------------------------------------------------------
  reac/.style={->,line width=0.5pt,draw=scReac,shorten >=0.8pt,shorten <=0.8pt},
  prod/.style={->,line width=0.5pt,draw=scProd,shorten >=0.8pt,shorten <=0.8pt},
  cat/.style={->,line width=0.5pt,draw=scHiE,shorten >=0.8pt,shorten <=0.8pt},
  flow/.style={-{Stealth[round,length=2.2mm,width=1.8mm]},line width=1.1pt,
               draw=scInkM!45},
  flab/.style={font=\scSmall,text=scInkM,inner sep=2pt},
  stoi/.style={font=\scSmall,text=scInkM,inner sep=0.8pt},
  % verdict marks: shape carries the meaning, colour only reinforces it
  vpass/.style={circle,fill=scAccentF,draw=scAccent,text=scAccent,
                font=\fontsize{6}{6}\selectfont\bfseries,
                minimum size=3.2mm,inner sep=0pt},
  vfail/.style={circle,fill=scAlertF,draw=scAlert,text=scAlert,
                font=\fontsize{6}{6}\selectfont\bfseries,
                minimum size=3.2mm,inner sep=0pt},
  callout/.style={font=\scSmall,text=scInkM,align=left},
}

% \scHead[<badge colour>]{<anchor coord>}{<letter>}{<title>}{<subtitle or empty>}
% One header treatment for every panel in the paper.
\providecommand{\scHead}[5][scInk]{%
  \node[hbullet,fill=#1,anchor=north west] (scbadge#3) at ($#2+(0,-0.015)$)
    {\MakeUppercase{#3}};
  \node[ctag,anchor=west] at ($(scbadge#3.east)+(0.14,0)$) {#4};
  \ifx\relax#5\relax\else
    \node[anchor=north west,font=\scSmall,text=scInkM]
      at ($#2+(0.56,-0.245)$) {#5};
  \fi}

% --- shared pgfplots look (only if the figure loaded pgfplots) ------------
\ifdefined\pgfplotsset
  \pgfplotsset{compat=1.18}
  \pgfplotsset{
    synaxis/.style={
      scale only axis,
      tick align=outside, tick pos=left,
      axis line style={draw=scInkM!55,line width=0.4pt},
      % One grid, on the value axis only.  No minor grid: at these sizes it
      % turns the plot frame into a grey wash.
      ymajorgrids=true, xmajorgrids=false,
      major grid style={draw=scRuleL,line width=0.3pt},
      major tick length=1.6pt,
      label style={font=\scBody,color=scInk},
      tick label style={font=\scSmall,color=scInkM},
      legend style={font=\scSmall,draw=none,fill=none,inner sep=1pt,
                    row sep=0.4pt,legend columns=1},
      legend cell align=left,
      every axis plot/.append style={line width=0.7pt,mark size=1.2pt},
    },
  }
\fi
}{% ---------------------------------------------------------------------------
% synstyle.tex --- shared look for every TikZ figure of the SynKit CRN paper.
%
% Input it from a figure file:
%     \IfFileExists{synstyle.tex}{\input{synstyle}}{\input{Figure/synstyle}}
% so that the figure compiles both from Figure/ and from the paper root.
%
% The visual vocabulary follows ../../../../Style/style.tex: the
% colour-blind-safe Okabe--Ito accents, quiet neutral rules and wash, circular
% panel badges, and a compact technical type scale are shared by the companion
% papers.  Semantic aliases below keep the CRN meaning explicit.
%
% DESIGN CONTRACT -----------------------------------------------------------
% 1. Geometry.  Every figure is laid out on a canvas exactly \scW wide (the
%    sn-jnl text width, 372 pt) and is included unscaled, so one point of type
%    is the same physical size in all five figures.  Open each picture with
%    \scCanvas{<height in cm>} to pin the bounding box.
% 2. Colour is semantic, never decorative.  Four roles only:
%       scAccent (green)       structure, the normal/expected path
%       scMark   (orange)      catalyst, certificate, or highlight
%       scAlert  (vermillion)  failure, violation, or uncovered set
%       scRef    (blue)        secondary or reference series
%    A panel gets at most one accent.  Titles are always ink; panel identity
%    is also encoded by a letter inside a circular badge, not by colour alone.
% 3. Type scale.  Three sizes, no others, and nothing below 6 pt:
%       \scTitle (7 pt bold)  panel titles
%       \scBody  (6.5 pt)     labels, values, in-figure prose
%       \scSmall (6 pt)       units, footnotes, dense graph labels
% ---------------------------------------------------------------------------
% Load guard.  The paper inputs every figure source inline, so this file is
% reached once per figure; without the guard the second \newlength below
% would abort the run.
\ifdefined\SynKitStyleLoaded\endinput\fi
\newcommand{\SynKitStyleLoaded}{}

% Load in the document preamble. XeLaTeX/LuaLaTeX and an installed Arial
% family are required; fontspec reports an error instead of substituting.
\RequirePackage[no-math]{fontspec}
\newfontfamily\scArial{Arial}[NFSSFamily=scArial,Ligatures=TeX]
% Keep manuscript mathematics unchanged; select Arial letters and numerals
% only inside figures. TeX supplies specialist mathematical symbols.
\RequirePackage[subdued,italic]{mathastext}
\MTDeclareVersion{synarial}{TU}{scArial}{m}{n}
\SetMathAlphabet{\Mathsf}{synarial}{TU}{scArial}{m}{n}
\SetMathAlphabet{\Mathtt}{synarial}{TU}{scArial}{m}{n}

\usetikzlibrary{positioning,calc,fit,backgrounds,arrows.meta,shapes.geometric,
                shapes.misc,decorations.pathreplacing,patterns,math}

% --- canvas ---------------------------------------------------------------
% 372 pt text width minus the 1 pt standalone border on each side.
\newlength{\scW}\setlength{\scW}{13.0326cm}
\providecommand{\scWnum}{13.0326}
% \scCanvas{<height>} pins the picture to the full text width, y running down
% from 0, so every figure shares one coordinate convention.  It *sets* the
% bounding box rather than contributing to it, so the emitted page is exactly
% the text width and every figure is included at 1:1 -- otherwise a stray
% invisible extent (pgfplots reserves a little around an axis) rescales one
% figure relative to the others and the same \scSmall prints at two sizes.
% Anything drawn outside the canvas is therefore cropped: keep content inside,
% and check with
%     \providecommand{\scCanvas}[1]{\path[draw=red] (0,0) rectangle (\scWnum,-#1);}
\providecommand{\scCanvas}[1]{%
  \path[use as bounding box] (0,0) rectangle (\scWnum,-#1);}

% --- palette --------------------------------------------------------------
% Canonical companion-paper values (../../../../Style/style.tex).
\definecolor{npgOrange}{HTML}{E69F00}
\definecolor{npgSky}{HTML}{56B4E9}
\definecolor{npgGreen}{HTML}{009E73}
\definecolor{npgYellow}{HTML}{F0E442}
\definecolor{npgBlue}{HTML}{0072B2}
\definecolor{npgVermillion}{HTML}{E64B35}
\definecolor{npgPurple}{HTML}{CC79A7}
\definecolor{figInk}{HTML}{1A1A1A}
\definecolor{figMuted}{HTML}{666666}
\definecolor{figHair}{HTML}{D9D9D9}
\definecolor{figWash}{HTML}{F7F7F7}

\colorlet{scInk}{figInk}              % primary text
\colorlet{scInkM}{figMuted}           % secondary text
\colorlet{scPaper}{figWash}           % quiet cards and note strips
\colorlet{scSpF}{npgSky!12}           % species fill
\colorlet{scSpE}{npgBlue}             % species edge
\colorlet{scRxF}{npgOrange!12}        % reaction fill
\colorlet{scRxE}{npgOrange!88!black}  % reaction edge
\colorlet{scReac}{npgBlue!78!black}   % reactant edge
\colorlet{scProd}{npgGreen!82!black}  % product edge
\colorlet{scHiF}{npgOrange!13}        % highlight fill
\colorlet{scHiE}{npgOrange!82!black}  % highlight edge
% Stage / accent tints.  stDE remains available for the fifth plot series.
\colorlet{stAF}{npgGreen!9}\colorlet{stAE}{npgGreen}
\colorlet{stBF}{npgOrange!11}\colorlet{stBE}{npgOrange}
\colorlet{stCF}{npgBlue!9}\colorlet{stCE}{npgBlue}
\colorlet{stDF}{npgPurple!10}\colorlet{stDE}{npgPurple!85!black}
\colorlet{stEF}{npgVermillion!9}\colorlet{stEE}{npgVermillion}

% Semantic aliases.  Figures should speak in these, not in stAE/stBE/...
\colorlet{scAccent}{stAE}\colorlet{scAccentF}{stAF}   % structure / normal
\colorlet{scMark}{stBE}\colorlet{scMarkF}{stBF}       % highlight / certificate
\colorlet{scRef}{stCE}\colorlet{scRefF}{stCF}         % secondary series
\colorlet{scAlert}{stEE}\colorlet{scAlertF}{stEF}     % failure / violation
\colorlet{scRule}{figHair}                            % card and axis hairlines
\colorlet{scRuleL}{figHair!55}                        % grid, dividers

% --- text styles ----------------------------------------------------------
% Three sizes, fixed in points so they do not drift with the document class.
% Arial is shared by labels, headings, panel badges, and chemistry labels.
\providecommand{\scTitle}{\scArial\fontsize{7}{8.4}\selectfont\bfseries\color{scInk}}
\providecommand{\scBody}{\scArial\fontsize{6.5}{7.8}\selectfont\color{scInkM}}
\providecommand{\scSmall}{\scArial\fontsize{6}{7.2}\selectfont\color{scInkM}}
\providecommand{\scSub}[1]{{\scSmall\itshape #1}}
\providecommand{\scPanel}[1]{{\scTitle(#1)}}
% Numerals that carry a result: ink, not colour, so the eye reads the value.
\providecommand{\scNum}{\scArial\fontsize{10}{11}\selectfont\bfseries\color{scInk}}

% --- shared TikZ styles ---------------------------------------------------
\tikzset{
  % text never hyphenates inside a figure
  every node/.append style={execute at begin node={%
      \hyphenpenalty=10000 \exhyphenpenalty=10000 \hbadness=10000 }},
  syn/.style={font=\scArial,color=scInk,
              execute at begin picture={\MTversion*{synarial}},
              >={Stealth[round,length=2.4mm,width=1.9mm]}},
  % containers ------------------------------------------------------------
  % One card chrome for the whole paper: white, 3 pt radius, one hairline.
  panel/.style={rounded corners=2.5pt,fill=none,draw=scRule,line width=0.45pt},
  card/.style={rounded corners=2.2pt,fill=scPaper,draw=scRule,line width=0.45pt,
               inner sep=5pt,align=left,anchor=north west},
  % a quiet strip for a panel's closing fact
  note/.style={rounded corners=2.5pt,fill=scPaper,draw=scRuleL,line width=0.4pt,
               inner xsep=6pt,inner ysep=3.5pt,align=left,font=\scSmall},
  % small labelled value pill
  chip/.style={rounded corners=2pt,fill=scPaper,draw=scRule,line width=0.4pt,
               inner xsep=4pt,inner ysep=2pt,font=\scSmall},
  % Circular letter badge: the same accessible panel marker used in Style/*.
  hbullet/.style={circle,fill=scInk,text=white,draw=none,
                  font=\scArial\fontsize{6}{6}\selectfont\bfseries,
                  minimum size=4.2mm,inner sep=0pt},
  ctag/.style={font=\scTitle,anchor=north west,text=scInk},
  % CRN primitives --------------------------------------------------------
  species/.style={rounded rectangle,fill=scSpF,draw=scSpE,line width=0.45pt,
                  inner xsep=2.4pt,inner ysep=1.6pt,font=\scSmall,
                  text=scInk,minimum height=3.6mm},
  cofactor/.style={species,fill=scRefF,draw=scRef},
  catalyst/.style={species,fill=scHiF,draw=scHiE},
  rxn/.style={rectangle,fill=scRxF,draw=scRxE,line width=0.45pt,
              inner sep=1.2pt,font=\scSmall,text=scInk,
              minimum width=3.0mm,minimum height=3.0mm},
  place/.style={circle,fill=scSpF,draw=scSpE,line width=0.45pt,
                minimum size=4.0mm,inner sep=0pt},
  trans/.style={rectangle,fill=scRxE,draw=scRxE,minimum width=1.1mm,
                minimum height=3.6mm,inner sep=0pt},
  token/.style={circle,fill=scInk,minimum size=1.1mm,inner sep=0pt},
  % edges -----------------------------------------------------------------
  reac/.style={->,line width=0.5pt,draw=scReac,shorten >=0.8pt,shorten <=0.8pt},
  prod/.style={->,line width=0.5pt,draw=scProd,shorten >=0.8pt,shorten <=0.8pt},
  cat/.style={->,line width=0.5pt,draw=scHiE,shorten >=0.8pt,shorten <=0.8pt},
  flow/.style={-{Stealth[round,length=2.2mm,width=1.8mm]},line width=1.1pt,
               draw=scInkM!45},
  flab/.style={font=\scSmall,text=scInkM,inner sep=2pt},
  stoi/.style={font=\scSmall,text=scInkM,inner sep=0.8pt},
  % verdict marks: shape carries the meaning, colour only reinforces it
  vpass/.style={circle,fill=scAccentF,draw=scAccent,text=scAccent,
                font=\fontsize{6}{6}\selectfont\bfseries,
                minimum size=3.2mm,inner sep=0pt},
  vfail/.style={circle,fill=scAlertF,draw=scAlert,text=scAlert,
                font=\fontsize{6}{6}\selectfont\bfseries,
                minimum size=3.2mm,inner sep=0pt},
  callout/.style={font=\scSmall,text=scInkM,align=left},
}

% \scHead[<badge colour>]{<anchor coord>}{<letter>}{<title>}{<subtitle or empty>}
% One header treatment for every panel in the paper.
\providecommand{\scHead}[5][scInk]{%
  \node[hbullet,fill=#1,anchor=north west] (scbadge#3) at ($#2+(0,-0.015)$)
    {\MakeUppercase{#3}};
  \node[ctag,anchor=west] at ($(scbadge#3.east)+(0.14,0)$) {#4};
  \ifx\relax#5\relax\else
    \node[anchor=north west,font=\scSmall,text=scInkM]
      at ($#2+(0.56,-0.245)$) {#5};
  \fi}

% --- shared pgfplots look (only if the figure loaded pgfplots) ------------
\ifdefined\pgfplotsset
  \pgfplotsset{compat=1.18}
  \pgfplotsset{
    synaxis/.style={
      scale only axis,
      tick align=outside, tick pos=left,
      axis line style={draw=scInkM!55,line width=0.4pt},
      % One grid, on the value axis only.  No minor grid: at these sizes it
      % turns the plot frame into a grey wash.
      ymajorgrids=true, xmajorgrids=false,
      major grid style={draw=scRuleL,line width=0.3pt},
      major tick length=1.6pt,
      label style={font=\scBody,color=scInk},
      tick label style={font=\scSmall,color=scInkM},
      legend style={font=\scSmall,draw=none,fill=none,inner sep=1pt,
                    row sep=0.4pt,legend columns=1},
      legend cell align=left,
      every axis plot/.append style={line width=0.7pt,mark size=1.2pt},
    },
  }
\fi
}
% so that the figure compiles both from Figure/ and from the paper root.
%
% The visual vocabulary follows ../../../../Style/style.tex: the
% colour-blind-safe Okabe--Ito accents, quiet neutral rules and wash, circular
% panel badges, and a compact technical type scale are shared by the companion
% papers.  Semantic aliases below keep the CRN meaning explicit.
%
% DESIGN CONTRACT -----------------------------------------------------------
% 1. Geometry.  Every figure is laid out on a canvas exactly \scW wide (the
%    sn-jnl text width, 372 pt) and is included unscaled, so one point of type
%    is the same physical size in all five figures.  Open each picture with
%    \scCanvas{<height in cm>} to pin the bounding box.
% 2. Colour is semantic, never decorative.  Four roles only:
%       scAccent (green)       structure, the normal/expected path
%       scMark   (orange)      catalyst, certificate, or highlight
%       scAlert  (vermillion)  failure, violation, or uncovered set
%       scRef    (blue)        secondary or reference series
%    A panel gets at most one accent.  Titles are always ink; panel identity
%    is also encoded by a letter inside a circular badge, not by colour alone.
% 3. Type scale.  Three sizes, no others, and nothing below 6 pt:
%       \scTitle (7 pt bold)  panel titles
%       \scBody  (6.5 pt)     labels, values, in-figure prose
%       \scSmall (6 pt)       units, footnotes, dense graph labels
% ---------------------------------------------------------------------------
% Load guard.  The paper inputs every figure source inline, so this file is
% reached once per figure; without the guard the second \newlength below
% would abort the run.
\ifdefined\SynKitStyleLoaded\endinput\fi
\newcommand{\SynKitStyleLoaded}{}

% Load in the document preamble. XeLaTeX/LuaLaTeX and an installed Arial
% family are required; fontspec reports an error instead of substituting.
\RequirePackage[no-math]{fontspec}
\newfontfamily\scArial{Arial}[NFSSFamily=scArial,Ligatures=TeX]
% Keep manuscript mathematics unchanged; select Arial letters and numerals
% only inside figures. TeX supplies specialist mathematical symbols.
\RequirePackage[subdued,italic]{mathastext}
\MTDeclareVersion{synarial}{TU}{scArial}{m}{n}
\SetMathAlphabet{\Mathsf}{synarial}{TU}{scArial}{m}{n}
\SetMathAlphabet{\Mathtt}{synarial}{TU}{scArial}{m}{n}

\usetikzlibrary{positioning,calc,fit,backgrounds,arrows.meta,shapes.geometric,
                shapes.misc,decorations.pathreplacing,patterns,math}

% --- canvas ---------------------------------------------------------------
% 372 pt text width minus the 1 pt standalone border on each side.
\newlength{\scW}\setlength{\scW}{13.0326cm}
\providecommand{\scWnum}{13.0326}
% \scCanvas{<height>} pins the picture to the full text width, y running down
% from 0, so every figure shares one coordinate convention.  It *sets* the
% bounding box rather than contributing to it, so the emitted page is exactly
% the text width and every figure is included at 1:1 -- otherwise a stray
% invisible extent (pgfplots reserves a little around an axis) rescales one
% figure relative to the others and the same \scSmall prints at two sizes.
% Anything drawn outside the canvas is therefore cropped: keep content inside,
% and check with
%     \providecommand{\scCanvas}[1]{\path[draw=red] (0,0) rectangle (\scWnum,-#1);}
\providecommand{\scCanvas}[1]{%
  \path[use as bounding box] (0,0) rectangle (\scWnum,-#1);}

% --- palette --------------------------------------------------------------
% Canonical companion-paper values (../../../../Style/style.tex).
\definecolor{npgOrange}{HTML}{E69F00}
\definecolor{npgSky}{HTML}{56B4E9}
\definecolor{npgGreen}{HTML}{009E73}
\definecolor{npgYellow}{HTML}{F0E442}
\definecolor{npgBlue}{HTML}{0072B2}
\definecolor{npgVermillion}{HTML}{E64B35}
\definecolor{npgPurple}{HTML}{CC79A7}
\definecolor{figInk}{HTML}{1A1A1A}
\definecolor{figMuted}{HTML}{666666}
\definecolor{figHair}{HTML}{D9D9D9}
\definecolor{figWash}{HTML}{F7F7F7}

\colorlet{scInk}{figInk}              % primary text
\colorlet{scInkM}{figMuted}           % secondary text
\colorlet{scPaper}{figWash}           % quiet cards and note strips
\colorlet{scSpF}{npgSky!12}           % species fill
\colorlet{scSpE}{npgBlue}             % species edge
\colorlet{scRxF}{npgOrange!12}        % reaction fill
\colorlet{scRxE}{npgOrange!88!black}  % reaction edge
\colorlet{scReac}{npgBlue!78!black}   % reactant edge
\colorlet{scProd}{npgGreen!82!black}  % product edge
\colorlet{scHiF}{npgOrange!13}        % highlight fill
\colorlet{scHiE}{npgOrange!82!black}  % highlight edge
% Stage / accent tints.  stDE remains available for the fifth plot series.
\colorlet{stAF}{npgGreen!9}\colorlet{stAE}{npgGreen}
\colorlet{stBF}{npgOrange!11}\colorlet{stBE}{npgOrange}
\colorlet{stCF}{npgBlue!9}\colorlet{stCE}{npgBlue}
\colorlet{stDF}{npgPurple!10}\colorlet{stDE}{npgPurple!85!black}
\colorlet{stEF}{npgVermillion!9}\colorlet{stEE}{npgVermillion}

% Semantic aliases.  Figures should speak in these, not in stAE/stBE/...
\colorlet{scAccent}{stAE}\colorlet{scAccentF}{stAF}   % structure / normal
\colorlet{scMark}{stBE}\colorlet{scMarkF}{stBF}       % highlight / certificate
\colorlet{scRef}{stCE}\colorlet{scRefF}{stCF}         % secondary series
\colorlet{scAlert}{stEE}\colorlet{scAlertF}{stEF}     % failure / violation
\colorlet{scRule}{figHair}                            % card and axis hairlines
\colorlet{scRuleL}{figHair!55}                        % grid, dividers

% --- text styles ----------------------------------------------------------
% Three sizes, fixed in points so they do not drift with the document class.
% Arial is shared by labels, headings, panel badges, and chemistry labels.
\providecommand{\scTitle}{\scArial\fontsize{7}{8.4}\selectfont\bfseries\color{scInk}}
\providecommand{\scBody}{\scArial\fontsize{6.5}{7.8}\selectfont\color{scInkM}}
\providecommand{\scSmall}{\scArial\fontsize{6}{7.2}\selectfont\color{scInkM}}
\providecommand{\scSub}[1]{{\scSmall\itshape #1}}
\providecommand{\scPanel}[1]{{\scTitle(#1)}}
% Numerals that carry a result: ink, not colour, so the eye reads the value.
\providecommand{\scNum}{\scArial\fontsize{10}{11}\selectfont\bfseries\color{scInk}}

% --- shared TikZ styles ---------------------------------------------------
\tikzset{
  % text never hyphenates inside a figure
  every node/.append style={execute at begin node={%
      \hyphenpenalty=10000 \exhyphenpenalty=10000 \hbadness=10000 }},
  syn/.style={font=\scArial,color=scInk,
              execute at begin picture={\MTversion*{synarial}},
              >={Stealth[round,length=2.4mm,width=1.9mm]}},
  % containers ------------------------------------------------------------
  % One card chrome for the whole paper: white, 3 pt radius, one hairline.
  panel/.style={rounded corners=2.5pt,fill=none,draw=scRule,line width=0.45pt},
  card/.style={rounded corners=2.2pt,fill=scPaper,draw=scRule,line width=0.45pt,
               inner sep=5pt,align=left,anchor=north west},
  % a quiet strip for a panel's closing fact
  note/.style={rounded corners=2.5pt,fill=scPaper,draw=scRuleL,line width=0.4pt,
               inner xsep=6pt,inner ysep=3.5pt,align=left,font=\scSmall},
  % small labelled value pill
  chip/.style={rounded corners=2pt,fill=scPaper,draw=scRule,line width=0.4pt,
               inner xsep=4pt,inner ysep=2pt,font=\scSmall},
  % Circular letter badge: the same accessible panel marker used in Style/*.
  hbullet/.style={circle,fill=scInk,text=white,draw=none,
                  font=\scArial\fontsize{6}{6}\selectfont\bfseries,
                  minimum size=4.2mm,inner sep=0pt},
  ctag/.style={font=\scTitle,anchor=north west,text=scInk},
  % CRN primitives --------------------------------------------------------
  species/.style={rounded rectangle,fill=scSpF,draw=scSpE,line width=0.45pt,
                  inner xsep=2.4pt,inner ysep=1.6pt,font=\scSmall,
                  text=scInk,minimum height=3.6mm},
  cofactor/.style={species,fill=scRefF,draw=scRef},
  catalyst/.style={species,fill=scHiF,draw=scHiE},
  rxn/.style={rectangle,fill=scRxF,draw=scRxE,line width=0.45pt,
              inner sep=1.2pt,font=\scSmall,text=scInk,
              minimum width=3.0mm,minimum height=3.0mm},
  place/.style={circle,fill=scSpF,draw=scSpE,line width=0.45pt,
                minimum size=4.0mm,inner sep=0pt},
  trans/.style={rectangle,fill=scRxE,draw=scRxE,minimum width=1.1mm,
                minimum height=3.6mm,inner sep=0pt},
  token/.style={circle,fill=scInk,minimum size=1.1mm,inner sep=0pt},
  % edges -----------------------------------------------------------------
  reac/.style={->,line width=0.5pt,draw=scReac,shorten >=0.8pt,shorten <=0.8pt},
  prod/.style={->,line width=0.5pt,draw=scProd,shorten >=0.8pt,shorten <=0.8pt},
  cat/.style={->,line width=0.5pt,draw=scHiE,shorten >=0.8pt,shorten <=0.8pt},
  flow/.style={-{Stealth[round,length=2.2mm,width=1.8mm]},line width=1.1pt,
               draw=scInkM!45},
  flab/.style={font=\scSmall,text=scInkM,inner sep=2pt},
  stoi/.style={font=\scSmall,text=scInkM,inner sep=0.8pt},
  % verdict marks: shape carries the meaning, colour only reinforces it
  vpass/.style={circle,fill=scAccentF,draw=scAccent,text=scAccent,
                font=\fontsize{6}{6}\selectfont\bfseries,
                minimum size=3.2mm,inner sep=0pt},
  vfail/.style={circle,fill=scAlertF,draw=scAlert,text=scAlert,
                font=\fontsize{6}{6}\selectfont\bfseries,
                minimum size=3.2mm,inner sep=0pt},
  callout/.style={font=\scSmall,text=scInkM,align=left},
}

% \scHead[<badge colour>]{<anchor coord>}{<letter>}{<title>}{<subtitle or empty>}
% One header treatment for every panel in the paper.
\providecommand{\scHead}[5][scInk]{%
  \node[hbullet,fill=#1,anchor=north west] (scbadge#3) at ($#2+(0,-0.015)$)
    {\MakeUppercase{#3}};
  \node[ctag,anchor=west] at ($(scbadge#3.east)+(0.14,0)$) {#4};
  \ifx\relax#5\relax\else
    \node[anchor=north west,font=\scSmall,text=scInkM]
      at ($#2+(0.56,-0.245)$) {#5};
  \fi}

% --- shared pgfplots look (only if the figure loaded pgfplots) ------------
\ifdefined\pgfplotsset
  \pgfplotsset{compat=1.18}
  \pgfplotsset{
    synaxis/.style={
      scale only axis,
      tick align=outside, tick pos=left,
      axis line style={draw=scInkM!55,line width=0.4pt},
      % One grid, on the value axis only.  No minor grid: at these sizes it
      % turns the plot frame into a grey wash.
      ymajorgrids=true, xmajorgrids=false,
      major grid style={draw=scRuleL,line width=0.3pt},
      major tick length=1.6pt,
      label style={font=\scBody,color=scInk},
      tick label style={font=\scSmall,color=scInkM},
      legend style={font=\scSmall,draw=none,fill=none,inner sep=1pt,
                    row sep=0.4pt,legend columns=1},
      legend cell align=left,
      every axis plot/.append style={line width=0.7pt,mark size=1.2pt},
    },
  }
\fi
}
% so that the figure compiles both from Figure/ and from the paper root.
%
% The visual vocabulary follows ../../../../Style/style.tex: the
% colour-blind-safe Okabe--Ito accents, quiet neutral rules and wash, circular
% panel badges, and a compact technical type scale are shared by the companion
% papers.  Semantic aliases below keep the CRN meaning explicit.
%
% DESIGN CONTRACT -----------------------------------------------------------
% 1. Geometry.  Every figure is laid out on a canvas exactly \scW wide (the
%    sn-jnl text width, 372 pt) and is included unscaled, so one point of type
%    is the same physical size in all five figures.  Open each picture with
%    \scCanvas{<height in cm>} to pin the bounding box.
% 2. Colour is semantic, never decorative.  Four roles only:
%       scAccent (green)       structure, the normal/expected path
%       scMark   (orange)      catalyst, certificate, or highlight
%       scAlert  (vermillion)  failure, violation, or uncovered set
%       scRef    (blue)        secondary or reference series
%    A panel gets at most one accent.  Titles are always ink; panel identity
%    is also encoded by a letter inside a circular badge, not by colour alone.
% 3. Type scale.  Three sizes, no others, and nothing below 6 pt:
%       \scTitle (7 pt bold)  panel titles
%       \scBody  (6.5 pt)     labels, values, in-figure prose
%       \scSmall (6 pt)       units, footnotes, dense graph labels
% ---------------------------------------------------------------------------
% Load guard.  The paper inputs every figure source inline, so this file is
% reached once per figure; without the guard the second \newlength below
% would abort the run.
\ifdefined\SynKitStyleLoaded\endinput\fi
\newcommand{\SynKitStyleLoaded}{}

% Load in the document preamble. XeLaTeX/LuaLaTeX and an installed Arial
% family are required; fontspec reports an error instead of substituting.
\RequirePackage[no-math]{fontspec}
\newfontfamily\scArial{Arial}[NFSSFamily=scArial,Ligatures=TeX]
% Keep manuscript mathematics unchanged; select Arial letters and numerals
% only inside figures. TeX supplies specialist mathematical symbols.
\RequirePackage[subdued,italic]{mathastext}
\MTDeclareVersion{synarial}{TU}{scArial}{m}{n}
\SetMathAlphabet{\Mathsf}{synarial}{TU}{scArial}{m}{n}
\SetMathAlphabet{\Mathtt}{synarial}{TU}{scArial}{m}{n}

\usetikzlibrary{positioning,calc,fit,backgrounds,arrows.meta,shapes.geometric,
                shapes.misc,decorations.pathreplacing,patterns,math}

% --- canvas ---------------------------------------------------------------
% 372 pt text width minus the 1 pt standalone border on each side.
\newlength{\scW}\setlength{\scW}{13.0326cm}
\providecommand{\scWnum}{13.0326}
% \scCanvas{<height>} pins the picture to the full text width, y running down
% from 0, so every figure shares one coordinate convention.  It *sets* the
% bounding box rather than contributing to it, so the emitted page is exactly
% the text width and every figure is included at 1:1 -- otherwise a stray
% invisible extent (pgfplots reserves a little around an axis) rescales one
% figure relative to the others and the same \scSmall prints at two sizes.
% Anything drawn outside the canvas is therefore cropped: keep content inside,
% and check with
%     \providecommand{\scCanvas}[1]{\path[draw=red] (0,0) rectangle (\scWnum,-#1);}
\providecommand{\scCanvas}[1]{%
  \path[use as bounding box] (0,0) rectangle (\scWnum,-#1);}

% --- palette --------------------------------------------------------------
% Canonical companion-paper values (../../../../Style/style.tex).
\definecolor{npgOrange}{HTML}{E69F00}
\definecolor{npgSky}{HTML}{56B4E9}
\definecolor{npgGreen}{HTML}{009E73}
\definecolor{npgYellow}{HTML}{F0E442}
\definecolor{npgBlue}{HTML}{0072B2}
\definecolor{npgVermillion}{HTML}{E64B35}
\definecolor{npgPurple}{HTML}{CC79A7}
\definecolor{figInk}{HTML}{1A1A1A}
\definecolor{figMuted}{HTML}{666666}
\definecolor{figHair}{HTML}{D9D9D9}
\definecolor{figWash}{HTML}{F7F7F7}

\colorlet{scInk}{figInk}              % primary text
\colorlet{scInkM}{figMuted}           % secondary text
\colorlet{scPaper}{figWash}           % quiet cards and note strips
\colorlet{scSpF}{npgSky!12}           % species fill
\colorlet{scSpE}{npgBlue}             % species edge
\colorlet{scRxF}{npgOrange!12}        % reaction fill
\colorlet{scRxE}{npgOrange!88!black}  % reaction edge
\colorlet{scReac}{npgBlue!78!black}   % reactant edge
\colorlet{scProd}{npgGreen!82!black}  % product edge
\colorlet{scHiF}{npgOrange!13}        % highlight fill
\colorlet{scHiE}{npgOrange!82!black}  % highlight edge
% Stage / accent tints.  stDE remains available for the fifth plot series.
\colorlet{stAF}{npgGreen!9}\colorlet{stAE}{npgGreen}
\colorlet{stBF}{npgOrange!11}\colorlet{stBE}{npgOrange}
\colorlet{stCF}{npgBlue!9}\colorlet{stCE}{npgBlue}
\colorlet{stDF}{npgPurple!10}\colorlet{stDE}{npgPurple!85!black}
\colorlet{stEF}{npgVermillion!9}\colorlet{stEE}{npgVermillion}

% Semantic aliases.  Figures should speak in these, not in stAE/stBE/...
\colorlet{scAccent}{stAE}\colorlet{scAccentF}{stAF}   % structure / normal
\colorlet{scMark}{stBE}\colorlet{scMarkF}{stBF}       % highlight / certificate
\colorlet{scRef}{stCE}\colorlet{scRefF}{stCF}         % secondary series
\colorlet{scAlert}{stEE}\colorlet{scAlertF}{stEF}     % failure / violation
\colorlet{scRule}{figHair}                            % card and axis hairlines
\colorlet{scRuleL}{figHair!55}                        % grid, dividers

% --- text styles ----------------------------------------------------------
% Three sizes, fixed in points so they do not drift with the document class.
% Arial is shared by labels, headings, panel badges, and chemistry labels.
\providecommand{\scTitle}{\scArial\fontsize{7}{8.4}\selectfont\bfseries\color{scInk}}
\providecommand{\scBody}{\scArial\fontsize{6.5}{7.8}\selectfont\color{scInkM}}
\providecommand{\scSmall}{\scArial\fontsize{6}{7.2}\selectfont\color{scInkM}}
\providecommand{\scSub}[1]{{\scSmall\itshape #1}}
\providecommand{\scPanel}[1]{{\scTitle(#1)}}
% Numerals that carry a result: ink, not colour, so the eye reads the value.
\providecommand{\scNum}{\scArial\fontsize{10}{11}\selectfont\bfseries\color{scInk}}

% --- shared TikZ styles ---------------------------------------------------
\tikzset{
  % text never hyphenates inside a figure
  every node/.append style={execute at begin node={%
      \hyphenpenalty=10000 \exhyphenpenalty=10000 \hbadness=10000 }},
  syn/.style={font=\scArial,color=scInk,
              execute at begin picture={\MTversion*{synarial}},
              >={Stealth[round,length=2.4mm,width=1.9mm]}},
  % containers ------------------------------------------------------------
  % One card chrome for the whole paper: white, 3 pt radius, one hairline.
  panel/.style={rounded corners=2.5pt,fill=none,draw=scRule,line width=0.45pt},
  card/.style={rounded corners=2.2pt,fill=scPaper,draw=scRule,line width=0.45pt,
               inner sep=5pt,align=left,anchor=north west},
  % a quiet strip for a panel's closing fact
  note/.style={rounded corners=2.5pt,fill=scPaper,draw=scRuleL,line width=0.4pt,
               inner xsep=6pt,inner ysep=3.5pt,align=left,font=\scSmall},
  % small labelled value pill
  chip/.style={rounded corners=2pt,fill=scPaper,draw=scRule,line width=0.4pt,
               inner xsep=4pt,inner ysep=2pt,font=\scSmall},
  % Circular letter badge: the same accessible panel marker used in Style/*.
  hbullet/.style={circle,fill=scInk,text=white,draw=none,
                  font=\scArial\fontsize{6}{6}\selectfont\bfseries,
                  minimum size=4.2mm,inner sep=0pt},
  ctag/.style={font=\scTitle,anchor=north west,text=scInk},
  % CRN primitives --------------------------------------------------------
  species/.style={rounded rectangle,fill=scSpF,draw=scSpE,line width=0.45pt,
                  inner xsep=2.4pt,inner ysep=1.6pt,font=\scSmall,
                  text=scInk,minimum height=3.6mm},
  cofactor/.style={species,fill=scRefF,draw=scRef},
  catalyst/.style={species,fill=scHiF,draw=scHiE},
  rxn/.style={rectangle,fill=scRxF,draw=scRxE,line width=0.45pt,
              inner sep=1.2pt,font=\scSmall,text=scInk,
              minimum width=3.0mm,minimum height=3.0mm},
  place/.style={circle,fill=scSpF,draw=scSpE,line width=0.45pt,
                minimum size=4.0mm,inner sep=0pt},
  trans/.style={rectangle,fill=scRxE,draw=scRxE,minimum width=1.1mm,
                minimum height=3.6mm,inner sep=0pt},
  token/.style={circle,fill=scInk,minimum size=1.1mm,inner sep=0pt},
  % edges -----------------------------------------------------------------
  reac/.style={->,line width=0.5pt,draw=scReac,shorten >=0.8pt,shorten <=0.8pt},
  prod/.style={->,line width=0.5pt,draw=scProd,shorten >=0.8pt,shorten <=0.8pt},
  cat/.style={->,line width=0.5pt,draw=scHiE,shorten >=0.8pt,shorten <=0.8pt},
  flow/.style={-{Stealth[round,length=2.2mm,width=1.8mm]},line width=1.1pt,
               draw=scInkM!45},
  flab/.style={font=\scSmall,text=scInkM,inner sep=2pt},
  stoi/.style={font=\scSmall,text=scInkM,inner sep=0.8pt},
  % verdict marks: shape carries the meaning, colour only reinforces it
  vpass/.style={circle,fill=scAccentF,draw=scAccent,text=scAccent,
                font=\fontsize{6}{6}\selectfont\bfseries,
                minimum size=3.2mm,inner sep=0pt},
  vfail/.style={circle,fill=scAlertF,draw=scAlert,text=scAlert,
                font=\fontsize{6}{6}\selectfont\bfseries,
                minimum size=3.2mm,inner sep=0pt},
  callout/.style={font=\scSmall,text=scInkM,align=left},
}

% \scHead[<badge colour>]{<anchor coord>}{<letter>}{<title>}{<subtitle or empty>}
% One header treatment for every panel in the paper.
\providecommand{\scHead}[5][scInk]{%
  \node[hbullet,fill=#1,anchor=north west] (scbadge#3) at ($#2+(0,-0.015)$)
    {\MakeUppercase{#3}};
  \node[ctag,anchor=west] at ($(scbadge#3.east)+(0.14,0)$) {#4};
  \ifx\relax#5\relax\else
    \node[anchor=north west,font=\scSmall,text=scInkM]
      at ($#2+(0.56,-0.245)$) {#5};
  \fi}

% --- shared pgfplots look (only if the figure loaded pgfplots) ------------
\ifdefined\pgfplotsset
  \pgfplotsset{compat=1.18}
  \pgfplotsset{
    synaxis/.style={
      scale only axis,
      tick align=outside, tick pos=left,
      axis line style={draw=scInkM!55,line width=0.4pt},
      % One grid, on the value axis only.  No minor grid: at these sizes it
      % turns the plot frame into a grey wash.
      ymajorgrids=true, xmajorgrids=false,
      major grid style={draw=scRuleL,line width=0.3pt},
      major tick length=1.6pt,
      label style={font=\scBody,color=scInk},
      tick label style={font=\scSmall,color=scInkM},
      legend style={font=\scSmall,draw=none,fill=none,inner sep=1pt,
                    row sep=0.4pt,legend columns=1},
      legend cell align=left,
      every axis plot/.append style={line width=0.7pt,mark size=1.2pt},
    },
  }
\fi
}}

\begingroup
% --- molecule drawing conventions -----------------------------------------
\setchemfig{atom sep=17pt,bond style={line width=0.5pt},bond offset=1.1pt,
            cram width=1.7pt}
\renewcommand*\printatom[1]{\ensuremath{\mathsf{#1}}}
% phosphoryl atoms carry the highlight colour; \pk keeps that in one place
\newcommand{\pk}[1]{\textcolor{scMark}{#1}}
\newcommand{\mfont}{\fontsize{6.5}{7.5}\selectfont\color{scInk}}
% phosphoester links, running right (R) or left (L) from the ester oxygen
\definesubmol{pbR}{-[:-30,0.85,,,draw=scMark]\pk{P}(=[:60,0.8,,,draw=scMark]\pk{O})%
  (-[:-120,0.8,,,draw=scMark]\pk{OH})-[:30,0.85,,,draw=scMark]\pk{O}}
\definesubmol{ptR}{-[:-30,0.85,,,draw=scMark]\pk{P}(=[:60,0.85,,,draw=scMark]\pk{O})%
  (-[:-120,0.8,,,draw=scMark]\pk{OH})-[:-15,0.85,,,draw=scMark]\pk{OH}}
\definesubmol{ptL}{-[:210,0.85,,,draw=scMark]\pk{P}(=[:120,0.85,,,draw=scMark]\pk{O})%
  (-[:-60,0.8,,,draw=scMark]\pk{OH})-[:195,0.85,,,draw=scMark]\pk{OH}}

\begin{tikzpicture}[syn]
% The sn-jnl text block of this manuscript is 466.63 pt = 16.40 cm wide (see
% Figure/arch.tex), so the canvas is that width and the figure is included
% unscaled: one point of type is the same size as in every other figure.
\def\W{16.40}
\path[use as bounding box] (0,0) rectangle (\W,-9.05);

\def\pw{9.20}          % width of the left (wide) panels
\def\qw{7.00}          % width of the right (narrow) panels
\def\gx{9.40}          % left edge of the right-hand column
\def\ra{-5.15}         % top edge of the second row

% =========================================================================
% A  ATP
% =========================================================================
\draw[panel] (0,-0.02) rectangle (\pw,-4.95);
\scHead{(0.22,-0.06)}{a}{ATP {\mdseries (C)}}{$\mathrm{C_{10}H_{16}N_5O_{13}P_3}$ --- phosphoryl donor of $r10$ and $r3$}
\node[anchor=center] at (4.60,-2.75) {\mfont\chemfig{
*6(-*5(-N=-N(-[:-70,1.3]
   ?[r]
   -[:-59.7,0.834,,,line width=0.95pt](-[:-120,0.7]HO)
   -[:0,0.66,,,line width=0.95pt](-[:-60,0.7]OH)
   -[:59.7,0.834,,,line width=0.95pt](-[:80,0.9]-[:10,0.9]\pk{O}!{pbR}!{pbR}!{ptR})
   -[:140.4,0.973]O?[r]
  )-)=(-[:52,0.95]NH_2)-N=-N=)
}};
\node[note,anchor=south west] at (0.22,-4.83) {two phosphoanhydride bonds; $r10$ leaves ADP};

% =========================================================================
% B  alpha-D-glucose
% =========================================================================
\draw[panel] (\gx,-0.02) rectangle (\gx+\qw,-4.95);
\scHead{(\gx+0.22,-0.06)}{b}{$\alpha$-D-Glucose {\mdseries (S)}}{$\mathrm{C_6H_{12}O_6}$ --- phosphorylated by $r10$}
\node[anchor=center] at (\gx+3.50,-2.75) {\mfont\chemfig{
?[a](-[:90,0.9]-[:30,0.9]OH)
-[:222,0.85](-[:-90,0.72]OH)
-[:-40,0.8,,,line width=0.95pt](-[:58,0.68]OH)
-[:0,1.9,,,line width=0.95pt](-[:-90,0.72]OH)
-[:40,0.8,,,line width=0.95pt](-[:-90,0.72]OH)
-[:138,0.85]O?[a]
}};
\node[note,anchor=south west] at (\gx+0.22,-4.83) {anomeric \ce{OH} down; no phosphate before $r10$};

% =========================================================================
% C  fructose 1,6-bisphosphate
% =========================================================================
\draw[panel] (0,\ra) rectangle (\pw,\ra-3.85);
\scHead{(0.22,\ra+0.02)}{c}{Fructose 1,6-bisphosphate {\mdseries (H)}}{$\mathrm{C_6H_{14}O_{12}P_2}$ --- product of $r3$ and $r12$}
\node[anchor=center] at (4.60,\ra-2.05) {\mfont\chemfig{
?[f](-[:100,0.9]-[:170,0.9]\pk{O}!{ptL})
-[:-59.7,0.834,,,line width=0.95pt](-[:-120,0.7]HO)
-[:0,0.66,,,line width=0.95pt](-[:108,0.62]OH)
-[:59.7,0.834,,,line width=0.95pt](-[:-25,0.78]OH)(-[:78,0.9]-[:12,0.9]\pk{O}!{ptR})
-[:140.4,0.973]O?[f]
}};
\node[note,anchor=south west] at (0.22,\ra-3.73) {cleaved by aldolase $r7$; C2 configuration not specified};

% =========================================================================
% D  D-glyceraldehyde 3-phosphate
% =========================================================================
\draw[panel] (\gx,\ra) rectangle (\gx+\qw,\ra-3.85);
\scHead{(\gx+0.22,\ra+0.02)}{d}{D-Glyceraldehyde 3-phosphate {\mdseries (J)}}{$\mathrm{C_3H_7O_6P}$ --- target of both routes}
\node[anchor=center] at (\gx+3.50,\ra-2.05) {\mfont\chemfig{
\pk{HO}-[:30,0.85,,,draw=scMark]\pk{P}(=[:90,0.85,,,draw=scMark]\pk{O})%
(-[:-90,0.8,,,draw=scMark]\pk{OH})-[:-30,0.85,,,draw=scMark]\pk{O}%
-[:-30,0.85]-[:30,0.85](<[:90,0.75]OH)-[:-30,0.85](=[:-60,0.8]O)-[:30,0.7]H
}};
\node[note,anchor=south west] at (\gx+0.22,\ra-3.73) {(\textit{R}) centre; released with DHAP by $r7$};

\end{tikzpicture}
\endgroup
